% Options for packages loaded elsewhere
\PassOptionsToPackage{unicode}{hyperref}
\PassOptionsToPackage{hyphens}{url}
\PassOptionsToPackage{dvipsnames,svgnames,x11names}{xcolor}
\documentclass[
  italian,
  11pt,
]{article}
\usepackage{xcolor}
\usepackage[top=3cm, bottom=3cm, left=2.5cm, right=2.5cm]{geometry}
\usepackage{amsmath,amssymb}
\setcounter{secnumdepth}{5}
\usepackage{iftex}
\ifPDFTeX
  \usepackage[T1]{fontenc}
  \usepackage[utf8]{inputenc}
  \usepackage{textcomp} % provide euro and other symbols
\else % if luatex or xetex
  \usepackage{unicode-math} % this also loads fontspec
  \defaultfontfeatures{Scale=MatchLowercase}
  \defaultfontfeatures[\rmfamily]{Ligatures=TeX,Scale=1}
\fi
\usepackage{lmodern}
\ifPDFTeX\else
  % xetex/luatex font selection
\fi
% Use upquote if available, for straight quotes in verbatim environments
\IfFileExists{upquote.sty}{\usepackage{upquote}}{}
\IfFileExists{microtype.sty}{% use microtype if available
  \usepackage[]{microtype}
  \UseMicrotypeSet[protrusion]{basicmath} % disable protrusion for tt fonts
}{}
\usepackage{setspace}
\makeatletter
\@ifundefined{KOMAClassName}{% if non-KOMA class
  \IfFileExists{parskip.sty}{%
    \usepackage{parskip}
  }{% else
    \setlength{\parindent}{0pt}
    \setlength{\parskip}{6pt plus 2pt minus 1pt}}
}{% if KOMA class
  \KOMAoptions{parskip=half}}
\makeatother
\usepackage{color}
\usepackage{fancyvrb}
\newcommand{\VerbBar}{|}
\newcommand{\VERB}{\Verb[commandchars=\\\{\}]}
\DefineVerbatimEnvironment{Highlighting}{Verbatim}{commandchars=\\\{\}}
% Add ',fontsize=\small' for more characters per line
\newenvironment{Shaded}{}{}
\newcommand{\AlertTok}[1]{\textcolor[rgb]{1.00,0.00,0.00}{\textbf{#1}}}
\newcommand{\AnnotationTok}[1]{\textcolor[rgb]{0.38,0.63,0.69}{\textbf{\textit{#1}}}}
\newcommand{\AttributeTok}[1]{\textcolor[rgb]{0.49,0.56,0.16}{#1}}
\newcommand{\BaseNTok}[1]{\textcolor[rgb]{0.25,0.63,0.44}{#1}}
\newcommand{\BuiltInTok}[1]{\textcolor[rgb]{0.00,0.50,0.00}{#1}}
\newcommand{\CharTok}[1]{\textcolor[rgb]{0.25,0.44,0.63}{#1}}
\newcommand{\CommentTok}[1]{\textcolor[rgb]{0.38,0.63,0.69}{\textit{#1}}}
\newcommand{\CommentVarTok}[1]{\textcolor[rgb]{0.38,0.63,0.69}{\textbf{\textit{#1}}}}
\newcommand{\ConstantTok}[1]{\textcolor[rgb]{0.53,0.00,0.00}{#1}}
\newcommand{\ControlFlowTok}[1]{\textcolor[rgb]{0.00,0.44,0.13}{\textbf{#1}}}
\newcommand{\DataTypeTok}[1]{\textcolor[rgb]{0.56,0.13,0.00}{#1}}
\newcommand{\DecValTok}[1]{\textcolor[rgb]{0.25,0.63,0.44}{#1}}
\newcommand{\DocumentationTok}[1]{\textcolor[rgb]{0.73,0.13,0.13}{\textit{#1}}}
\newcommand{\ErrorTok}[1]{\textcolor[rgb]{1.00,0.00,0.00}{\textbf{#1}}}
\newcommand{\ExtensionTok}[1]{#1}
\newcommand{\FloatTok}[1]{\textcolor[rgb]{0.25,0.63,0.44}{#1}}
\newcommand{\FunctionTok}[1]{\textcolor[rgb]{0.02,0.16,0.49}{#1}}
\newcommand{\ImportTok}[1]{\textcolor[rgb]{0.00,0.50,0.00}{\textbf{#1}}}
\newcommand{\InformationTok}[1]{\textcolor[rgb]{0.38,0.63,0.69}{\textbf{\textit{#1}}}}
\newcommand{\KeywordTok}[1]{\textcolor[rgb]{0.00,0.44,0.13}{\textbf{#1}}}
\newcommand{\NormalTok}[1]{#1}
\newcommand{\OperatorTok}[1]{\textcolor[rgb]{0.40,0.40,0.40}{#1}}
\newcommand{\OtherTok}[1]{\textcolor[rgb]{0.00,0.44,0.13}{#1}}
\newcommand{\PreprocessorTok}[1]{\textcolor[rgb]{0.74,0.48,0.00}{#1}}
\newcommand{\RegionMarkerTok}[1]{#1}
\newcommand{\SpecialCharTok}[1]{\textcolor[rgb]{0.25,0.44,0.63}{#1}}
\newcommand{\SpecialStringTok}[1]{\textcolor[rgb]{0.73,0.40,0.53}{#1}}
\newcommand{\StringTok}[1]{\textcolor[rgb]{0.25,0.44,0.63}{#1}}
\newcommand{\VariableTok}[1]{\textcolor[rgb]{0.10,0.09,0.49}{#1}}
\newcommand{\VerbatimStringTok}[1]{\textcolor[rgb]{0.25,0.44,0.63}{#1}}
\newcommand{\WarningTok}[1]{\textcolor[rgb]{0.38,0.63,0.69}{\textbf{\textit{#1}}}}
\usepackage{longtable,booktabs,array}
\newcounter{none} % for unnumbered tables
\usepackage{calc} % for calculating minipage widths
% Correct order of tables after \paragraph or \subparagraph
\usepackage{etoolbox}
\makeatletter
\patchcmd\longtable{\par}{\if@noskipsec\mbox{}\fi\par}{}{}
\makeatother
% Allow footnotes in longtable head/foot
\IfFileExists{footnotehyper.sty}{\usepackage{footnotehyper}}{\usepackage{footnote}}
\makesavenoteenv{longtable}
\usepackage{graphicx}
\makeatletter
\newsavebox\pandoc@box
\newcommand*\pandocbounded[1]{% scales image to fit in text height/width
  \sbox\pandoc@box{#1}%
  \Gscale@div\@tempa{\textheight}{\dimexpr\ht\pandoc@box+\dp\pandoc@box\relax}%
  \Gscale@div\@tempb{\linewidth}{\wd\pandoc@box}%
  \ifdim\@tempb\p@<\@tempa\p@\let\@tempa\@tempb\fi% select the smaller of both
  \ifdim\@tempa\p@<\p@\scalebox{\@tempa}{\usebox\pandoc@box}%
  \else\usebox{\pandoc@box}%
  \fi%
}
% Set default figure placement to htbp
\def\fps@figure{htbp}
\makeatother
\ifLuaTeX
\usepackage[bidi=basic,shorthands=off]{babel}
\else
\usepackage[bidi=default,shorthands=off]{babel}
\fi
\ifLuaTeX
  \usepackage{selnolig} % disable illegal ligatures
\fi
\setlength{\emergencystretch}{3em} % prevent overfull lines
\providecommand{\tightlist}{%
  \setlength{\itemsep}{0pt}\setlength{\parskip}{0pt}}
\usepackage{mathpazo}
\usepackage[T1]{fontenc}
\usepackage{mathtools}
\usepackage{bm}
\usepackage{booktabs}
\usepackage{longtable}
\usepackage{array}
\usepackage{multirow}
\usepackage{fancyhdr}
\pagestyle{fancy}
\fancyhead[L]{}
\fancyhead[R]{\leftmark}
\fancyfoot[C]{\thepage}
\usepackage[font=small,labelfont=bf]{caption}
\usepackage{framed}
\definecolor{shadecolor}{RGB}{248,248,248}
\usepackage{float}
\let\origfigure\figure
\let\endorigfigure\endfigure
\renewenvironment{figure}[1][2] {\expandafter\origfigure\expandafter[H]} {\endorigfigure}
\usepackage{etoolbox}
\AtBeginEnvironment{longtable}{\small}
\usepackage{bookmark}
\IfFileExists{xurl.sty}{\usepackage{xurl}}{} % add URL line breaks if available
\urlstyle{same}
\hypersetup{
  pdftitle={Machine Learning Notes},
  pdfauthor={Nicola Moscufo},
  pdflang={it},
  colorlinks=true,
  linkcolor={NavyBlue},
  filecolor={Maroon},
  citecolor={Blue},
  urlcolor={NavyBlue},
  pdfcreator={LaTeX via pandoc}}

\title{Machine Learning Notes}
\author{Nicola Moscufo}
\date{4 Marzo 2026}

\begin{document}
\maketitle

{
\hypersetup{linkcolor=Black}
\setcounter{tocdepth}{3}
\tableofcontents
}
\setstretch{1.25}
\newpage

\section{The Statistical Learning
Framework}\label{the-statistical-learning-framework}

The Statistical Learning Framework provides the formal mathematical
foundation for machine learning. It defines the components of the
learning problem, the assumptions made about data generation, and the
criteria used to evaluate the success of a learning algorithm.

\subsection{1. The Learner's Input}\label{the-learners-input}

In the basic statistical learning setting (specifically Supervised
Learning), the learner has access to the following components:

\subsubsection{\texorpdfstring{A. The Domain Set
(\(\mathcal{X}\))}{A. The Domain Set (\textbackslash mathcal\{X\})}}\label{a.-the-domain-set-mathcalx}

The set of all possible objects that we wish to classify or label. This
is often represented as a vector space of features.

\begin{itemize}
\tightlist
\item
  \textbf{Notation:} \(\mathcal{X}\)
\item
  \textbf{Example:} In an image classification task, \(\mathcal{X}\) is
  the set of all possible images (e.g., \(\mathbb{R}^{d}\) where \(d\)
  is the number of pixels).
\end{itemize}

\subsubsection{\texorpdfstring{B. The Label Set
(\(\mathcal{Y}\))}{B. The Label Set (\textbackslash mathcal\{Y\})}}\label{b.-the-label-set-mathcaly}

The set of all possible outcomes or labels for the objects in the
domain.

\begin{itemize}
\tightlist
\item
  \textbf{Notation:} \(\mathcal{Y}\)
\item
  \textbf{Example:}

  \begin{itemize}
  \tightlist
  \item
    Binary Classification: \(\mathcal{Y} = \{0, 1\}\) or \(\{-1, +1\}\).
  \item
    Regression: \(\mathcal{Y} = \mathbb{R}\).
  \end{itemize}
\end{itemize}

\subsubsection{\texorpdfstring{C. Training Data
(\(S\))}{C. Training Data (S)}}\label{c.-training-data-s}

A finite sequence of labeled examples given to the learner.

\begin{itemize}
\tightlist
\item
  \textbf{Notation:}
  \(S = \{(x_1, y_1), (x_2, y_2), \dots, (x_m, y_m)\}\)
\item
  \textbf{Structure:} Each pair consists of an input
  \(x_i \in \mathcal{X}\) and a label \(y_i \in \mathcal{Y}\).
\item
  \textbf{Sample Size:} The number of examples \(m\) is called the
  sample size.
\end{itemize}

\subsection{2. The Data Generation Model (Hidden
Components)}\label{the-data-generation-model-hidden-components}

A crucial aspect of the framework is that the learner does not see the
full picture. The data is generated by processes unknown to the learner.

\subsubsection{\texorpdfstring{A. The Underlying Distribution
(\(\mathcal{D}\))}{A. The Underlying Distribution (\textbackslash mathcal\{D\})}}\label{a.-the-underlying-distribution-mathcald}

There exists a probability distribution \(\mathcal{D}\) over the domain
set \(\mathcal{X}\).

\begin{itemize}
\tightlist
\item
  \textbf{Role:} It determines how likely we are to encounter different
  instances in the real world.
\item
  \textbf{Assumption:} The marginal distribution is unknown to the
  learner.
\end{itemize}

\subsubsection{\texorpdfstring{B. The Labeling Function
(\(f\))}{B. The Labeling Function (f)}}\label{b.-the-labeling-function-f}

There is a ``true'' target function that assigns the correct label to
every point in the domain.

\begin{itemize}
\tightlist
\item
  \textbf{Notation:} \(f: \mathcal{X} \to \mathcal{Y}\)
\item
  \textbf{Assumption:} For any \(x_i\) in the training data,
  \(y_i = f(x_i)\).
\item
  \textbf{Blindness:} The learner is blind to both \(\mathcal{D}\) and
  \(f\). The learner's goal is to discover or approximate \(f\).
\end{itemize}

\subsubsection{C. The i.i.d. Assumption}\label{c.-the-i.i.d.-assumption}

We assume that the training data \(S\) is generated by sampling points
\(x\) independently and identically distributed (i.i.d.) according to
\(\mathcal{D}\), and then labeling them using \(f\).
\[ S \sim \mathcal{D}^m \quad \text{and} \quad y_i = f(x_i) \]

\subsection{3. The Learner's Output}\label{the-learners-output}

The learner processes the training data \(S\) to produce a model.

\subsubsection{\texorpdfstring{The Prediction Rule
(\(h\))}{The Prediction Rule (h)}}\label{the-prediction-rule-h}

The output is a function \(h: \mathcal{X} \to \mathcal{Y}\).

\begin{itemize}
\tightlist
\item
  \textbf{Terminology:} This function is variously called a
  \textbf{predictor}, a \textbf{hypothesis}, or a \textbf{classifier}.
\item
  \textbf{Purpose:} The predictor is used to predict the label of new
  domain points \(x\) that were not present in the training set.
\item
  \textbf{Hypothesis Class (\(\mathcal{H}\)):} Often, the learner
  selects \(h\) from a predefined set of functions \(\mathcal{H}\)
  (e.g., the set of all linear classifiers).
\end{itemize}

\subsection{4. Measure of Success}\label{measure-of-success}

To evaluate a classifier, we quantify how well it performs on the entire
domain \(\mathcal{X}\), weighted by the probability of seeing each
instance.

\subsubsection{The Generalization Error (True
Risk)}\label{the-generalization-error-true-risk}

We define the \textbf{error} (or risk) of a classifier \(h\) as the
probability that it predicts the incorrect label for a random data point
generated by the underlying distribution \(\mathcal{D}\).

\textbf{Formal Definition (\(L_{\mathcal{D}}(h)\)):}
\[ L_{\mathcal{D}, f}(h) \overset{\text{def}}{=} \mathbb{P}_{x \sim \mathcal{D}} [h(x) \neq f(x)] \]

Mathematically, using the indicator function \(\mathbb{1}\):
\[ L_{\mathcal{D}}(h) = \mathbb{E}_{x \sim \mathcal{D}} [\mathbb{1}_{\{h(x) \neq f(x)\}}] \]

\begin{itemize}
\tightlist
\item
  \textbf{Intuition:} This represents the expected failure rate on
  \emph{unseen} data from the real world.
\item
  \textbf{Goal:} The learner aims to find \(h\) such that
  \(L_{\mathcal{D}}(h)\) is minimized.
\end{itemize}

\subsubsection{Empirical Error (Training
Error)}\label{empirical-error-training-error}

Since \(\mathcal{D}\) is unknown, the learner cannot calculate the True
Risk directly. Instead, it can calculate the error on the training data
\(S\).

\textbf{Formal Definition (\(L_S(h)\)):}
\[ L_S(h) \overset{\text{def}}{=} \frac{1}{m} \sum_{i=1}^{m} \mathbb{1}_{\{h(x_i) \neq y_i\}} \]

\begin{itemize}
\tightlist
\item
  \textbf{Relationship:} In Empirical Risk Minimization (ERM), we hope
  that minimizing \(L_S(h)\) (which we can see) leads to a low
  \(L_{\mathcal{D}}(h)\) (which we cannot see).
\end{itemize}

\subsection{Key Takeaways}\label{key-takeaways}

\begin{enumerate}
\def\labelenumi{\arabic{enumi}.}
\tightlist
\item
  \textbf{Inputs:} Domain \(\mathcal{X}\), Label set \(\mathcal{Y}\),
  Training set \(S\).
\item
  \textbf{Hidden World:} Probability distribution \(\mathcal{D}\) and
  Labeling function \(f\) are fixed but unknown.
\item
  \textbf{Output:} A hypothesis \(h\) that approximates \(f\).
\item
  \textbf{Goal:} Minimize the \textbf{Generalization Error} (True Risk),
  which is the probability of error on new data sampled from
  \(\mathcal{D}\).
\item
  \textbf{Challenge:} The learner must generalize from a finite sample
  \(S\) to the entire domain \(\mathcal{X}\) without knowing
  \(\mathcal{D}\).
\end{enumerate}

\newpage

\section{Empirical Risk Minimization
(ERM)}\label{empirical-risk-minimization-erm}

The ultimate goal of a learning algorithm is to find a hypothesis \(h\)
that minimizes the \textbf{True Risk} (Generalization Error) with
respect to the unknown distribution \(\mathcal{D}\) and labeling
function \(f\):
\[ L_{\mathcal{D},f}(h) = \mathbb{P}_{x \sim \mathcal{D}}[h(x) \neq f(x)] \]

Since the learner does not know \(\mathcal{D}\) or \(f\), the true error
is \textbf{not directly available}. The learner cannot minimize what it
cannot calculate. However, a useful notion that \emph{can} be calculated
is the \textbf{Training Error} (or Empirical Error)---the error the
classifier incurs over the available training sample \(S\).

\subsubsection{\texorpdfstring{Definition: Empirical Risk
(\(L_S(h)\))}{Definition: Empirical Risk (L\_S(h))}}\label{definition-empirical-risk-l_sh}

Given a training set \(S = \{(x_1, y_1), \dots, (x_m, y_m)\}\), the
empirical risk (or error) of a hypothesis \(h\) is the average error
over the sample:

\[ L_S(h) \overset{\text{def}}{=} \frac{1}{m} \sum_{i=1}^{m} \mathbb{1}_{\{h(x_i) \neq y_i\}} \]
Where \(\mathbb{1}_{\{\cdot\}}\) is the indicator function (1 if the
condition is true, 0 otherwise).

\subsubsection{The ERM Rule}\label{the-erm-rule}

The Empirical Risk Minimization (ERM) rule is a learning strategy that
instructs the learner to select a hypothesis \(h_S\) that minimizes the
empirical error on the training set \(S\).

\[ h_S \in \underset{h}{\text{argmin}} \ L_S(h) \]

\textbf{Intuition:} Since the training data is a ``window'' into the
true distribution \(\mathcal{D}\), we hope that a predictor performing
well on \(S\) will also perform well on \(\mathcal{D}\).

\subsection{The Danger of Overfitting}\label{the-danger-of-overfitting}

While ERM is intuitive, blindly minimizing \(L_S(h)\) can lead to
failure. This phenomenon is called \textbf{Overfitting}.

\textbf{Definition:} Overfitting occurs when the ERM rule outputs a
hypothesis \(h\) such that:

\begin{enumerate}
\def\labelenumi{\arabic{enumi}.}
\tightlist
\item
  The empirical error \(L_S(h)\) is very small (or zero).
\item
  The true error \(L_{\mathcal{D}}(h)\) is large.
\end{enumerate}

\textbf{Why does this happen?} Without constraints, a learner might
``memorize'' the training data rather than learning the underlying
pattern. Consider a predictor that simply stores the training set:
\[ h(x) = \begin{cases} y_i & \text{if } x = x_i \in S \\ 0 & \text{otherwise} \end{cases} \]
This predictor has \(L_S(h) = 0\) (perfect training accuracy), but on
any new data point not in \(S\), it outputs a default label, likely
resulting in high generalization error.

\subsection{ERM with Inductive Bias}\label{erm-with-inductive-bias}

To prevent overfitting, we must restrict the learner from choosing
\emph{any} possible function. We require the learner to choose a
predictor from a specific, predefined set.

\subsubsection{\texorpdfstring{Hypothesis Class
(\(\mathcal{H}\))}{Hypothesis Class (\textbackslash mathcal\{H\})}}\label{hypothesis-class-mathcalh}

The learner chooses in advance a set of \textbf{candidate predictors},
denoted by \(\mathcal{H}\). This is the ``search space'' of the
algorithm.

\begin{itemize}
\tightlist
\item
  Examples of \(\mathcal{H}\): The set of all linear classifiers, the
  set of decision trees of depth 3, etc.
\end{itemize}

\subsubsection{Constrained ERM Formula}\label{constrained-erm-formula}

The ERM learner restricts the search to \(\mathcal{H}\):

\[ ERM_{\mathcal{H}}(S) \in \underset{h \in \mathcal{H}}{\text{argmin}} \ L_S(h) \]

\subsubsection{Inductive Bias}\label{inductive-bias}

The restriction to a specific class \(\mathcal{H}\) is called an
\textbf{Inductive Bias}.

\begin{itemize}
\tightlist
\item
  \textbf{Bias:} We are ``biased'' toward believing the truth lies
  within \(\mathcal{H}\).
\item
  \textbf{Trade-off:}

  \begin{itemize}
  \tightlist
  \item
    A \textbf{stronger} bias (smaller \(\mathcal{H}\)) \textbf{reduces
    the risk of overfitting} but \emph{increases the risk that the true
    target function} \(f\) is not in \(\mathcal{H}\) (underfitting).
  \item
    A \textbf{weaker} bias (larger \(\mathcal{H}\)) increases the
    capacity to approximate complex functions but increases the risk of
    overfitting.
  \end{itemize}
\end{itemize}

\subsection{Finite Hypothesis Class}\label{finite-hypothesis-class}

The simplest type of inductive bias is assuming that the hypothesis
class \(\mathcal{H}\) has a finite number of elements.

\[ |\mathcal{H}| < \infty \] \#\#\# The Realizability Assumption To
analyze performance in this setting, we often make the
\textbf{Realizability Assumption}: We assume that the true labeling
function \(f\) actually exists within our hypothesis class (or
effectively, that there is some \(h^* \in \mathcal{H}\) that makes no
errors).
\[ \exists h^* \in \mathcal{H} \text{ such that } L_{\mathcal{D}}(h^*) = 0 \]
Note that this implies \(L_S(h^*) = 0\) as well.

\subsubsection{Why Finite Classes Help}\label{why-finite-classes-help}

\textbf{If \(\mathcal{H}\) is finite, the learner cannot ``invent'' a
complex function to memorize noise}. * If \(|\mathcal{H}|\) is small and
the sample size \(m\) is large enough, it becomes statistically unlikely
that a ``bad'' hypothesis (one with high true error) will accidentally
classify all training examples perfectly. * Therefore, if we find an
\(h \in \mathcal{H}\) with \(L_S(h) = 0\), it is likely that
\(L_{\mathcal{D}}(h)\) is also low.

\subsubsection{Theoretical Guarantee
(Preview)}\label{theoretical-guarantee-preview}

For a finite hypothesis class \(\mathcal{H}\), the sample size \(m\)
required to guarantee a generalization error of at most \(\epsilon\)
(with probability \(1-\delta\)) grows logarithmically with the size of
the class: \[ m \propto \log(|\mathcal{H}|) \]

\subsection{Key Takeaways}\label{key-takeaways-1}

\begin{enumerate}
\def\labelenumi{\arabic{enumi}.}
\tightlist
\item
  \textbf{ERM Principle:} We minimize Training Error (\(L_S\)) because
  True Error (\(L_{\mathcal{D}}\)) is unknown.
\item
  \textbf{Overfitting:} Perfect performance on training data does not
  guarantee good performance on unseen data.
\item
  \textbf{Inductive Bias:} We must restrict the search space to a
  specific Hypothesis Class \(\mathcal{H}\) to enable generalization.
\item
  \textbf{Finite \(\mathcal{H}\):} Restricting \(\mathcal{H}\) to a
  finite set is a basic form of inductive bias that makes learning
  theoretically feasible under the Realizability Assumption.
\end{enumerate}

\newpage

\section{The No-Free Lunch Theorem}\label{the-no-free-lunch-theorem}

The \textbf{No-Free Lunch (NFL) Theorem} is a fundamental impossibility
result in statistical learning theory. It asserts that no learning
algorithm can be universally good. It provides the formal justification
for why \textbf{inductive bias} is necessary for learning.

\subsection{Formal Statement}\label{formal-statement}

Let \(\mathcal{A}\) be any learning algorithm for the task of binary
classification with respect to the 0-1 loss over a domain
\(\mathcal{X}\). Let the training set size be \(m < |\mathcal{X}|/2\).

\textbf{Theorem:} For every learning algorithm \(\mathcal{A}\), there
exists a distribution \(\mathcal{D}\) over
\(\mathcal{X} \times \{0,1\}\) and a target labeling function \(f\) such
that:

\begin{enumerate}
\def\labelenumi{\arabic{enumi}.}
\tightlist
\item
  \textbf{High Error:} The generalization error of the hypothesis
  \(h_S\) returned by the algorithm is high (specifically, at least
  1/4).
  \[ \mathbb{E}_{S \sim \mathcal{D}^m} [L_{\mathcal{D}}(h_S)] \geq \frac{1}{4} \]
\item
  \textbf{Learnability:} There exists \emph{another} algorithm that
  could have learned this specific task with zero error.
\end{enumerate}

\subsection{Intuition: The Limits of
Data}\label{intuition-the-limits-of-data}

The theorem addresses the fundamental problem of induction: \textbf{How
can we know the label of a point we have never seen?}

\begin{enumerate}
\def\labelenumi{\arabic{enumi}.}
\tightlist
\item
  \textbf{Unseen Data is Arbitrary:} Suppose the learner observes a
  training set \(S\) covering half the domain. For the remaining half of
  the domain (unseen points), the true labels could be \emph{anything}.
\item
  \textbf{No Pattern Without Assumptions:} Without prior assumptions
  (inductive bias), all possible labelings of the unseen points are
  equally likely.
\item
  \textbf{Random Guessing:} Therefore, on the unseen points, any
  algorithm \(\mathcal{A}\) is effectively guessing. Averaged over all
  possible target functions, the error rate on unseen data will be 0.5
  (like a coin toss).
\end{enumerate}

This implies that purely data-driven learning (without assumptions about
the shape or structure of \(f\)) is impossible.

\subsection{Relative Superiority of
Learners}\label{relative-superiority-of-learners}

\begin{quote}
``For every learner, there exists a task on which it fails, even though
that task can be successfully learned by another learner.''
\end{quote}

This means:

\begin{itemize}
\tightlist
\item
  There is no ``Master Algorithm'' that is superior to all other
  algorithms on all possible distributions.
\item
  If Algorithm A outperforms Algorithm B on a certain set of tasks
  (e.g., image recognition), there must exist a different set of tasks
  (e.g., noise patterns) where Algorithm B outperforms Algorithm A.
\item
  \textbf{Success depends on the fit between the algorithm's inductive
  bias and the actual task structure.}
\end{itemize}

\subsection{Mathematical Sketch}\label{mathematical-sketch}

Consider a finite domain \(\mathcal{X}\).

\begin{enumerate}
\def\labelenumi{\arabic{enumi}.}
\tightlist
\item
  We fix a training set \(S\).
\item
  We consider the set of all possible target functions
  \(F = \{f : \mathcal{X} \to \{0,1\}\}\).
\item
  For any point \(x \notin S\), exactly half of the functions in \(F\)
  label \(x\) as 1, and the other half label it as 0.
\item
  Regardless of what hypothesis \(h\) the algorithm outputs based on
  \(S\), \(h(x)\) will disagree with exactly 50\% of the possible target
  functions.
\item
  Therefore, averaged over all possible target functions, the error is
  high.
\end{enumerate}

\subsection{Key Takeaways}\label{key-takeaways-2}

\begin{enumerate}
\def\labelenumi{\arabic{enumi}.}
\tightlist
\item
  \textbf{Universal Learning is Impossible:} We cannot learn a model
  that guarantees good performance on \emph{every} possible problem.
\item
  \textbf{Inductive Bias is Mandatory:} To learn successfully, we must
  make assumptions about the data (e.g., ``the data is linearly
  separable'' or ``the data is smooth'').
\item
  \textbf{Context Matters:} We cannot say an algorithm is ``good'' in
  isolation; we can only say it is good for a specific class of problems
  (e.g., Convolutional Neural Networks are good for \emph{images}, not
  necessarily for \emph{tabular financial data}).
\item
  \textbf{No Free Lunch:} You cannot get a low error rate ``for free''
  (without leveraging specific properties of the target distribution).
\end{enumerate}

\newpage

\section{Error Decomposition}\label{error-decomposition}

To understand the performance of a learning algorithm, it is useful to
\textbf{decompose} the generalization error (True Risk) into distinct
components. This decomposition helps diagnose whether a model is
suffering from underfitting or overfitting.

Let \(h_S\) be the hypothesis returned by an ERM algorithm operating on
a sample \(S\) using hypothesis class \(\mathcal{H}\). Let \(h^*\) be
the best possible hypothesis within that class (the one that minimizes
the true risk):
\[ h^* \in \underset{h \in \mathcal{H}}{\text{argmin}} \ L_{\mathcal{D}}(h) \]

We can decompose the error of our learned hypothesis
\(L_{\mathcal{D}}(h_S)\) as follows:

\[ L_{\mathcal{D}}(h_S) = \underbrace{L_{\mathcal{D}}(h^*)}_{\text{Approximation Error}} + \underbrace{(L_{\mathcal{D}}(h_S) - L_{\mathcal{D}}(h^*))}_{\text{Estimation Error}} \]

\subsection{\texorpdfstring{1. The Approximation Error
(\(\epsilon_{\text{app}}\))}{1. The Approximation Error (\textbackslash epsilon\_\{\textbackslash text\{app\}\})}}\label{the-approximation-error-epsilon_textapp}

The Approximation Error measures how much risk we incur because we
restricted ourselves to a specific hypothesis class \(\mathcal{H}\),
rather than searching the entire universe of possible functions.

\[ \epsilon_{\text{app}} = L_{\mathcal{D}}(h^*) \] \emph{(Note: If the
true labeling function \(f\) is deterministic and noise-free, the ideal
error is 0. Thus, this term represents the gap between the best possible
model in \(\mathcal{H}\) and the truth.)}

\begin{itemize}
\tightlist
\item
  \textbf{Cause:} This error arises from the \textbf{Inductive Bias}. If
  \(\mathcal{H}\) is not ``rich'' or complex enough to contain the true
  target function \(f\), even the best possible hypothesis \(h^*\) will
  have non-zero error.
\item
  \textbf{Intuition:} This corresponds to \textbf{Underfitting}. If you
  try to separate a circular pattern using only straight lines (linear
  classifiers), you will have high approximation error regardless of how
  much data you have.
\item
  \textbf{Dependency:} It depends on the choice of the hypothesis class
  \(\mathcal{H}\), but it is \textbf{independent} of the training data
  size \(m\).
\end{itemize}

\subsection{\texorpdfstring{2. The Estimation Error
(\(\epsilon_{\text{est}}\))}{2. The Estimation Error (\textbackslash epsilon\_\{\textbackslash text\{est\}\})}}\label{the-estimation-error-epsilon_textest}

The Estimation Error measures the additional risk incurred because we do
not know the true distribution \(\mathcal{D}\) and must rely on a finite
training sample \(S\) to choose our hypothesis.

\[ \epsilon_{\text{est}} = L_{\mathcal{D}}(h_S) - L_{\mathcal{D}}(h^*) \]

\begin{itemize}
\tightlist
\item
  \textbf{Cause:} Because the training set \(S\) is a random sample, the
  empirical error \(L_S(h)\) is only an \emph{estimate} of the true
  error \(L_{\mathcal{D}}(h)\). The ERM algorithm minimizes the
  estimate, not the true value. Consequently, \(h_S\) is likely strictly
  worse than \(h^*\).
\item
  \textbf{Intuition:} This corresponds to \textbf{Overfitting} (or
  variance). We chose \(h_S\) because it looked good on the training
  data, but it might not be the true best predictor.
\item
  \textbf{Dependency:} This error decreases as the sample size \(m\)
  increases (by the Law of Large Numbers). However, it generally
  increases as the complexity of \(\mathcal{H}\) increases (more complex
  classes are harder to estimate).
\end{itemize}

\subsection{The Bias-Complexity
Tradeoff}\label{the-bias-complexity-tradeoff}

{\def\LTcaptype{none} % do not increment counter
\begin{longtable}[]{@{}
  >{\raggedright\arraybackslash}p{(\linewidth - 4\tabcolsep) * \real{0.3333}}
  >{\raggedright\arraybackslash}p{(\linewidth - 4\tabcolsep) * \real{0.3333}}
  >{\raggedright\arraybackslash}p{(\linewidth - 4\tabcolsep) * \real{0.3333}}@{}}
\toprule\noalign{}
\begin{minipage}[b]{\linewidth}\raggedright
Property
\end{minipage} & \begin{minipage}[b]{\linewidth}\raggedright
Small / Simple \(\mathcal{H}\)
\end{minipage} & \begin{minipage}[b]{\linewidth}\raggedright
Large / Complex \(\mathcal{H}\)
\end{minipage} \\
\midrule\noalign{}
\endhead
\bottomrule\noalign{}
\endlastfoot
\textbf{Example} & Linear Models, Shallow Trees & Deep Neural Networks,
High-degree Polynomials \\
\textbf{Approximation Error} & \textbf{High} (Underfitting) \newline The
model is too rigid to capture reality. & \textbf{Low} \newline The model
is flexible enough to fit \(f\). \\
\textbf{Estimation Error} & \textbf{Low} \newline Easy to find the best
\(h\) with few samples. & \textbf{High} (Overfitting) \newline Hard to
distinguish signal from noise. \\
\textbf{Total Error} & Dominated by Bias & Dominated by Variance \\
\end{longtable}
}

\subsubsection{Summary of the Tradeoff}\label{summary-of-the-tradeoff}

\begin{itemize}
\tightlist
\item
  \textbf{If \(\mathcal{H}\) is too small:} We fail to learn the pattern
  (\(\epsilon_{\text{app}}\) is large).
\item
  \textbf{If \(\mathcal{H}\) is too large:} We memorize the noise in the
  specific sample \(S\) (\(\epsilon_{\text{est}}\) is large).
\item
  \textbf{Goal:} The goal of model selection is to find the ``sweet
  spot'' for \(\mathcal{H}\) that minimizes the sum
  \(\epsilon_{\text{app}} + \epsilon_{\text{est}}\).
\end{itemize}

\subsection{Key Takeaways}\label{key-takeaways-3}

\begin{enumerate}
\def\labelenumi{\arabic{enumi}.}
\tightlist
\item
  \textbf{Approximation Error:} The limit of your chosen model class.
  Fixed by design choices.
\item
  \textbf{Estimation Error:} The penalty for having finite data. Reduces
  with more data.
\item
  \textbf{Decomposition:} Total Error = (How far the best model is from
  the truth) + (How far your learned model is from the best model).
\item
  \textbf{Optimization Note:} In practice, there is often a third term,
  \textbf{Optimization Error}, if the algorithm fails to find the exact
  minimum of the training error (e.g., getting stuck in a local minimum
  during Gradient Descent).
  \[ L_{\mathcal{D}}(h_{\text{final}}) = \epsilon_{\text{app}} + \epsilon_{\text{est}} + \epsilon_{\text{opt}} \]
\end{enumerate}

\newpage

\section{The VC-Dimension}\label{the-vc-dimension}

Our current goal is to determine which hypothesis classes
\(\mathcal{H}\) are PAC learnable. Previously, we established that
finite classes are learnable with a sample complexity depending on
\(\log(|\mathcal{H}|)\). However, most practical classes (e.g., linear
classifiers, neural networks) contain an infinite number of hypotheses.

\subsection{1. Why Infinite Classes Can Be
Learnable}\label{why-infinite-classes-can-be-learnable}

It may seem intuitive that if \(\mathcal{H}\) is infinite, learning is
impossible because the search space is unbounded. However, this is
incorrect. The ``size'' (cardinality) of the class is a crude measure of
its complexity. What matters is the class's \textbf{expressive power} on
a finite set of data points.

\subsubsection{The Concept of ``Effective''
Size}\label{the-concept-of-effective-size}

Consider the class of threshold functions on the \textbf{real line}:
\(h_t(x) = 1\) if \(x > t\), else \(0\). * There are infinitely many
real numbers \(t\), so \(|\mathcal{H}| = \infty\). * However, given a
dataset of \(m\) points, there are at most \(m+1\) ways to label them
using this class (placing the threshold between points). * Even though
the \emph{parameters} are continuous, the number of distinct
\textbf{behaviors} (dichotomies) on a finite sample is small and finite.

This insight leads to the Vapnik-Chervonenkis (VC) theory: we measure
complexity not by counting functions, but by measuring how
``shatterable'' data is.

\subsection{2. Shattering}\label{shattering}

To define VC-Dimension, we first need the concept of
\textbf{Shattering}.

Let \(C = \{c_1, \dots, c_m\}\) be a finite set of instances from the
domain \(\mathcal{X}\). The restriction of \(\mathcal{H}\) to \(C\),
denoted \(\mathcal{H}_C\), is the set of all possible labelings that
hypotheses in \(\mathcal{H}\) can assign to \(C\).

\textbf{Definition:} A set \(C\) is \textbf{shattered} by
\(\mathcal{H}\) if \(\mathcal{H}\) can realize \emph{all possible} label
configurations on \(C\). \[ |\mathcal{H}_C| = 2^{|C|} \]

\textbf{Intuition:} If a set of points is shattered, it means the
hypothesis class is ``rich'' enough to classify these specific points in
any arbitrary way (e.g., all 0s, all 1s, alternating, random noise,
etc.).

\subsection{3. The VC-Dimension
Definition}\label{the-vc-dimension-definition}

The \textbf{VC-Dimension} of a hypothesis class \(\mathcal{H}\), denoted
as \(\text{VCdim}(\mathcal{H})\), is the size of the largest finite
subset of \(\mathcal{X}\) that can be shattered by \(\mathcal{H}\).

\textbf{Formal Definition:}
\[ \text{VCdim}(\mathcal{H}) = \sup \{ |C| : C \subset \mathcal{X} \text{ and } C \text{ is shattered by } \mathcal{H} \} \]

\textbf{To show that \(\text{VCdim}(\mathcal{H}) = d\), you must prove
two things:} 1. \textbf{Lower Bound:} There \textbf{exists} at least one
set of size \(d\) that is shattered by \(\mathcal{H}\). 2. \textbf{Upper
Bound:} \textbf{No} set of size \(d+1\) can be shattered by
\(\mathcal{H}\).

If \(\mathcal{H}\) can shatter sets of arbitrarily large size, then
\(\text{VCdim}(\mathcal{H}) = \infty\).

\subsection{4. Intuitive Meaning of
VC-Dimension}\label{intuitive-meaning-of-vc-dimension}

The VC-Dimension measures the \textbf{capacity} or \textbf{complexity}
of a learning model.

\subsubsection{``Falsifiability'' and
Generalization}\label{falsifiability-and-generalization}

\begin{itemize}
\tightlist
\item
  If \(\text{VCdim}(\mathcal{H})\) is high, the model is very flexible.
  It can fit complex patterns, but it can also fit pure noise.
\item
  If \(\text{VCdim}(\mathcal{H})\) is low, the model is rigid (strong
  inductive bias).
\end{itemize}

\textbf{The Combinatorial Cliff:} Imagine the VC-dimension is \(d\). *
For \(m \le d\) points, the model can essentially ``memorize'' any
labeling. We cannot distinguish between true learning and memorization
here. * For \(m > d\) points, the model \textbf{cannot} fit every
possible labeling. The ``growth function'' (number of possible
behaviors) stops growing exponentially (\(2^m\)) and starts growing
polynomially (\(m^d\))---this is known as Sauer's Lemma. *
\textbf{Crucial Insight:} Once the number of data points exceeds the
VC-dimension, the model is \emph{forced} to find structure because it
lacks the capacity to just memorize random labels. This restriction is
what enables generalization.

\subsection{5. Examples}\label{examples}

{\def\LTcaptype{none} % do not increment counter
\begin{longtable}[]{@{}
  >{\raggedright\arraybackslash}p{(\linewidth - 6\tabcolsep) * \real{0.2500}}
  >{\raggedright\arraybackslash}p{(\linewidth - 6\tabcolsep) * \real{0.2500}}
  >{\raggedright\arraybackslash}p{(\linewidth - 6\tabcolsep) * \real{0.2500}}
  >{\raggedright\arraybackslash}p{(\linewidth - 6\tabcolsep) * \real{0.2500}}@{}}
\toprule\noalign{}
\begin{minipage}[b]{\linewidth}\raggedright
Hypothesis Class
\end{minipage} & \begin{minipage}[b]{\linewidth}\raggedright
Domain
\end{minipage} & \begin{minipage}[b]{\linewidth}\raggedright
VC-Dimension
\end{minipage} & \begin{minipage}[b]{\linewidth}\raggedright
Explanation
\end{minipage} \\
\midrule\noalign{}
\endhead
\bottomrule\noalign{}
\endlastfoot
: & : & : & : \\
\textbf{Thresholds} & \(\mathbb{R}\) & \textbf{1} & Can shatter 1 point
(can be 0 or 1). Cannot shatter 2 points (cannot label them \(1, 0\) if
\(x_1 < x_2\)). \\
\textbf{Intervals} & \(\mathbb{R}\) & \textbf{2} & Can capture
\([a,b]\). Can shatter 2 points. Cannot shatter 3 points (cannot label
\(1, 0, 1\)). \\
\textbf{Linear Classifiers} & \(\mathbb{R}^2\) & \textbf{3} & Can
shatter 3 non-collinear points. Cannot shatter 4 points (the XOR problem
requires non-linear bounds). \\
\textbf{Linear Classifiers} & \(\mathbb{R}^d\) & \textbf{d + 1} & In
general, hyperplanes in \(d\) dimensions have VC-dim \(d+1\). \\
\textbf{Sine Function} & \(\mathbb{R}\) & \textbf{\(\infty\)} &
\(h_\theta(x) = \text{sign}(\sin(\theta x))\). Can oscillate infinitely
fast to fit any binary sequence. \\
\end{longtable}
}

\subsection{Key Takeaways}\label{key-takeaways-4}

\begin{enumerate}
\def\labelenumi{\arabic{enumi}.}
\tightlist
\item
  \textbf{Infinite Classes:} Learnability depends on the ``effective''
  combinatorial complexity, not the number of parameters.
\item
  \textbf{Shattering:} The ability to classify a set of points in all
  \(2^m\) possible ways.
\item
  \textbf{VC-Dimension:} The maximum number of points the class can
  shatter.
\item
  \textbf{Fundamental Theorem of Statistical Learning:} A hypothesis
  class \(\mathcal{H}\) is PAC learnable \textbf{if and only if} its
  VC-dimension is finite.
\item
  \textbf{Sample Complexity:} The number of samples required to learn is
  proportional to the VC-dimension:
  \(m \propto \frac{d + \log(1/\delta)}{\epsilon}\).
\end{enumerate}

\newpage

\section{Perceptron for Halfspaces}\label{perceptron-for-halfspaces}

The Perceptron, introduced by Frank Rosenblatt in 1958, is one of the
earliest and most fundamental algorithms for Supervised Learning. It is
an iterative algorithm designed to learn a \textbf{Linear Classifier}
(or Halfspace) for binary classification tasks where
\(\mathcal{Y} = \{-1, +1\}\).

\subsection{1. The Setting: Halfspaces}\label{the-setting-halfspaces}

We focus on the hypothesis class of \textbf{Linear Halfspaces} in
\(\mathbb{R}^d\). A hypothesis \(h_w\) is parameterized by a weight
vector \(w \in \mathbb{R}^d\) and (optionally) a bias term \(b\).
\[ h_{w,b}(x) = \text{sign}(\langle w, x \rangle + b) \] \emph{Note: We
often use the ``bias trick'' by appending a 1 to the input vector
\(x \to [x, 1]\) and absorbing \(b\) into \(w\), reducing the problem to
learning a homogeneous halfspace passing through the origin:
\(h_w(x) = \text{sign}(\langle w, x \rangle)\).}

\textbf{Goal:} Find a vector \(w\) that separates the data perfectly (if
possible). This is equivalent to finding a \(w\) such that for all
training examples \((x_i, y_i)\): \[ y_i \langle w, x_i \rangle > 0 \]

\subsection{2. The Perceptron Algorithm}\label{the-perceptron-algorithm}

The Perceptron is an implementation of \textbf{Empirical Risk
Minimization (ERM)} for linearly separable data. It attempts to find a
separating hyperplane by iteratively correcting errors.

\subsubsection{The Algorithm Steps}\label{the-algorithm-steps}

\begin{enumerate}
\def\labelenumi{\arabic{enumi}.}
\tightlist
\item
  \textbf{Initialization:} Start with the all-zero vector
  \(w^{(1)} = (0, \dots, 0)\).
\item
  \textbf{Iterative Updates:} At each step \(t\), the algorithm looks
  for a sample \((x_i, y_i)\) that is \textbf{mislabeled} by the current
  weight vector \(w^{(t)}\).

  \begin{itemize}
  \tightlist
  \item
    Condition for misclassification:
    \(\text{sign}(\langle w^{(t)}, x_i \rangle) \neq y_i\) (or
    equivalently \(y_i \langle w^{(t)}, x_i \rangle \leq 0\)).
  \end{itemize}
\item
  \textbf{Update Rule:} If such a point is found, update the weight
  vector: \[ w^{(t+1)} \leftarrow w^{(t)} + y_i x_i \]
\end{enumerate}

\subsubsection{Intuition of the Update}\label{intuition-of-the-update}

\begin{itemize}
\tightlist
\item
  If \(y_i = +1\) and we predicted \(-1\), we add \(x_i\) to \(w\). This
  rotates \(w\) \emph{towards} \(x_i\), increasing the dot product
  \(\langle w, x_i \rangle\).
\item
  If \(y_i = -1\) and we predicted \(+1\), we subtract \(x_i\). This
  rotates \(w\) \emph{away} from \(x_i\), decreasing the dot product.
\end{itemize}

!{[}{[}Pasted image 20260215181640.png \textbar{} 400{]}{]}

\subsection{3. The Batch Perceptron
Pseudo-Code}\label{the-batch-perceptron-pseudo-code}

While the ``Online Perceptron'' processes data one by one from a stream,
the \textbf{Batch Perceptron} iterates over a fixed training set \(S\)
until convergence.

\begin{Shaded}
\begin{Highlighting}[]
\AnnotationTok{Algorithm:}\CommentTok{ Batch Perceptron}

\NormalTok{Input: }
\SpecialStringTok{  {-} }\NormalTok{Training data S = \{(x\_1, y\_1), ..., (x\_m, y\_m)\} where x\_i \textbackslash{}in R\^{}d, y\_i \textbackslash{}in \{{-}1, +1\}}
\SpecialStringTok{  {-} }\NormalTok{Max epochs T (optional, to prevent infinite loops if data is not separable)}

\NormalTok{Initialize:}
\NormalTok{  w ← (0, 0, ..., 0)  // Zero vector of dimension d}
\NormalTok{  t ← 0}

\NormalTok{Loop:}
\NormalTok{  mistake\_found ← FALSE}
  
\NormalTok{  For i = 1 to m:}
\NormalTok{    // Check if x\_i is misclassified}
\NormalTok{    If y\_i * (dot\_product(w, x\_i)) \textless{}= 0:}
      
\InformationTok{      // Update rule}
\InformationTok{      w ← w + y\_i * x\_i}
\InformationTok{      }
\InformationTok{      mistake\_found ← TRUE}
\InformationTok{      t ← t + 1}
\InformationTok{      }
\NormalTok{  // Termination condition}
\NormalTok{  If mistake\_found == FALSE:}
\NormalTok{    Break Loop (Convergence reached: all points correctly classified)}

\NormalTok{Output: w}
\end{Highlighting}
\end{Shaded}

\subsubsection{Convergence Theorem
(Novikoff)}\label{convergence-theorem-novikoff}

If the training data \(S\) is linearly separable with a margin
\(\gamma\) (i.e., there exists a \(w^*\) such that
\(y_i \langle w^*, x_i \rangle \ge \gamma\) for all \(i\)), then the
Perceptron is guaranteed to converge in a finite number of steps,
specifically at most \((R/\gamma)^2\) updates, where
\(R = \max_i \|x_i\|\).

\subsection{4. VC Dimension of
Halfspaces}\label{vc-dimension-of-halfspaces}

To understand the learnability of linear classifiers, we analyze their
VC-Dimension.

\subsubsection{Theorem}\label{theorem}

The VC-dimension of the class of Affine Halfspaces in \(\mathbb{R}^d\)
is \textbf{\(d + 1\)}.
\[ \text{VCdim}(\mathcal{H}_{\text{halfspaces}}) = d + 1 \] \emph{(Note:
If we restrict to homogeneous halfspaces passing through the origin,
VC-dim is \(d\).)}

\subsubsection{Intuitive Proof Sketch}\label{intuitive-proof-sketch}

\begin{enumerate}
\def\labelenumi{\arabic{enumi}.}
\item
  \textbf{Lower Bound (Why \(\ge d+1\)):} We need to show there exists a
  set of \(d+1\) points that can be shattered. Consider the set of basis
  vectors in \(\mathbb{R}^d\) plus the origin (e.g., in 2D, a triangle
  of points). It is geometrically possible to draw a line separating any
  subset of these points from the others. Thus, we can realize all
  \(2^{d+1}\) labelings.
\item
  \textbf{Upper Bound (Why \(< d+2\)):} We need to show that \textbf{no}
  set of \(d+2\) points can be shattered. This relies on Radon's Theorem
  or linear dependence. In a \(d\)-dimensional space, any set of \(d+2\)
  points must be linearly dependent.

  \emph{Example in \(\mathbb{R}^2\) (2D Plane):} Take any 4 points
  (\(d+2\)).

  \begin{itemize}
  \tightlist
  \item
    Case A: One point is inside the triangle formed by the other three.
    If we label the outer points \(+1\) and the inner point \(-1\), no
    linear line can separate them (XOR problem).
  \item
    Case B: The points form a convex quadrilateral. If we label diagonal
    opposites as \(+1\) and the other pair as \(-1\), no line can
    separate them. Since there is at least one labeling we \emph{cannot}
    achieve, the set is not shattered.
  \end{itemize}
\end{enumerate}

\subsubsection{Key Takeaways}\label{key-takeaways-5}

\begin{itemize}
\tightlist
\item
  \textbf{Perceptron:} A simple, guaranteed algorithm for finding linear
  separators if they exist.
\item
  \textbf{Complexity:} The complexity of learning linear models is
  controlled by the dimension of the input space \(d\).
\item
  \textbf{Sample Complexity:} Since \(\text{VCdim} = d+1\), the number
  of samples needed to learn a linear classifier grows linearly with the
  number of features \(d\).
\end{itemize}

\newpage

\section{Linear Regression}\label{linear-regression}

\subsection{1. Introduction}\label{introduction}

Linear Regression is a fundamental supervised learning algorithm used to
predict a \textbf{continuous} target variable \(y \in \mathbb{R}\) based
on a feature vector \(x \in \mathbb{R}^d\). This contrasts with
classification, where the target is discrete.

\subsubsection{The Hypothesis Class}\label{the-hypothesis-class}

We assume the relationship between input and output is linear (or
affine). The hypothesis class \(\mathcal{H}_{reg}\) consists of
functions of the form: \[ h_w(x) = \langle w, x \rangle + b \] where
\(w \in \mathbb{R}^d\) is the weight vector and \(b \in \mathbb{R}\) is
the bias. To simplify notation, we use the ``bias trick'': append a 1 to
each \(x\) (\(x \to [x, 1]\)) and include \(b\) in \(w\)
(\(w \to [w, b]\)). Then: \[ h_w(x) = \langle w, x \rangle \]

\subsubsection{The Loss Function: Squared
Error}\label{the-loss-function-squared-error}

To evaluate how well a hypothesis fits the data, we need a loss function
that \emph{penalizes deviations} from the true value. The most common
choice is the \textbf{Squared Error Loss}:
\[ \ell(h_w(x), y) = (h_w(x) - y)^2 = (\langle w, x \rangle - y)^2 \]

\subsubsection{Empirical Risk Minimization
(ERM)}\label{empirical-risk-minimization-erm-1}

The goal of Linear Regression is to find the vector \(w\) that minimizes
the Empirical Risk (Mean Squared Error) over the training set
\(S = \{(x_1, y_1), \dots, (x_m, y_m)\}\).

\[ L_S(h_w) = \frac{1}{m} \sum_{i=1}^{m} (\langle w, x_i \rangle - y_i)^2 \]

\[ w^* = \underset{w}{\text{argmin}} \ L_S(h_w) \]

\subsection{2. Least Squares (Matrix
Form)}\label{least-squares-matrix-form}

To solve the minimization problem efficiently, we can rewrite the
objective function using matrix notation. This formulation is often
called \textbf{Ordinary Least Squares (OLS)}.

\subsubsection{Matrix Definitions}\label{matrix-definitions}

Let \(m\) be the number of samples and \(d\) be the number of features.

\begin{enumerate}
\def\labelenumi{\arabic{enumi}.}
\item
  \textbf{Design Matrix (\(X\)):} An \(m \times d\) matrix where each
  row \(i\) is the feature vector \(x_i^\top\).
  \[ X = \begin{pmatrix} x_1^\top \\ x_2^\top \\ \vdots \\ x_m^\top \end{pmatrix} \in \mathbb{R}^{m \times d} \]
\item
  \textbf{Target Vector (\(y\)):} An \(m \times 1\) column vector
  containing all target values.
  \[ y = \begin{pmatrix} y_1 \\ y_2 \\ \vdots \\ y_m \end{pmatrix} \in \mathbb{R}^m \]
\item
  \textbf{Weight Vector (\(w\)):} A \(d \times 1\) column vector of
  parameters.
  \[ w = \begin{pmatrix} w_1 \\ \vdots \\ w_d \end{pmatrix} \in \mathbb{R}^d \]
\end{enumerate}

\subsubsection{The Objective in Matrix
Form}\label{the-objective-in-matrix-form}

The vector of \textbf{predictions for all training points} is given by
\(Xw\). The error vector (residuals) is \(Xw - y\). The sum of squared
errors can be written as the \textbf{squared Euclidean norm} of the
residual vector:

\[ L_S(w) = \frac{1}{m} \| Xw - y \|^2 = \frac{1}{m} (Xw - y)^\top (Xw - y) \]

We want to find \(w\) that minimizes \(\| Xw - y \|^2\).

\subsection{3. Derivation of the Normal
Equation}\label{derivation-of-the-normal-equation}

To find the minimum, we take the gradient of the objective function with
respect to \(w\) and set it to zero.

\subsubsection{Gradient Derivation}\label{gradient-derivation}

Let \(f(w) = \| Xw - y \|^2 = (Xw - y)^\top (Xw - y)\). Expand the
terms: \[ f(w) = (w^\top X^\top - y^\top)(Xw - y) \]
\[ f(w) = w^\top X^\top X w - w^\top X^\top y - y^\top X w + y^\top y \]
Since \(w^\top X^\top y\) is a scalar, it is equal to its transpose
\((w^\top X^\top y)^\top = y^\top X w\).
\[ f(w) = w^\top (X^\top X) w - 2 (X^\top y)^\top w + y^\top y \]

Now, calculate the gradient \(\nabla_w f(w)\): 1. The derivative of
\(w^\top A w\) with respect to \(w\) is \(2Aw\) (assuming A is
symmetric, which \(X^\top X\) is). 2. The derivative of \(b^\top w\)
with respect to \(w\) is \(b\).

\[ \nabla_w f(w) = 2 X^\top X w - 2 X^\top y \]

\subsubsection{Setting Gradient to Zero}\label{setting-gradient-to-zero}

To find the optimum \(w^*\): \[ 2 X^\top X w^* - 2 X^\top y = 0 \]
\[ X^\top X w^* = X^\top y \]

This is known as the \textbf{Normal Equation}.

\subsubsection{The Solution}\label{the-solution}

Assuming \(X^\top X\) is invertible (which requires full rank, i.e.,
features are linearly independent): \[ w^* = (X^\top X)^{-1} X^\top y \]

\subsection{4. Geometric Interpretation: The Projection
Theorem}\label{geometric-interpretation-the-projection-theorem}

The least squares solution has a beautiful geometric interpretation.

\begin{enumerate}
\def\labelenumi{\arabic{enumi}.}
\tightlist
\item
  \textbf{Column Space:} The vector of predictions \(\hat{y} = Xw\) is a
  linear combination of the columns of \(X\). Thus, \(\hat{y}\) must lie
  in the \textbf{Column Space} (Range) of \(X\), denoted
  \(\text{Col}(X)\).
\item
  \textbf{Target Vector:} The true target vector \(y\) generally does
  not lie in \(\text{Col}(X)\) (due to noise or model mismatch).
\item
  \textbf{Closest Point:} We want to find the point
  \(\hat{y} \in \text{Col}(X)\) that is closest to \(y\) in Euclidean
  distance.
\end{enumerate}

\pandocbounded{\includegraphics[keepaspectratio]{/home/nicola/Documents/Local Obs/Obsidian Locale/images/1445268e3ac976688c6da3932b974f1c3d153e21.png}}

\textbf{The Orthogonality Principle:} The distance \(\| \hat{y} - y \|\)
is minimized when the residual vector \(r = y - \hat{y}\) is
\textbf{orthogonal} (perpendicular) to the subspace \(\text{Col}(X)\).
This means the residual must be orthogonal to every column of \(X\).

Mathematically: \[ X^\top (y - Xw) = 0 \]
\[ X^\top y - X^\top X w = 0 \] \[ X^\top X w = X^\top y \]

This recovers the Normal Equation!

\textbf{Intuition:} The optimal prediction \(\hat{y} = Xw^*\) is the
\textbf{Orthogonal Projection} of the observed vector \(y\) onto the
subspace spanned by the features \(X\). The matrix
\(P = X(X^\top X)^{-1}X^\top\) is essentially a projection matrix.

\subsection{Key Takeaways}\label{key-takeaways-6}

\begin{enumerate}
\def\labelenumi{\arabic{enumi}.}
\tightlist
\item
  \textbf{Model:} Linear function \(h(x) = \langle w, x \rangle\).
\item
  \textbf{Loss:} Squared Loss (convex and differentiable).
\item
  \textbf{Solution:} Closed-form solution exists via the Normal
  Equation: \(w^* = (X^\top X)^{-1} X^\top y\).
\item
  \textbf{Geometry:} Least Squares projects \(y\) onto the column space
  of \(X\) to minimize the residual norm.
\item
  \textbf{Uniqueness:} The solution is unique if \(X^\top X\) is
  invertible. If features are collinear, regularization (like Ridge
  Regression) is needed.
\end{enumerate}

\newpage

\section{Regularized Linear
Regression}\label{regularized-linear-regression}

Standard Linear Regression (Ordinary Least Squares - OLS) minimizes the
empirical risk \(L_S(w) = \|Xw - y\|^2\). However, OLS often suffers
from \textbf{overfitting}, particularly when: 1. The number of features
\(d\) is large relative to the number of samples \(m\) (\(d \approx m\)
or \(d > m\)). 2. The features are highly correlated
(multicollinearity), causing the matrix \(X^\top X\) to be close to
singular, which makes the weights \(w\) unstable and very large.

To address this, we apply \textbf{Regularization}, which constrains the
magnitude of the weight vector. This introduces a \textbf{Bias-Variance
Tradeoff}: we accept a small amount of bias (deviation from the training
data) to significantly reduce variance (sensitivity to noise), thereby
lowering the total Generalization Error.

\pandocbounded{\includegraphics[keepaspectratio]{/home/nicola/Documents/Local Obs/Obsidian Locale/images/bdde770ce8e8f43b914101085fac38920270d135.png}}

\subsection{1. Ridge Regression (L2
Regularization)}\label{ridge-regression-l2-regularization}

Ridge Regression adds an L2 penalty term to the Squared Error loss.

\subsubsection{The Objective Function}\label{the-objective-function}

\[ \min_{w} \ J(w) = \| Xw - y \|^2 + \lambda \| w \|_2^2 \]
\[ J(w) = \sum_{i=1}^{m} (\langle w, x_i \rangle - y_i)^2 + \lambda \sum_{j=1}^{d} w_j^2 \]
where \(\lambda \ge 0\) is the regularization hyperparameter.

\subsubsection{Closed-Form Solution}\label{closed-form-solution}

Unlike Logistic Regression, Ridge Regression has an exact analytic
solution. We take the gradient of \(J(w)\) with respect to \(w\) and set
it to zero: \[ \nabla_w J(w) = 2 X^\top (Xw - y) + 2 \lambda w = 0 \]
\[ X^\top X w + \lambda I w = X^\top y \]
\[ (X^\top X + \lambda I) w = X^\top y \]

The solution is:
\[ w_{\text{ridge}} = (X^\top X + \lambda I)^{-1} X^\top y \]

\subsubsection{Key Properties}\label{key-properties}

\begin{enumerate}
\def\labelenumi{\arabic{enumi}.}
\tightlist
\item
  \textbf{Invertibility:} The matrix \((X^\top X + \lambda I)\) is
  \textbf{always invertible} for \(\lambda > 0\), even if \(X^\top X\)
  is singular (non-invertible). This mathematically solves the problem
  of multicollinearity.
\item
  \textbf{Weight Shrinkage:} As \(\lambda \to \infty\), the penalty
  dominates, and \(\|w\| \to 0\). As \(\lambda \to 0\), we recover the
  OLS solution.
\item
  \textbf{Probabilistic View:} Ridge corresponds to MAP estimation
  assuming a Gaussian Prior on weights:
  \(w \sim \mathcal{N}(0, \tau^2 I)\).
\end{enumerate}

\subsection{2. Lasso Regression (L1
Regularization)}\label{lasso-regression-l1-regularization}

Lasso (Least Absolute Shrinkage and Selection Operator) adds an L1
penalty term.

\subsubsection{The Objective Function}\label{the-objective-function-1}

\[ \min_{w} \ J(w) = \| Xw - y \|^2 + \lambda \| w \|_1 \]
\[ J(w) = \sum_{i=1}^{m} (\langle w, x_i \rangle - y_i)^2 + \lambda \sum_{j=1}^{d} |w_j| \]

\subsubsection{Solution and Sparsity}\label{solution-and-sparsity}

\begin{enumerate}
\def\labelenumi{\arabic{enumi}.}
\tightlist
\item
  \textbf{No Closed Form:} Because \(|w_j|\) is not differentiable at 0,
  there is no simple formula like \((X^\top X + \lambda I)^{-1}\). We
  must use iterative optimization algorithms (e.g., Coordinate Descent
  or Subgradient Descent).
\item
  \textbf{Feature Selection:} Just like in Logistic Regression, Lasso
  drives many coefficients to \textbf{exactly zero}.

  \begin{itemize}
  \tightlist
  \item
    This implies that Lasso performs \textbf{automatic feature
    selection}.
  \item
    If \(w_j = 0\), the feature \(j\) is effectively removed from the
    model.
  \end{itemize}
\item
  \textbf{Geometric Interpretation:} The constraint region
  \(\sum |w_j| \le C\) is a polytope (diamond shape in 2D). The
  elliptical contours of the Squared Error loss function often hit the
  ``corners'' of this diamond first. At the corners, one or more
  coordinates are zero.
\end{enumerate}

\subsection{3. Data Normalization (Crucial
Step)}\label{data-normalization-crucial-step}

Before applying Regularization (L1 or L2), it is \textbf{mandatory} to
normalize or standardize the features (e.g., subtract mean, divide by
standard deviation).
\[ x_{ij} \leftarrow \frac{x_{ij} - \mu_j}{\sigma_j} \]

\textbf{Reason:} The penalty term \(\lambda \|w\|\) treats all weights
equally. If features have different scales (e.g., ``Age'' in years
vs.~``Income'' in dollars), their corresponding weights \(w_j\) will
have vastly different natural scales. Without normalization, the
regularization will heavily penalize the weights for small-scale
features (which require large \(w\) to have effect) while ignoring
large-scale features.

{\def\LTcaptype{none} % do not increment counter
\begin{longtable}[]{@{}
  >{\raggedright\arraybackslash}p{(\linewidth - 4\tabcolsep) * \real{0.3333}}
  >{\raggedright\arraybackslash}p{(\linewidth - 4\tabcolsep) * \real{0.3333}}
  >{\raggedright\arraybackslash}p{(\linewidth - 4\tabcolsep) * \real{0.3333}}@{}}
\toprule\noalign{}
\begin{minipage}[b]{\linewidth}\raggedright
Feature
\end{minipage} & \begin{minipage}[b]{\linewidth}\raggedright
Ridge Regression (L2)
\end{minipage} & \begin{minipage}[b]{\linewidth}\raggedright
Lasso Regression (L1)
\end{minipage} \\
\midrule\noalign{}
\endhead
\bottomrule\noalign{}
\endlastfoot
\textbf{Formula} & \(\|Xw-y\|^2 + \lambda \|w\|^2\) &
\(\|Xw-y\|^2 + \lambda \|w\|_1\) \\
\textbf{Solution} & Closed-form:
\((X^\top X + \lambda I)^{-1} X^\top y\) & Iterative optimization (no
formula) \\
\textbf{Sparsity} & Dense (all \(w_j \neq 0\)) & \textbf{Sparse} (many
\(w_j = 0\)) \\
\textbf{Use Case} & Multicollinearity; Better prediction when all
features matter. & Feature Selection; Interpretability. \\
\textbf{Prior} & Gaussian Prior & Laplace Prior \\
\end{longtable}
}

\subsection{Key Takeaways}\label{key-takeaways-7}

\begin{enumerate}
\def\labelenumi{\arabic{enumi}.}
\tightlist
\item
  \textbf{Bias-Variance Tradeoff:} Regularization increases Bias
  (training error) but decreases Variance (generalization error).
\item
  \textbf{Singularity:} Ridge Regression makes the normal equation
  solvable even when features are correlated.
\item
  \textbf{Sparsity:} Lasso produces sparse models, identifying the most
  important features.
\item
  \textbf{Hyperparameter \(\lambda\):} Must be tuned using
  Cross-Validation.

  \begin{itemize}
  \tightlist
  \item
    Too small \(\to\) Overfitting (OLS).
  \item
    Too large \(\to\) Underfitting (Constant predictor \(w=0\)).
  \end{itemize}
\end{enumerate}

\newpage

\section{9 Logistic Regression}\label{logistic-regression}

Logistic Regression is a discriminative probabilistic model used for
classification. Unlike Linear Regression, which outputs continuous
values, Logistic Regression outputs a probability score between 0 and 1,
representing the likelihood that a sample belongs to a particular class
(usually the ``positive'' class).

\subsection{9.1 The Sigmoid Function}\label{the-sigmoid-function}

To map the output of a linear function \(\langle w, x \rangle\) (which
ranges from \(-\infty\) to \(+\infty\)) to a probability (which must be
in \([0, 1]\)), we use the Sigmoid function (also known as the Logistic
function).

\subsubsection{9.1.1 Definition}\label{definition}

\[\sigma(z) = \frac{1}{1 + e^{-z}}\]

\subsubsection{9.1.2 Properties}\label{properties}

\begin{itemize}
\tightlist
\item
  \textbf{Range:} \(\sigma(z) \in (0, 1)\) for all \(z \in \mathbb{R}\).
\item
  \textbf{Symmetry:} \(\sigma(-z) = 1 - \sigma(z)\).
\item
  \textbf{Asymptotes:} \(\lim_{z \to \infty} \sigma(z) = 1\) and
  \(\lim_{z \to -\infty} \sigma(z) = 0\).
\item
  \textbf{Derivative:} The derivative can be expressed in terms of the
  function itself (useful for gradient descent):
  \[\sigma'(z) = \sigma(z)(1 - \sigma(z))\]
\end{itemize}

\subsubsection{9.1.3 The Model}\label{the-model}

We model the conditional probability of the label \(Y=1\) given input
\(x\) as:
\[P(Y=1 | x; w) = h_w(x) = \sigma(\langle w, x \rangle) = \frac{1}{1 + e^{-\langle w, x \rangle}}\]
Consequently, \(P(Y=0 | x; w) = 1 - h_w(x)\).

\subsection{9.2 ERM for Logistic
Regression}\label{erm-for-logistic-regression}

We cannot use the Squared Error loss for classification because it is
non-convex when applied to the sigmoid function, leading to local
minima. Instead, we use the Logistic Loss (or Log-Loss / Cross-Entropy
Loss).

\subsubsection{9.2.1 The Loss Function}\label{the-loss-function}

For a single example \((x, y)\) where \(y \in \{0, 1\}\):
\[\ell(h_w(x), y) = -y \log(h_w(x)) - (1-y) \log(1 - h_w(x))\] *
\textbf{Case y=1:} Loss is \(-\log(h_w(x))\). We want
\(h_w(x) \approx 1\). If \(h_w(x) \to 0\), \(\text{Loss} \to \infty\). *
\textbf{Case y=0:} Loss is \(-\log(1 - h_w(x))\). We want
\(h_w(x) \approx 0\). If \(h_w(x) \to 1\), \(\text{Loss} \to \infty\).

\subsubsection{9.2.2 The Empirical Risk}\label{the-empirical-risk}

The ERM rule seeks to minimize the average loss over the training set
\(S\):
\[L_S(w) = -\frac{1}{m} \sum_{i=1}^{m} \left[ y_i \log(h_w(x_i)) + (1-y_i) \log(1 - h_w(x_i)) \right]\]
This function is convex, ensuring that gradient-based optimization (like
Gradient Descent) will converge to the global minimum.

\subsection{9.3 Equivalence to Maximum Likelihood Estimation
(MLE)}\label{equivalence-to-maximum-likelihood-estimation-mle}

Why do we choose the Log-Loss? It arises naturally from the principle of
Maximum Likelihood Estimation.

\subsubsection{9.3.1 The Likelihood
Function}\label{the-likelihood-function}

Assume the training labels \(y_i\) are independent samples from a
Bernoulli distribution parameterized by \(h_w(x_i)\).

The probability of observing a single label \(y_i\) is:
\[P(y_i | x_i; w) = h_w(x_i)^{y_i} (1 - h_w(x_i))^{1 - y_i}\]

The Likelihood of observing the entire training set labels
\(Y = (y_1, \dots, y_m)\) given \(X\) is the product of individual
probabilities:
\[\mathcal{L}(w) = \prod_{i=1}^{m} P(y_i | x_i; w) = \prod_{i=1}^{m} h_w(x_i)^{y_i} (1 - h_w(x_i))^{1 - y_i}\]

\subsubsection{9.3.2 Maximizing
Log-Likelihood}\label{maximizing-log-likelihood}

To make the math easier (turning products into sums) and avoid numerical
underflow, we maximize the Log-Likelihood \(\ell(w)\):
\[\log(\mathcal{L}(w)) = \sum_{i=1}^{m} \left[ y_i \log(h_w(x_i)) + (1-y_i) \log(1 - h_w(x_i)) \right]\]

Notice that: \[\log(\mathcal{L}(w)) = -m \cdot L_S(w)\]

Therefore:
\[\underset{w}{\text{argmax}} \ \mathcal{L}(w) \iff \underset{w}{\text{argmin}} \ L_S(w)\]
\textbf{Conclusion:} Minimizing the Logistic Loss (ERM) is
mathematically equivalent to Maximizing the Likelihood of the data under
the assumed Bernoulli model.

\subsection{9.4 Connection to KL
Divergence}\label{connection-to-kl-divergence}

The Log-Loss can also be interpreted through the lens of Information
Theory, specifically using the Kullback-Leibler (KL) Divergence.

\subsubsection{9.4.1 KL Divergence
Definition}\label{kl-divergence-definition}

The KL Divergence measures the ``distance'' between two probability
distributions \(P\) (true distribution) and \(Q\) (approximated
distribution).
\[D_{KL}(P \| Q) = \sum_{x} P(x) \log \left( \frac{P(x)}{Q(x)} \right) = \mathbb{E}_{x \sim P} [\log P(x) - \log Q(x)]\]

\subsubsection{9.4.2 Relation to Logistic
Regression}\label{relation-to-logistic-regression}

In our context: * \(P\) is the ``true'' label distribution. For a
specific example \(x_i\), \(P(y=1) = y_i\) (deterministic). * \(Q\) is
the model's prediction. \(Q(y=1) = h_w(x_i)\).

The KL Divergence for a single sample is:
\[D_{KL}(y_i \| h_w(x_i)) = y_i \log \frac{y_i}{h_w(x_i)} + (1-y_i) \log \frac{1-y_i}{1-h_w(x_i)}\]

Since \(y_i \in \{0, 1\}\), the numerator terms \(y_i \log y_i\)
represent the entropy of the true label, which is \(0\) (since it's
deterministic). The expression simplifies to:
\[D_{KL} = - \left[ y_i \log h_w(x_i) + (1-y_i) \log (1-h_w(x_i)) \right]\]
This is exactly the Cross-Entropy Loss. Minimizing the Log-Loss is
equivalent to minimizing the KL Divergence between the empirical
distribution of labels and the model's predicted distribution.

\subsection{9.5 Properties of KL
Divergence}\label{properties-of-kl-divergence}

Although often called a ``distance,'' the KL Divergence is not a true
metric in the mathematical sense.

\begin{itemize}
\tightlist
\item
  \textbf{Non-Negativity (Gibbs' Inequality):} \[D_{KL}(P \| Q) \ge 0\]
  The divergence is \(0\) if and only if \(P = Q\) almost everywhere.
  This confirms that minimizing Log-Loss drives our predictions \(Q\)
  toward the truth \(P\).
\item
  \textbf{Asymmetry:} \[D_{KL}(P \| Q) \neq D_{KL}(Q \| P)\] The cost of
  approximating \(P\) with \(Q\) is different from approximating \(Q\)
  with \(P\). In Logistic Regression, we specifically minimize
  \(D_{KL}(\text{Truth} \| \text{Prediction})\).
\item
  \textbf{Triangle Inequality Fails:} It does not satisfy
  \(D_{KL}(P \| R) \le D_{KL}(P \| Q) + D_{KL}(Q \| R)\).
\end{itemize}

\subsection{9.6 Key Takeaways}\label{key-takeaways-8}

\begin{itemize}
\tightlist
\item
  \textbf{Sigmoid:} Maps linear outputs to probabilities
  qwen3.5:0.8b\((0, 1)\).
\item
  \textbf{Loss:} We minimize Cross-Entropy (Log-Loss) because it is
  convex.
\item
  \textbf{MLE:} ERM on Log-Loss is equivalent to finding the Maximum
  Likelihood Estimate for a Bernoulli distribution.
\item
  \textbf{KL Divergence:} Minimizing Log-Loss is equivalent to
  minimizing the information distance between the true labels and
  predicted probabilities.
\end{itemize}

\begin{center}\rule{0.5\linewidth}{0.5pt}\end{center}

\section{10 Regularization in Logistic
Regression}\label{regularization-in-logistic-regression}

While minimizing the Empirical Risk (Log-Loss) is a powerful strategy,
it can lead to Overfitting, especially when the data is linearly
separable. In such cases, the algorithm can drive the norm of the weight
vector \(\|w\| \to \infty\) to push the predicted probabilities
\(h_w(x)\) arbitrarily close to 0 or 1.

To prevent this and control the model complexity, we add a
Regularization Term \(R(w)\) to the loss function. This technique is
known as Regularized Empirical Risk Minimization.

The new objective function becomes:
\[\min_{w} \ J(w) = L_S(w) + \lambda R(w)\] where \(\lambda > 0\) is a
hyperparameter that controls the trade-off between fitting the data (low
\(L_S\)) and keeping the model simple (low \(R\)).

\subsection{10.1 L2 Regularization
(Ridge)}\label{l2-regularization-ridge}

L2 Regularization adds a penalty equal to the square of the magnitude of
coefficients. This is the most common form of regularization.

\subsubsection{10.1.1 The Objective}\label{the-objective}

\[J(w) = -\frac{1}{m} \sum_{i=1}^{m} \log P(y_i | x_i; w) + \lambda \|w\|_2^2\]
where \(\|w\|_2^2 = \sum_{j=1}^{d} w_j^2\).

\subsubsection{10.1.2 Effects and
Properties}\label{effects-and-properties}

\begin{itemize}
\tightlist
\item
  \textbf{Weight Decay:} The gradient of the regularization term is
  \(2\lambda w\). In gradient descent, this term constantly subtracts a
  fraction of the weight vector at each step
  (\(w \leftarrow w(1 - 2\eta\lambda) - \eta \nabla L_S\)). This shrinks
  weights toward zero but rarely makes them exactly zero.
\item
  \textbf{Prevents Divergence:} It ensures the optimization problem has
  a unique, finite solution even if the data is perfectly separable.
\item
  \textbf{Probabilistic Interpretation (MAP):} L2 regularization is
  equivalent to finding the Maximum A Posteriori (MAP) estimate where we
  assume a Gaussian Prior on the weights:
  \[w_j \sim \mathcal{N}(0, \sigma^2)\] Minimizing \(J(w)\) corresponds
  to maximizing the posterior \(P(w | S) \propto P(S | w) P(w)\).
\end{itemize}

\subsection{10.2 L1 Regularization
(Lasso)}\label{l1-regularization-lasso}

L1 Regularization adds a penalty equal to the sum of the absolute values
of the coefficients.

\subsubsection{10.2.1 The Objective}\label{the-objective-1}

\[J(w) = -\frac{1}{m} \sum_{i=1}^{m} \log P(y_i | x_i; w) + \lambda \|w\|_1\]
where \(\|w\|_1 = \sum_{j=1}^{d} |w_j|\).

\subsubsection{10.2.2 Effects and
Properties}\label{effects-and-properties-1}

\begin{itemize}
\tightlist
\item
  \textbf{Sparsity (Feature Selection):} The key feature of L1 is that
  it forces many weight coefficients to become exactly zero. This
  performs automatic feature selection, resulting in a sparse model that
  is easier to interpret.
\item
  \textbf{Geometric Intuition:} The constraint region for L1
  (\(\|w\|_1 \le C\)) is a ``diamond'' (polytope) with corners on the
  axes. The contours of the loss function \(L_S(w)\) are likely to touch
  the diamond at a corner first, where some coordinates are zero. In
  contrast, the L2 constraint is a sphere, and touching points usually
  happen away from the axes.
\item
  \textbf{Probabilistic Interpretation (MAP):} L1 regularization is
  equivalent to MAP estimation with a Laplacian Prior on the weights:
  \[P(w_j) = \frac{1}{2b} \exp\left(-\frac{|w_j|}{b}\right)\] The sharp
  peak of the Laplace distribution at zero encourages sparsity.
\end{itemize}

\subsection{10.3 Comparison and Key
Takeaways}\label{comparison-and-key-takeaways}

\pandocbounded{\includegraphics[keepaspectratio]{/home/nicola/Documents/Local Obs/Obsidian Locale/images/b795c07304c42cc254e12fb250fda72320f1c900.png}}

{\def\LTcaptype{none} % do not increment counter
\begin{longtable}[]{@{}
  >{\raggedright\arraybackslash}p{(\linewidth - 4\tabcolsep) * \real{0.3333}}
  >{\raggedright\arraybackslash}p{(\linewidth - 4\tabcolsep) * \real{0.3333}}
  >{\raggedright\arraybackslash}p{(\linewidth - 4\tabcolsep) * \real{0.3333}}@{}}
\toprule\noalign{}
\begin{minipage}[b]{\linewidth}\raggedright
Feature
\end{minipage} & \begin{minipage}[b]{\linewidth}\raggedright
L2 Regularization (Ridge)
\end{minipage} & \begin{minipage}[b]{\linewidth}\raggedright
L1 Regularization (Lasso)
\end{minipage} \\
\midrule\noalign{}
\endhead
\bottomrule\noalign{}
\endlastfoot
\textbf{Penalty Term} & \(\lambda \sum w_j^2\) & \(\lambda \sum\) \\
\textbf{Solution} & Non-sparse (weights are small but \(\neq 0\)) &
Sparse (many weights = \(0\)) \\
\textbf{Feature Selection} & No (keeps all features) & Yes (removes
irrelevant features) \\
\textbf{Differentiability} & Differentiable everywhere. &
Non-differentiable at \(0\) (requires subgradients). \\
\textbf{Outliers} & More sensitive to outliers (due to squaring). & More
robust. \\
\textbf{Prior Belief} & Weights are normally distributed. & Weights are
Laplace distributed. \\
\end{longtable}
}

\subsubsection{Key Takeaways}\label{key-takeaways-9}

\begin{itemize}
\tightlist
\item
  \textbf{Necessity:} Regularization is crucial in Logistic Regression
  to prevent weights from exploding (\(\infty\)) on separable data.
\item
  \textbf{L2 (Ridge):} Use for general performance improvement and
  handling collinear features. Smooths the weights.
\item
  \textbf{L1 (Lasso):} Use when you suspect only a few features are
  relevant (sparse signal) or need an interpretable model.
\item
  \textbf{\(\lambda\) Selection:} The regularization strength
  \(\lambda\) is a hyperparameter usually tuned via Cross-Validation.

  \begin{itemize}
  \tightlist
  \item
    Small \(\lambda\): High variance (Overfitting).
  \item
    Large \(\lambda\): High bias (Underfitting).
  \end{itemize}
\end{itemize}

\newpage

\section{Boosting and Bagging}\label{boosting-and-bagging}

\textbf{Boosting} and \textbf{Bagging} are \emph{ensemble} learning
techniques. Instead of trying to learn one single, highly complex
predictor, these methods combine multiple simple predictors (hypotheses)
to form a stronger aggregate predictor.

\subsection{1. What is Boosting?}\label{what-is-boosting}

Boosting is a general method for converting rough rules of thumb into
highly accurate prediction rules. * \textbf{Concept:} It trains a
sequence of ``weak'' classifiers iteratively. Each new classifier
focuses on the examples that \textbf{the previous classifiers got
wrong}. * \textbf{Aggregation:} The final prediction is a
\textbf{weighted majority vote} of all the weak classifiers. *
\textbf{Goal:} Primarily to reduce \textbf{Bias} (and also Variance). It
can turn a model that slightly underfits into one that fits the data
well.

\subsection{2. Weak Learnability}\label{weak-learnability}

To formalize boosting, we must define what a ``weak'' learner is.

\textbf{Definition (Strong Learnability - PAC):} A concept class is
strongly learnable if there exists an algorithm that can achieve
arbitrarily low error \(\epsilon\) with high probability, for any
distribution \(\mathcal{D}\).

\textbf{Definition (Weak Learnability):} A concept class is weakly
learnable if there exists an algorithm that can produce a hypothesis
\(h\) whose error \textbf{is slightly better than random guessing}.
\[ L_{\mathcal{D}}(h) \leq \frac{1}{2} - \gamma \] where \(\gamma > 0\)
is a small constant (called the ``edge'').

\textbf{The Fundamental Question:} * \emph{Does Weak Learnability imply
Strong Learnability?} * \textbf{Answer:} Yes. (Proven by Schapire in
1990). * \textbf{Implication:} If we can consistently find a rule just
slightly better than a coin toss, we can combine enough of them to
create an arbitrarily accurate classifier. Boosting is the constructive
proof of this theorem.

\subsection{3. AdaBoost (Adaptive
Boosting)}\label{adaboost-adaptive-boosting}

The most popular boosting algorithm is \textbf{AdaBoost} (Freund \&
Schapire, 1995). It solves the problem of how to force the weak learner
to focus on difficult examples.

\subsubsection{The Algorithm}\label{the-algorithm}

\textbf{Input:} Training set \(S = \{(x_1, y_1), \dots, (x_m, y_m)\}\);
Weak learning algorithm \(\mathcal{A}\).

\textbf{Initialization:} Set observation weights
\(D_1(i) = \frac{1}{m}\) for all \(i=1, \dots, m\).

\textbf{For \(t = 1, \dots, T\):} 1. \textbf{Train:} Call weak learner
\(\mathcal{A}\) using distribution (weights) \(D_t\). Get hypothesis
\(h_t: \mathcal{X} \to \{-1, +1\}\). 2. \textbf{Calculate Error:}
Compute the weighted error of \(h_t\):
\[ \epsilon_t = \sum_{i=1}^{m} D_t(i) \mathbb{1}_{\{h_t(x_i) \neq y_i\}} \]
3. \textbf{Compute Importance (\(\alpha_t\)):} Calculate the weight of
this classifier in the final vote:
\[ \alpha_t = \frac{1}{2} \ln \left( \frac{1 - \epsilon_t}{\epsilon_t} \right) \]
\emph{(Note: If \(\epsilon_t \approx 0\), \(\alpha_t\) is large. If
\(\epsilon_t \approx 0.5\), \(\alpha_t \approx 0\).)} 4. \textbf{Update
Weights:} Update distribution for the next iteration:
\[ D_{t+1}(i) = \frac{D_t(i) \exp(-\alpha_t y_i h_t(x_i))}{Z_t} \] where
\(Z_t\) is a normalization factor to ensure \(\sum D_{t+1}(i) = 1\). *
If \(y_i = h_t(x_i)\) (Correct): Weight decreases (multiplied by
\(e^{-\alpha} < 1\)). * If \(y_i \neq h_t(x_i)\) (Incorrect): Weight
increases (multiplied by \(e^{\alpha} > 1\)).

\textbf{Output:} The final hypothesis is a weighted sign vote:
\[ H(x) = \text{sign} \left( \sum_{t=1}^{T} \alpha_t h_t(x) \right) \]

\subsubsection{Intuition}\label{intuition}

\begin{enumerate}
\def\labelenumi{\arabic{enumi}.}
\tightlist
\item
  \textbf{Reweighting:} Samples that are hard to classify get higher
  weights. The next weak learner is forced to focus on them to minimize
  the weighted error.
\item
  \textbf{Combination:} Accurate classifiers get more say (\(\alpha_t\))
  in the final decision.
\item
  \textbf{Exponential Loss:} It can be shown that AdaBoost minimizes the
  exponential loss function \(L(H) = \sum e^{-y_i H(x_i)}\) via
  coordinate descent.
\end{enumerate}

\pandocbounded{\includegraphics[keepaspectratio]{/home/nicola/Documents/Local Obs/Obsidian Locale/images/e9e3c9070bd77dc45e50109e83b8340c8cf50467.jpg}}

\subsection{4. Bagging (Bootstrap
Aggregating)}\label{bagging-bootstrap-aggregating}

\textbf{Bagging} is a different ensemble approach designed to reduce
\textbf{Variance}. It is particularly useful for unstable learners like
Decision Trees (which are prone to overfitting).

\subsubsection{The Algorithm}\label{the-algorithm-1}

\begin{enumerate}
\def\labelenumi{\arabic{enumi}.}
\tightlist
\item
  \textbf{Bootstrap Sampling:} Given a training set \(S\) of size \(m\),
  create \(T\) new datasets \(S_1, \dots, S_T\). Each \(S_t\) is created
  by sampling \(m\) items from \(S\) uniformly at random \textbf{with
  replacement}.

  \begin{itemize}
  \tightlist
  \item
    \emph{Note:} Each \(S_t\) contains about \(63.2\%\) of the original
    unique examples; the rest are duplicates.
  \end{itemize}
\item
  \textbf{Parallel Training:} Train a base learner (e.g., a deep
  Decision Tree) independently on each \(S_t\) to obtain
  \(h_1, \dots, h_T\).
\item
  \textbf{Aggregation:}

  \begin{itemize}
  \tightlist
  \item
    \textbf{Classification:} Majority Vote:
    \(H(x) = \text{mode}\{h_1(x), \dots, h_T(x)\}\).
  \item
    \textbf{Regression:} Average: \(H(x) = \frac{1}{T} \sum h_t(x)\).
  \end{itemize}
\end{enumerate}

\subsubsection{Why it Works (Variance
Reduction)}\label{why-it-works-variance-reduction}

If we have independent random variables \(Z_1, \dots, Z_T\) each with
variance \(\sigma^2\), the variance of their mean \(\bar{Z}\) is
\(\frac{\sigma^2}{T}\). While the models \(h_t\) are not perfectly
independent (they share the original dataset \(S\)), Bagging simulates
this effect. By averaging many ``overfitted'' (high variance) models,
the noise cancels out, yielding a stable predictor.

\textbf{Famous Example:} \textbf{Random Forest} is essentially Bagging
applied to Decision Trees, with an extra trick (random feature
selection) to further decorrelate the trees.

\subsection{Comparison: Boosting
vs.~Bagging}\label{comparison-boosting-vs.-bagging}

{\def\LTcaptype{none} % do not increment counter
\begin{longtable}[]{@{}
  >{\raggedright\arraybackslash}p{(\linewidth - 4\tabcolsep) * \real{0.3333}}
  >{\raggedright\arraybackslash}p{(\linewidth - 4\tabcolsep) * \real{0.3333}}
  >{\raggedright\arraybackslash}p{(\linewidth - 4\tabcolsep) * \real{0.3333}}@{}}
\toprule\noalign{}
\begin{minipage}[b]{\linewidth}\raggedright
Feature
\end{minipage} & \begin{minipage}[b]{\linewidth}\raggedright
\textbf{Boosting} (e.g., AdaBoost, XGBoost)
\end{minipage} & \begin{minipage}[b]{\linewidth}\raggedright
\textbf{Bagging} (e.g., Random Forest)
\end{minipage} \\
\midrule\noalign{}
\endhead
\bottomrule\noalign{}
\endlastfoot
\textbf{Goal} & Reduces \textbf{Bias} (and Variance) & Reduces
\textbf{Variance} \\
\textbf{Base Learner} & Must be \textbf{Weak} (high bias, low variance,
e.g., shallow tree/stump) & Must be \textbf{Strong/Complex} (low bias,
high variance, e.g., deep tree) \\
\textbf{Training} & \textbf{Sequential} (Model \(t\) depends on Model
\(t-1\)) & \textbf{Parallel} (Independent training) \\
\textbf{Weights} & Examples are re-weighted based on difficulty. & All
examples effectively have equal weight (via resampling). \\
\textbf{Outliers} & Sensitive (tries hard to fix errors). & Robust
(averages them out). \\
\textbf{Overfitting} & Can overfit if \(T\) is too large (though often
robust). & Harder to overfit; adding more trees usually helps. \\
\end{longtable}
}

\subsection{Key Takeaways}\label{key-takeaways-10}

\begin{enumerate}
\def\labelenumi{\arabic{enumi}.}
\item
  \textbf{Weak Learnability:} A concept is learnable if we can perform
  slightly better than random guessing.
\item
  \textbf{AdaBoost:} An iterative algorithm that focuses on ``hard''
  examples, turning weak learners into a strong one.
\item
  \textbf{Bagging:} Uses bootstrap sampling and averaging to stabilize
  high-variance models.
\item
  \textbf{Trade-off:} Use Boosting to improve underfitting models; use
  Bagging to improve overfitting models.

  \newpage
\end{enumerate}

\section{Model Selection and Performance
Evaluation}\label{model-selection-and-performance-evaluation}

Model selection is the process of choosing the best machine learning
model from a set of candidate models, which may differ in terms of
algorithm (e.g., SVM vs.~Logistic Regression) or hyperparameters (e.g.,
polynomial degree \(d\), regularization strength \(\lambda\)).

The primary goal is to estimate the \textbf{Generalization Error}---the
error on unseen data---rather than optimizing the Training Error, to
avoid overfitting.

\subsection{1. Train-Test Split (Hold-out
Method)}\label{train-test-split-hold-out-method}

The simplest method for evaluating a model's performance is to split the
available dataset \(\mathcal{D}\) into two mutually exclusive subsets: a
\textbf{Training Set} and a \textbf{Test Set}.

\subsubsection{1.1 Methodology}\label{methodology}

\begin{enumerate}
\def\labelenumi{\arabic{enumi}.}
\item
  \textbf{Shuffle} the dataset \(\mathcal{D}\) randomly to remove any
  ordering bias.
\item
  \textbf{Split} \(\mathcal{D}\) into \(\mathcal{D}_{train}\) and
  \(\mathcal{D}_{test}\) based on a split ratio (e.g., 70/30 or 80/20).

  \begin{itemize}
  \item
    \(|\mathcal{D}_{train}| \approx (1 - \alpha) |\mathcal{D}|\)
  \item
    \(|\mathcal{D}_{test}| \approx \alpha |\mathcal{D}|\)
  \end{itemize}
\item
  \textbf{Train} the model \(h_\theta\) on \(\mathcal{D}_{train}\) to
  find optimal parameters \(\theta^*\).
\item
  \textbf{Evaluate} \(h_{\theta^*}\) on \(\mathcal{D}_{test}\) to
  estimate generalization performance.
\end{enumerate}

\subsubsection{1.2 Mathematical
Formulation}\label{mathematical-formulation}

Let \(L(y, \hat{y})\) be the loss function (e.g., Squared Error for
regression, 0-1 Loss for classification).

\begin{itemize}
\item
  \textbf{Training Error:}

  \[E_{train}(\theta) = \frac{1}{|\mathcal{D}_{train}|} \sum_{(x, y) \in \mathcal{D}_{train}} L(y, h_\theta(x))\]
\item
  \textbf{Test Error (Hold-out Estimate):}

  \[E_{test}(\theta) = \frac{1}{|\mathcal{D}_{test}|} \sum_{(x, y) \in \mathcal{D}_{test}} L(y, h_\theta(x))\]
\end{itemize}

The hold-out estimate \(E_{test}\) provides an unbiased estimate of the
true generalization error, provided the test set is never used during
training or hyperparameter tuning.

\subsubsection{1.3 Validation Set
Extension}\label{validation-set-extension}

If we need to perform \textbf{Model Selection} (hyperparameter tuning),
a simple Train-Test split is insufficient because using the Test set to
select the best model biases the error estimate. We introduce a third
set:

\begin{itemize}
\item
  \textbf{Training Set:} Used to learn parameters \(\theta\).
\item
  \textbf{Validation Set:} Used to tune hyperparameters and select the
  best model structure.
\item
  \textbf{Test Set:} Used \emph{only once} at the very end to estimate
  the generalization error of the final chosen model.
\end{itemize}

\subsubsection{1.4 Pros and Cons}\label{pros-and-cons}

{\def\LTcaptype{none} % do not increment counter
\begin{longtable}[]{@{}
  >{\raggedright\arraybackslash}p{(\linewidth - 2\tabcolsep) * \real{0.5000}}
  >{\raggedright\arraybackslash}p{(\linewidth - 2\tabcolsep) * \real{0.5000}}@{}}
\toprule\noalign{}
\begin{minipage}[b]{\linewidth}\raggedright
\textbf{Pros}
\end{minipage} & \begin{minipage}[b]{\linewidth}\raggedright
\textbf{Cons}
\end{minipage} \\
\midrule\noalign{}
\endhead
\bottomrule\noalign{}
\endlastfoot
Computationally efficient (trains only once). & High variance:
Performance estimate depends heavily on exactly which points end up in
the test set. \\
Simple to implement and interpret. & Wasteful: A significant portion of
data is withheld from training. \\
Suitable for very large datasets where \(N\) is massive. & Problematic
for small datasets (insufficient training data). \\
\end{longtable}
}

\subsection{2. k-fold Cross-Validation
(CV)}\label{k-fold-cross-validation-cv}

To address the variance and data-wastage issues of the Hold-out method,
particularly for smaller datasets, we use \(k\)-fold Cross-Validation.

\subsubsection{2.1 Methodology}\label{methodology-1}

\begin{enumerate}
\def\labelenumi{\arabic{enumi}.}
\item
  Randomly partition the dataset \(\mathcal{D}\) into \(k\) equal-sized
  disjoint subsets (folds)
  \(\mathcal{F}_1, \mathcal{F}_2, \dots, \mathcal{F}_k\).
\item
  For each iteration \(i = 1\) to \(k\):

  \begin{itemize}
  \item
    \textbf{Train} the model on all folds except \(\mathcal{F}_i\)
    (i.e.,
    \(\mathcal{D}_{train}^{(i)} = \mathcal{D} \setminus \mathcal{F}_i\)).
  \item
    \textbf{Validate} the model on the held-out fold \(\mathcal{F}_i\)
    to obtain error \(E_i\).
  \end{itemize}
\item
  \textbf{Average} the \(k\) error estimates to get the final
  Cross-Validation Error.
\end{enumerate}

\pandocbounded{\includegraphics[keepaspectratio]{/home/nicola/Documents/Local Obs/Obsidian Locale/images/3681c0d63a8099f6b428903d4cfeed2b8fe818db.webp}}

\subsubsection{2.2 Mathematical
Formulation}\label{mathematical-formulation-1}

The Cross-Validation Score (CV Error) is defined as:

\[CV_{(k)} = \frac{1}{k} \sum_{i=1}^{k} E_i\]

where \(E_i\) is the validation error on fold \(\mathcal{F}_i\).

To estimate the standard error of the mean (uncertainty of the
estimate):

\[SE_{CV} = \frac{\sigma}{\sqrt{k}} \quad \text{where} \quad \sigma = \sqrt{\frac{1}{k-1}\sum_{i=1}^k (E_i - CV_{(k)})^2}\]

\subsubsection{\texorpdfstring{2.3 Choosing
\(k\)}{2.3 Choosing k}}\label{choosing-k}

\begin{itemize}
\item
  \textbf{\(k=5\) or \(k=10\):} Standard choices. Offers a good
  trade-off between bias and variance.
\item
  \textbf{\(k=N\) (Leave-One-Out CV - LOOCV):}

  \begin{itemize}
  \item
    \textbf{Description:} Each data point is used as a single test set
    once.
  \item
    \textbf{Pros:} Unbiased estimate of the true expected error (uses
    almost all data for training). Deterministic (no random splitting
    variance).
  \item
    \textbf{Cons:} Computationally expensive (\(N\) training runs). High
    variance in the error estimator itself (training sets are very
    highly correlated).
  \end{itemize}
\end{itemize}

\subsubsection{2.4 Properties}\label{properties-1}

\begin{itemize}
\item
  \textbf{Robustness:} Uses all data points for both training and
  validation (in different iterations).
\item
  \textbf{Model Selection:} Commonly used to select hyperparameters. We
  compute \(CV_{(k)}\) for each hyperparameter combination and choose
  the one with the lowest average error.
\item
  \textbf{Final Model:} After selecting the best hyperparameters via CV,
  the final model is typically retrained on the \textbf{entire} dataset
  \(\mathcal{D}\) before deployment.
\end{itemize}

\subsection{3. Learning Curves}\label{learning-curves}

Learning curves are diagnostic plots that visualize the model's
performance (Training Error and Validation Error) as a function of the
\textbf{Training Set Size} (\(m\)). They help diagnose \textbf{Bias
(Underfitting)} vs.~\textbf{Variance (Overfitting)}.

\subsubsection{3.1 Structure of the Plot}\label{structure-of-the-plot}

\begin{itemize}
\item
  \textbf{X-axis:} Number of training examples (\(m\)).
\item
  \textbf{Y-axis:} Error (Loss).
\item
  \textbf{Curves:} Two curves are plotted:

  \begin{enumerate}
  \def\labelenumi{\arabic{enumi}.}
  \item
    \(J_{train}(\theta)\): Average error on the training set.
  \item
    \(J_{val}(\theta)\) (or \(J_{cv}(\theta)\)): Average error on the
    validation/CV set.
  \end{enumerate}
\end{itemize}

\subsubsection{3.2 Diagnosing Bias
vs.~Variance}\label{diagnosing-bias-vs.-variance}

\paragraph{Case A: High Bias
(Underfitting)}\label{case-a-high-bias-underfitting}

The model is too simple to capture the underlying structure of the data
(e.g., fitting a line to quadratic data).

\begin{itemize}
\item
  \textbf{Behavior as \(m \to \infty\):}

  \begin{itemize}
  \item
    \textbf{Training Error:} Increases rapidly and plateaus at a high
    value. (The model cannot even fit the training data well).
  \item
    \textbf{Validation Error:} Decreases initially but plateaus at a
    similarly high value.
  \item
    \textbf{Gap:} The gap between Training and Validation error is
    \textbf{small} or non-existent.
  \end{itemize}
\item
  \textbf{Solution:} Increase model complexity (add features, increase
  polynomial degree), decrease regularization. Adding more data
  \textbf{will not} help.
\end{itemize}

\pandocbounded{\includegraphics[keepaspectratio]{/home/nicola/Documents/Local Obs/Obsidian Locale/images/d2bbb46f2611a7c0c546f4bae900263c9737624f.png}}

\paragraph{Case B: High Variance
(Overfitting)}\label{case-b-high-variance-overfitting}

The model is too complex and fits the noise in the training data (e.g.,
high-degree polynomial).

\begin{itemize}
\item
  \textbf{Behavior as \(m \to \infty\):}

  \begin{itemize}
  \item
    \textbf{Training Error:} Remains very low (model memorizes training
    data).
  \item
    \textbf{Validation Error:} Remains high (poor generalization).
  \item
    \textbf{Gap:} A \textbf{large gap} exists between Training and
    Validation error.
  \end{itemize}
\item
  \textbf{Solution:} Get more training data (helps validation error
  converge to training error), reduce model complexity (feature
  selection), increase regularization.
\end{itemize}

\subsubsection{3.3 Ideal Learning Curve}\label{ideal-learning-curve}

\begin{itemize}
\item
  Both errors converge to a low value.
\item
  The gap between \(J_{train}\) and \(J_{val}\) is small.
\item
  This indicates the model has low bias and low variance.
\end{itemize}

\subsubsection{Key Takeaways}\label{key-takeaways-11}

\begin{itemize}
\item
  \textbf{Train-Test Split:} Fast and simple, but sensitive to how the
  data is split. Good for large \(N\).
\item
  \textbf{k-fold CV:} The gold standard for model selection on
  small-to-medium datasets. Reduces variance of the error estimate.
\item
  \textbf{Hold-out for Selection vs.~Assessment:} Use a Validation set
  (or CV) for \emph{selecting} the model, and a separate independent
  Test set for the final \emph{assessment}.
\item
  \textbf{Learning Curves:} Essential for debugging.

  \begin{itemize}
  \item
    High Error + Small Gap \(\rightarrow\) \textbf{High Bias}.
  \item
    Low Train Error + High Val Error + Large Gap \(\rightarrow\)
    \textbf{High Variance}.
  \end{itemize}

  \newpage
\end{itemize}

\section{Gradient Descent
Optimization}\label{gradient-descent-optimization}

Gradient Descent (GD) is an iterative optimization algorithm used to
minimize a function. In Machine Learning, it is the primary method used
to minimize the \textbf{Cost Function} \(J(\theta)\) (also called Loss
Function) by iteratively adjusting the model parameters \(\theta\).

The goal is to find the optimal parameters \(\theta^*\) such that:

\[\theta^* = \underset{\theta}{\text{argmin}} \ J(\theta)\]

\subsection{1. The Gradient (Vector of Partial
Derivatives)}\label{the-gradient-vector-of-partial-derivatives}

To understand how to minimize the cost function, we first define the
\textbf{Gradient}.

\subsubsection{1.1 Mathematical
Definition}\label{mathematical-definition}

For a model with \(n\) parameters
\(\theta = [\theta_0, \theta_1, \dots, \theta_n]^T\), the cost function
\(J(\theta)\) is a function of \(\mathbb{R}^{n+1} \to \mathbb{R}\).

The gradient, denoted by \(\nabla_\theta J(\theta)\), is a
\textbf{vector} containing the partial derivatives of the cost function
with respect to \emph{each} parameter.

\[\nabla_\theta J(\theta) = \begin{bmatrix} \frac{\partial J}{\partial \theta_0} \\ \frac{\partial J}{\partial \theta_1} \\ \vdots \\ \frac{\partial J}{\partial \theta_n} \end{bmatrix}\]

\subsubsection{1.2 Intuitive Meaning}\label{intuitive-meaning}

\begin{itemize}
\item
  \textbf{Direction:} The gradient vector points in the direction of the
  \textbf{steepest ascent} (greatest increase) of the function \(J\).
\item
  \textbf{Magnitude:} The magnitude \(||\nabla J||\) indicates the
  steepness of the slope.
\item
  \textbf{Goal:} Since we want to \emph{minimize} \(J\), we must move in
  the direction \textbf{opposite} to the gradient (negative gradient).
\end{itemize}

\pandocbounded{\includegraphics[keepaspectratio]{/home/nicola/Documents/Local Obs/Obsidian Locale/images/a1205455276678d62b52ba9bfae60302b90314f4.webp}}

\subsection{2. The Update Rule}\label{the-update-rule}

The core of the algorithm is the update rule, which adjusts the
parameters \(\theta\) step-by-step.

\subsubsection{2.1 The Algorithm}\label{the-algorithm-2}

\begin{enumerate}
\def\labelenumi{\arabic{enumi}.}
\item
  \textbf{Initialization:} Start with random values for \(\theta\)
  (e.g., \(\theta = \vec{0}\) or small random numbers).
\item
  \textbf{Iteration:} Repeat the update step until convergence (when
  \(J(\theta)\) stops decreasing significantly or
  \(\nabla J \approx 0\)).
\end{enumerate}

\subsubsection{2.2 The General Update Formula
(Vectorized)}\label{the-general-update-formula-vectorized}

\[\theta := \theta - \alpha \nabla_\theta J(\theta)\]

Where:

\begin{itemize}
\item
  \(\theta\): The parameter vector.
\item
  \(\alpha\): The \textbf{Learning Rate} (a positive scalar
  hyperparameter).
\item
  \(\nabla_\theta J(\theta)\): The gradient vector.
\item
  \(:=\): Assignment operator (update the old value with the new one).
\end{itemize}

\subsubsection{2.3 Component-wise Update}\label{component-wise-update}

If we expand the vectors, the update for the specific \(j\)-th parameter
\(\theta_j\) is:

\[\theta_j := \theta_j - \alpha \frac{\partial}{\partial \theta_j} J(\theta)\]

This update is performed \textbf{simultaneously} for all
\(j = 0, \dots, n\) in a single iteration.

\subsection{\texorpdfstring{3. The Learning Rate
(\(\alpha\))}{3. The Learning Rate (\textbackslash alpha)}}\label{the-learning-rate-alpha}

The learning rate \(\alpha\) controls the \textbf{step size} of the
iteration. It is a critical hyperparameter.

\begin{itemize}
\item
  \textbf{If \(\alpha\) is too small:}

  \begin{itemize}
  \item
    The algorithm makes tiny steps.
  \item
    Convergence is slow and computationally expensive.
  \end{itemize}
\item
  \textbf{If \(\alpha\) is too large:}

  \begin{itemize}
  \item
    The algorithm may overshoot the minimum.
  \item
    It may fail to converge or even diverge (move away from the
    minimum).
  \end{itemize}
\end{itemize}

\subsection{4. Concrete Example: Linear
Regression}\label{concrete-example-linear-regression}

To make this concrete, let's apply GD to Linear Regression with the Mean
Squared Error (MSE) cost function.

\textbf{Hypothesis:} \(h_\theta(x) = \theta_0 + \theta_1 x\)

\textbf{Cost Function:}
\(J(\theta_0, \theta_1) = \frac{1}{2m} \sum_{i=1}^{m} (h_\theta(x^{(i)}) - y^{(i)})^2\)

\subsubsection{4.1 Deriving the Partial
Derivatives}\label{deriving-the-partial-derivatives}

We calculate the derivative of \(J\) with respect to a specific weight
\(\theta_j\):

\[\frac{\partial}{\partial \theta_j} J(\theta) = \frac{1}{m} \sum_{i=1}^{m} (h_\theta(x^{(i)}) - y^{(i)}) \cdot x_j^{(i)}\]

\subsubsection{4.2 Specific Update Rules}\label{specific-update-rules}

Substituting the derivative back into the general update rule:

\begin{itemize}
\item
  \textbf{For Bias (\(\theta_0\)):} (Note that \(x_0 = 1\))

  \[\theta_0 := \theta_0 - \alpha \frac{1}{m} \sum_{i=1}^{m} (h_\theta(x^{(i)}) - y^{(i)})\]
\item
  \textbf{For Weight (\(\theta_1\)):}

  \[\theta_1 := \theta_1 - \alpha \frac{1}{m} \sum_{i=1}^{m} (h_\theta(x^{(i)}) - y^{(i)}) \cdot x^{(i)}\]
\end{itemize}

\subsection{5. Variations of Gradient
Descent}\label{variations-of-gradient-descent}

{\def\LTcaptype{none} % do not increment counter
\begin{longtable}[]{@{}
  >{\raggedright\arraybackslash}p{(\linewidth - 8\tabcolsep) * \real{0.2000}}
  >{\raggedright\arraybackslash}p{(\linewidth - 8\tabcolsep) * \real{0.2000}}
  >{\raggedright\arraybackslash}p{(\linewidth - 8\tabcolsep) * \real{0.2000}}
  >{\raggedright\arraybackslash}p{(\linewidth - 8\tabcolsep) * \real{0.2000}}
  >{\raggedright\arraybackslash}p{(\linewidth - 8\tabcolsep) * \real{0.2000}}@{}}
\toprule\noalign{}
\begin{minipage}[b]{\linewidth}\raggedright
\textbf{Variant}
\end{minipage} & \begin{minipage}[b]{\linewidth}\raggedright
\textbf{Data used per step}
\end{minipage} & \begin{minipage}[b]{\linewidth}\raggedright
\textbf{Gradient Accuracy}
\end{minipage} & \begin{minipage}[b]{\linewidth}\raggedright
\textbf{Speed per step}
\end{minipage} & \begin{minipage}[b]{\linewidth}\raggedright
\textbf{Convergence}
\end{minipage} \\
\midrule\noalign{}
\endhead
\bottomrule\noalign{}
\endlastfoot
\textbf{Batch GD} & Entire dataset (\(m\)) & Exact & Slow & Stable,
direct path to minima \\
\textbf{Stochastic GD (SGD)} & Single random example & Noisy
approximation & Very Fast & Fluctuates around minima \\
\textbf{Mini-Batch GD} & Small batch (e.g., 32, 64) & Good approximation
& Fast & Stable (Best of both worlds) \\
\end{longtable}
}

\subsubsection{Key Takeaways}\label{key-takeaways-12}

\begin{itemize}
\item
  \textbf{Gradient Vector (\(\nabla J\)):} A vector of partial
  derivatives pointing in the direction of steepest increase.
\item
  \textbf{Descent:} We subtract the gradient to move down the ``error
  mountain.''
\item
  \textbf{Update Rule:}
  \(\theta_{new} = \theta_{old} - \alpha \cdot \text{gradient}\).
\item
  \textbf{Convergence:} Occurs when the gradient becomes zero (flat
  slope at the minimum).
\item
  \textbf{Convexity:} If \(J(\theta)\) is a convex function (like MSE),
  Gradient Descent is guaranteed to find the Global Minimum (provided
  \(\alpha\) is not too large). For non-convex functions (like Neural
  Networks), it may get stuck in Local Minima.

  \newpage
\end{itemize}

\section{Subgradients and Non-Differentiable
Optimization}\label{subgradients-and-non-differentiable-optimization}

In standard Gradient Descent, we rely on the gradient \(\nabla f(x)\) to
point us toward the minimum. However, many important Machine Learning
loss functions are \textbf{non-differentiable} (non-smooth) at certain
points.

To optimize these functions, we generalize the concept of the derivative
to the \textbf{Subgradient}.

\subsection{1. The Problem:
Non-Differentiability}\label{the-problem-non-differentiability}

A function is non-differentiable if it has ``kinks'' or sharp corners
where a unique tangent plane cannot be defined.

\subsubsection{Common Examples in ML}\label{common-examples-in-ml}

\begin{enumerate}
\def\labelenumi{\arabic{enumi}.}
\item
  \textbf{L1 Regularization (Lasso):} \(f(w) = ||w||_1 = \sum |w_i|\).
  The absolute value function \(|x|\) has a sharp corner at \(x=0\).
\item
  \textbf{Hinge Loss (SVM):} \(L(y, f(x)) = \max(0, 1 - y \cdot f(x))\).
  This function has a ``kink'' at the point where
  \(1 - y \cdot f(x) = 0\).
\item
  \textbf{ReLU Activation:} \(f(x) = \max(0, x)\). Non-differentiable at
  \(x=0\).
\end{enumerate}

At these points, the gradient \(\nabla f(x)\) is undefined. The
\textbf{Subgradient Method} allows us to run optimization algorithms
like SGD on these functions.

\subsection{2. The Subgradient}\label{the-subgradient}

\subsubsection{2.1 Formal Definition}\label{formal-definition}

Let \(f: \mathbb{R}^n \to \mathbb{R}\) be a \textbf{convex} function. A
vector \(g \in \mathbb{R}^n\) is a \textbf{subgradient} of \(f\) at
point \(x\) if for all \(z \in \mathbb{R}^n\):

\[f(z) \geq f(x) + g^T (z - x)\]

\pandocbounded{\includegraphics[keepaspectratio]{/home/nicola/Documents/Local Obs/Obsidian Locale/images/761a520db6509eb4fbe522fc4f9b94880c18b9ca.png}}
\#\#\# 2.2 Intuitive Explanation

\begin{itemize}
\item
  \textbf{Geometric Interpretation:} In standard calculus, the
  derivative defines a \emph{unique} tangent line that touches the curve
  at \(x\) and stays below the convex function.
\item
  \textbf{Subgradient Interpretation:} At a ``kink,'' there isn't just
  one tangent line; there are infinitely many lines that pass through
  \(x\) and stay below the function \(f(z)\).
\item
  \textbf{The Subdifferential (\(\partial f(x)\)):} This is the
  \textbf{set} of all possible subgradients at point \(x\). If the
  function is differentiable at \(x\), the set contains only one vector:
  \textbf{the gradient itself.}
\end{itemize}

\subsubsection{2.3 Example: Absolute
Value}\label{example-absolute-value}

Consider \(f(x) = |x|\).

\begin{itemize}
\item
  If \(x > 0\), the slope is unique: \(g = 1\).
\item
  If \(x < 0\), the slope is unique: \(g = -1\).
\item
  If \(x = 0\), any line with a slope between \(-1\) and \(1\) stays
  below the graph.

  \[\partial f(0) = [-1, 1]\]

  Therefore, any value \(g \in [-1, 1]\) is a valid subgradient at
  \(x=0\).
\end{itemize}

\subsection{3. The Subgradient Method}\label{the-subgradient-method}

The Subgradient Method is a generalization of Gradient Descent.

\subsubsection{3.1 The Update Rule}\label{the-update-rule-1}

To minimize a non-differentiable convex function \(J(\theta)\), we
iterate:

\[\theta^{(k+1)} = \theta^{(k)} - \alpha_k \cdot g^{(k)}\]

Where:

\begin{itemize}
\item
  \(g^{(k)}\) is \textbf{any} subgradient chosen from the set
  \(\partial J(\theta^{(k)})\).
\item
  \(\alpha_k\) is the step size.
\end{itemize}

\subsubsection{3.2 Key Differences from Gradient
Descent}\label{key-differences-from-gradient-descent}

\begin{enumerate}
\def\labelenumi{\arabic{enumi}.}
\item
  \textbf{Direction:} A subgradient \(g\) is not necessarily the
  direction of \emph{steepest} descent. It merely makes an angle less
  than 90 degrees with the direction toward the minimum.
\item
  \textbf{No Monotonicity:} In standard GD, proper step sizes guarantee
  \(J(\theta^{(k+1)}) < J(\theta^{(k)})\). In the subgradient method,
  the objective function value \textbf{may increase} temporarily, even
  with small steps. We usually keep track of the best value seen so far
  (\(J_{best}\)).
\item
  \textbf{Convergence:} The method converges, but generally slower than
  GD on smooth functions (\(O(1/\sqrt{k})\) vs \(O(1/k)\)).
\end{enumerate}

\subsection{4. Importance for Stochastic Gradient Descent
(SGD)}\label{importance-for-stochastic-gradient-descent-sgd}

The concept of the subgradient is crucial for SGD because \textbf{Online
Learning} and large-scale algorithms often rely on loss functions that
are sparse (L1) or margin-based (SVM).

\subsubsection{4.1 Case Study: SGD for Support Vector Machines
(Pegasos)}\label{case-study-sgd-for-support-vector-machines-pegasos}

The classic algorithm for training linear SVMs on massive datasets is
essentially \textbf{SGD with Subgradients}.

\textbf{Objective (Primal SVM with Regularization):}

\[J(w) = \frac{\lambda}{2} ||w||^2 + \frac{1}{m} \sum_{i=1}^m \max(0, 1 - y_i w^T x_i)\]

When we pick a random example \((x_i, y_i)\) for SGD, we need the
gradient of the Hinge Loss term \(L_i(w) = \max(0, 1 - y_i w^T x_i)\).

The subgradient of \(L_i(w)\) with respect to \(w\) is:

\begin{itemize}
\item
  \textbf{If \(y_i w^T x_i < 1\) (Misclassified or within margin):} The
  function is linear with slope \(-y_i x_i\).

  \[g = -y_i x_i\]
\item
  \textbf{If \(y_i w^T x_i > 1\) (Correctly classified):} The function
  is flat (0).

  \[g = 0\]
\item
  \textbf{If \(y_i w^T x_i = 1\) (Exactly on the margin):} The
  subgradient is undefined in the classical sense, but we can choose any
  vector in the convex hull. Usually, we simplify and choose
  \(g = -y_i x_i\) or \(0\).
\end{itemize}

\textbf{SGD Update Rule for SVM:}

\[w := w - \alpha \left( \lambda w - \mathbb{I}[y_i w^T x_i < 1] \cdot y_i x_i \right)\]

Without the theory of subgradients, we mathematically could not justify
applying gradient-based updates to the ``kink'' in the Hinge Loss.

\subsubsection{4.2 L1 Regularization and
Sparsity}\label{l1-regularization-and-sparsity}

When using L1 regularization (\(J(\theta) + \lambda ||\theta||_1\)), the
gradient at \(\theta_j = 0\) is technically undefined.

Using subgradients allows the algorithm to theoretically ``land'' on
exactly zero. (Note: In practice, proximal gradient methods are often
preferred for L1 to enforce exact sparsity, but subgradients provide the
theoretical basis for why the descent works).

\subsubsection{Key Takeaways}\label{key-takeaways-13}

\begin{itemize}
\item
  \textbf{Subgradient:} A generalized derivative for non-differentiable
  convex functions. It is a vector defining a plane that supports the
  function from below.
\item
  \textbf{Subdifferential:} The set of all valid subgradients at a point
  (e.g., at a sharp corner).
\item
  \textbf{Necessity:} We cannot train SVMs (Hinge Loss) or
  L1-regularized models using standard calculus because the gradient is
  undefined at the ``kinks.''
\item
  \textbf{Convergence:} Subgradient methods are robust but slower than
  differentiable methods. They require a diminishing step size (e.g.,
  \(\alpha_k \propto 1/k\)) to guarantee convergence to the global
  minimum.

  \newpage
\end{itemize}

\section{Variations of Gradient
Descent}\label{variations-of-gradient-descent-1}

The standard Gradient Descent algorithm updates parameters based on the
gradient of the cost function \(J(\theta)\). However, calculating the
gradient requires summing errors over data points. The difference
between Batch, Stochastic, and Mini-Batch lies in \textbf{how many data
points} are used to compute the gradient in a single update step.

\subsection{1. Batch Gradient Descent
(BGD)}\label{batch-gradient-descent-bgd}

This is the ``standard'' theoretical version of the algorithm.

\subsubsection{1.1 Mechanism}\label{mechanism}

It computes the gradient using the \textbf{entire training dataset}
(\(m\) examples) for every single step.

\[\theta_{new} = \theta_{old} - \alpha \cdot \frac{1}{m} \sum_{i=1}^{m} \nabla_\theta L(h_\theta(x^{(i)}), y^{(i)})\]

\subsubsection{1.2 Characteristics}\label{characteristics}

\begin{itemize}
\item
  \textbf{Path to Minimum:} Moves directly downhill. The path is smooth
  and deterministic.
\item
  \textbf{Convergence:} Guaranteed to converge to the global minimum
  (for convex surfaces) or a local minimum (non-convex).
\item
  \textbf{Computational Cost:} extremely high. If \(m = 10,000,000\), we
  must process 10 million examples just to take \emph{one} tiny step.
\item
  \textbf{Memory:} Requires loading the entire dataset into memory
  (RAM/VRAM), which is often impossible for Big Data.
\end{itemize}

\subsection{2. Stochastic Gradient Descent
(SGD)}\label{stochastic-gradient-descent-sgd}

At the other extreme, SGD updates the parameters using \textbf{a single
training example} \((x^{(i)}, y^{(i)})\) chosen uniformly at random.

\subsubsection{2.1 Mechanism}\label{mechanism-1}

\[\theta_{new} = \theta_{old} - \alpha \cdot \nabla_\theta L(h_\theta(x^{(i)}), y^{(i)})\]

\subsubsection{2.2 Characteristics}\label{characteristics-1}

\begin{itemize}
\item
  \textbf{Path to Minimum:} extremely noisy. The gradient of a single
  example is a poor approximation of the true gradient. The path
  ``zig-zags'' violently.
\item
  \textbf{Convergence:} It does not converge to the exact minimum but
  wanders around it. (Convergence requires a diminishing learning rate
  schedule, e.g., \(\alpha_k \propto 1/k\)).
\item
  \textbf{Computational Cost:} Very low per step.
\item
  \textbf{Efficiency:} Bad for hardware. Processors (CPUs/GPUs) are
  optimized for vector operations (SIMD), not sequential scalar
  operations.
\end{itemize}

\subsection{3. Mini-Batch Gradient
Descent}\label{mini-batch-gradient-descent}

This is the ``sweet spot'' between Batch and Stochastic GD. It processes
data in small groups called \textbf{batches}.

\subsubsection{3.1 Mechanism}\label{mechanism-2}

We split the training set into small batches of size \(b\) (typically a
power of 2: 32, 64, 128, 256).

\[\theta_{new} = \theta_{old} - \alpha \cdot \frac{1}{b} \sum_{i=k}^{k+b} \nabla_\theta L(h_\theta(x^{(i)}), y^{(i)})\]

\subsubsection{3.2 Terminology}\label{terminology}

\begin{itemize}
\item
  \textbf{Batch Size (\(b\)):} The number of examples per update
  (hyperparameter).
\item
  \textbf{Iteration:} One update step (processing one batch).
\item
  \textbf{Epoch:} One complete pass through the entire dataset.

  \begin{itemize}
  \tightlist
  \item
    \emph{Example:} If \(m=1000\) and \(b=100\), one epoch consists of
    10 iterations.
  \end{itemize}
\end{itemize}

\subsubsection{3.3 Characteristics}\label{characteristics-2}

\begin{itemize}
\item
  \textbf{Path:} Less noisy than SGD, but faster than Batch GD.
\item
  \textbf{Efficiency:} Highly efficient. GPUs are massive parallel
  matrix multipliers. Computing the gradient for 64 examples
  simultaneously takes roughly the same time as for 1 example, but
  provides a much more accurate gradient.
\end{itemize}

\pandocbounded{\includegraphics[keepaspectratio]{/home/nicola/Documents/Local Obs/Obsidian Locale/images/c8ed71f04f4b53b1ef0b1572cd679665917588a0.png}}

\subsection{4. Comparative Analysis}\label{comparative-analysis}

{\def\LTcaptype{none} % do not increment counter
\begin{longtable}[]{@{}
  >{\raggedright\arraybackslash}p{(\linewidth - 6\tabcolsep) * \real{0.2500}}
  >{\raggedright\arraybackslash}p{(\linewidth - 6\tabcolsep) * \real{0.2500}}
  >{\raggedright\arraybackslash}p{(\linewidth - 6\tabcolsep) * \real{0.2500}}
  >{\raggedright\arraybackslash}p{(\linewidth - 6\tabcolsep) * \real{0.2500}}@{}}
\toprule\noalign{}
\begin{minipage}[b]{\linewidth}\raggedright
\textbf{Feature}
\end{minipage} & \begin{minipage}[b]{\linewidth}\raggedright
\textbf{Batch GD}
\end{minipage} & \begin{minipage}[b]{\linewidth}\raggedright
\textbf{Stochastic GD (SGD)}
\end{minipage} & \begin{minipage}[b]{\linewidth}\raggedright
\textbf{Mini-Batch GD}
\end{minipage} \\
\midrule\noalign{}
\endhead
\bottomrule\noalign{}
\endlastfoot
\textbf{Dataset Usage} & Entire dataset (\(m\)) & 1 example & Batch size
(\(b \approx 32-512\)) \\
\textbf{Update Frequency} & Once per epoch & \(m\) times per epoch &
\(m/b\) times per epoch \\
\textbf{Speed per Step} & Very Slow & Very Fast & Fast \\
\textbf{Convergence} & Stable, Smooth & Noisy, Oscillating & Stable
enough \\
\textbf{Memory Usage} & High & Low & Moderate (User defined) \\
\textbf{Vectorization} & Good & Poor & \textbf{Excellent} \\
\end{longtable}
}

\subsection{5. What is Used in Production
Today?}\label{what-is-used-in-production-today}

In modern Machine Learning and Deep Learning production environments,
\textbf{Mini-Batch Gradient Descent is the standard.}

However, it is rarely used in its ``vanilla'' (plain) form. It is almost
always combined with \textbf{optimization algorithms} that accelerate
convergence.

\subsubsection{5.1 Why Mini-Batch
Dominates}\label{why-mini-batch-dominates}

\begin{enumerate}
\def\labelenumi{\arabic{enumi}.}
\item
  \textbf{Hardware Architecture:} Modern hardware (NVIDIA GPUs, Google
  TPUs) relies on \textbf{SIMD (Single Instruction, Multiple Data)}.
  They are designed to multiply large matrices. Mini-batches allow us to
  utilize the full parallel processing power of the GPU. Pure SGD leaves
  the GPU idle; Batch GD exceeds GPU memory (VRAM).
\item
  \textbf{Regularization Effect:} The slight noise introduced by
  mini-batches (compared to full batch) acts as a form of
  regularization, helping the model escape sharp local minima or saddle
  points, which Batch GD might get stuck in.
\end{enumerate}

\subsubsection{5.2 The ``Real'' Production
Algorithms}\label{the-real-production-algorithms}

When a Data Scientist selects ``SGD'' in a framework like PyTorch or
TensorFlow, they are actually running \textbf{Mini-Batch SGD}.

In practice, we use advanced variations of Mini-Batch GD:

\paragraph{A. Mini-Batch with
Momentum}\label{a.-mini-batch-with-momentum}

Instead of using only the current gradient, we use a moving average of
past gradients to ``gain speed'' in relevant directions and dampen
oscillations.

\pandocbounded{\includegraphics[keepaspectratio]{/home/nicola/Documents/Local Obs/Obsidian Locale/images/286a41830e443632af7fc39c6ce1d225b2df44c7.png}}

\[v_t = \gamma v_{t-1} + \eta \nabla J(\theta)\]

\[\theta = \theta - v_t\]

\paragraph{B. Adaptive Learning Rate Methods
(Adam)}\label{b.-adaptive-learning-rate-methods-adam}

\textbf{Adam (Adaptive Moment Estimation)} is currently the default
optimizer for most Deep Learning applications (LLMs, Computer Vision).

\begin{itemize}
\item
  It computes individual adaptive learning rates for different
  parameters.
\item
  It combines the benefits of Momentum and RMSProp.
\item
  It requires less tuning of the learning rate hyperparameter
  \(\alpha\).
\end{itemize}

\subsubsection{Key Takeaways}\label{key-takeaways-14}

\begin{itemize}
\item
  \textbf{Batch GD} is theoretically pure but computationally
  impractical for Big Data.
\item
  \textbf{SGD} (single sample) is fast but noisy and computationally
  inefficient on GPUs.
\item
  \textbf{Mini-Batch GD} is the industry standard because it balances
  gradient accuracy with hardware efficiency (Vectorization).
\item
  \textbf{Production Reality:} We use Mini-Batch GD wrapped in
  sophisticated optimizers like \textbf{Adam} or \textbf{SGD with
  Momentum}.
\item
  \textbf{Note on Frameworks:} In code (e.g., \texttt{torch.optim.SGD}),
  the term ``SGD'' almost always refers to Mini-Batch GD, where the
  batch size is determined by the Data Loader, not the optimizer itself.

  \newpage
\end{itemize}

\section{Support Vector Machines (SVM) - Part 1:
Foundations}\label{support-vector-machines-svm---part-1-foundations}

Support Vector Machines (SVM) are a powerful class of supervised
learning algorithms for classification and regression. Developed by
Vladimir Vapnik and colleagues (1963, 1990s) at AT\&T Bell Labs, they
are rooted in \textbf{Statistical Learning Theory} (VC Theory).

Unlike heuristics that just ``find a line that separates data,'' SVMs
are based on a rigorous mathematical framework to maximize
generalization performance.

\subsection{1. Vapnik's Idea: The Geometric
Intuition}\label{vapniks-idea-the-geometric-intuition}

In a binary classification problem, we have a dataset
\(\mathcal{D} = \{(x_i, y_i)\}_{i=1}^N\), where \(x_i \in \mathbb{R}^d\)
and \(y_i \in \{-1, +1\}\).

We assume the data is \textbf{linearly separable} (for now).

There are infinite hyperplanes that can separate the two classes with
zero training error. Which one is the best?

\subsubsection{1.1 The Margin}\label{the-margin}

Vapnik's key insight is the \textbf{Margin}: the smallest distance
between the decision boundary (hyperplane) and any of the training
samples.

\begin{itemize}
\item
  \textbf{Intuition:} A larger margin implies ``safer'' classification.
  If we place the boundary as far as possible from the data points, the
  model is more robust to noise in future unseen data.
\item
  \textbf{Goal:} Find the hyperplane that separates the classes and
  \textbf{maximizes the margin}.
\end{itemize}

\subsubsection{1.2 Support Vectors}\label{support-vectors}

Not all data points matter. The position of the optimal hyperplane is
determined \emph{only} by the data points closest to it.

\begin{itemize}
\item
  These critical points are called \textbf{Support Vectors}.
\item
  If we remove all other data points and keep only the support vectors,
  the optimal hyperplane remains exactly the same.
\item
  This property makes SVMs memory efficient (the solution is sparse in
  the dual space).
\end{itemize}

\subsection{2. Mathematical
Formulation}\label{mathematical-formulation-2}

\subsubsection{2.1 The Hyperplane}\label{the-hyperplane}

A hyperplane is defined by a normal vector \(w\) (weights) and a bias
term \(b\):

\[h(x) = w^T x + b = 0\]

The decision rule for a new point \(x\) is:

\[f(x) = \text{sign}(w^T x + b)\]

\subsubsection{2.2 Canonical Hyperplanes}\label{canonical-hyperplanes}

The distance from a point \(x_i\) to the hyperplane \((w, b)\) is given
by:

\[\text{dist}(x_i, w, b) = \frac{|w^T x_i + b|}{||w||}\]

Since we want to classify correctly, \(y_i(w^T x_i + b) \geq 0\) for all
\(i\).

To fix the scale of \(w\) and \(b\) (which are scale-invariant), we
introduce the \textbf{Canonical constraints}:

We require that for the \emph{closest} points (support vectors), the
functional margin is exactly 1.

\[y_i (w^T x_i + b) \geq 1 \quad \forall i=1, \dots, N\]

\begin{itemize}
\item
  For Support Vectors: \(y_i (w^T x_i + b) = 1\).
\item
  For all other points: \(y_i (w^T x_i + b) > 1\).
\end{itemize}

\subsubsection{2.3 The Geometric Margin}\label{the-geometric-margin}

With the canonical constraint, the distance of the closest point (margin
\(\gamma\)) becomes:

\[\gamma = \frac{1}{||w||}\]

Therefore, maximizing the margin \(\gamma\) is equivalent to
\textbf{minimizing the norm \(||w||\)}.

\subsection{3. The Primal Optimization Problem
(ERM)}\label{the-primal-optimization-problem-erm}

We now have a constrained optimization problem. This is a convex
quadratic programming (QP) problem.

\subsubsection{3.1 Hard Margin SVM (Linearly
Separable)}\label{hard-margin-svm-linearly-separable}

To maximize the margin \(\frac{1}{||w||}\), we minimize \(||w||^2\):

\[\begin{aligned} & \underset{w, b}{\text{minimize}} & & \frac{1}{2} ||w||^2 \\ & \text{subject to} & & y_i (w^T x_i + b) \geq 1, \quad i = 1, \dots, N \end{aligned}\]

\begin{itemize}
\item
  \textbf{Objective:} Minimize the complexity (norm) of the model
  (structural risk).
\item
  \textbf{Constraints:} Ensure every single training point is correctly
  classified (empirical risk = 0).
\end{itemize}

\subsubsection{3.2 Soft Margin SVM
(Non-Separable)}\label{soft-margin-svm-non-separable}

Real-world data is rarely perfectly separable (noise, outliers). We
introduce \textbf{Slack Variables} \(\xi_i \geq 0\) to allow some
misclassification.

\[y_i (w^T x_i + b) \geq 1 - \xi_i\]

The optimization problem becomes:

\[\begin{aligned} & \underset{w, b, \xi}{\text{minimize}} & & \frac{1}{2} ||w||^2 + C \sum_{i=1}^N \xi_i \\ & \text{subject to} & & y_i (w^T x_i + b) \geq 1 - \xi_i, \quad \forall i \\ & & & \xi_i \geq 0, \quad \forall i \end{aligned}\]

\begin{itemize}
\item
  \textbf{\(\frac{1}{2}||w||^2\):} Regularization term (Max Margin).
\item
  \textbf{\(\sum \xi_i\):} Hinge Loss (Empirical Error).
\item
  \textbf{\(C\):} Hyperparameter controlling the trade-off.

  \begin{itemize}
  \item
    High \(C\): Strict penalty for errors (Low Bias, High Variance).
  \item
    Low \(C\): Wider margin, tolerates errors (High Bias, Low Variance).
  \end{itemize}
\end{itemize}

This formulation is exactly \textbf{Regularized Empirical Risk
Minimization (ERM)} with Hinge Loss.

\subsection{4. The Lagrangian
Formulation}\label{the-lagrangian-formulation}

To solve this constrained optimization efficiently (and to allow for
Kernels later), we switch to the \textbf{Dual Problem} using Lagrange
Multipliers.

\subsubsection{4.1 The Primal Lagrangian}\label{the-primal-lagrangian}

We introduce Lagrange multipliers \(\alpha_i \geq 0\) for each
inequality constraint \(g_i(w, b) = 1 - y_i(w^T x_i + b) \leq 0\).

The Lagrangian \(\mathcal{L}(w, b, \alpha)\) is:

\[\mathcal{L}(w, b, \alpha) = \underbrace{\frac{1}{2} ||w||^2}_{\text{Objective}} - \sum_{i=1}^N \alpha_i \underbrace{[y_i (w^T x_i + b) - 1]}_{\text{Constraint}}\]

We want to minimize \(\mathcal{L}\) with respect to the primal variables
(\(w, b\)) and maximize it with respect to the dual variables
(\(\alpha_i\)).

\[\min_{w, b} \max_{\alpha \geq 0} \mathcal{L}(w, b, \alpha)\]

\subsubsection{4.2 Deriving the Dual}\label{deriving-the-dual}

To find the minimum w.r.t \(w\) and \(b\), we set the gradients to zero:

\begin{enumerate}
\def\labelenumi{\arabic{enumi}.}
\item
  \textbf{Gradient w.r.t \(w\):}

  \[\nabla_w \mathcal{L} = w - \sum_{i=1}^N \alpha_i y_i x_i = 0 \implies w = \sum_{i=1}^N \alpha_i y_i x_i\]

  \emph{Note: The weight vector \(w\) is a linear combination of the
  training data!}
\item
  \textbf{Gradient w.r.t \(b\):}

  \[\nabla_b \mathcal{L} = - \sum_{i=1}^N \alpha_i y_i = 0 \implies \sum_{i=1}^N \alpha_i y_i = 0\]
\end{enumerate}

\subsubsection{4.3 The Dual Optimization
Problem}\label{the-dual-optimization-problem}

Substitute these results back into the Primal Lagrangian. The primal
variables \(w\) and \(b\) vanish, leaving only \(\alpha\).

\[\begin{aligned} & \underset{\alpha}{\text{maximize}} & & \sum_{i=1}^N \alpha_i - \frac{1}{2} \sum_{i=1}^N \sum_{j=1}^N \alpha_i \alpha_j y_i y_j (x_i^T x_j) \\ & \text{subject to} & & \alpha_i \geq 0, \quad \forall i \\ & & & \sum_{i=1}^N \alpha_i y_i = 0 \end{aligned}\]

\emph{(For Soft Margin, the constraint becomes
\(0 \leq \alpha_i \leq C\))}.

\subsubsection{Key Takeaways}\label{key-takeaways-15}

\begin{itemize}
\item
  \textbf{Vapnik's Idea:} Generalization improves by maximizing the
  \textbf{Margin} (distance to nearest points).
\item
  \textbf{Support Vectors:} Only the difficult points (on the margin)
  determine the model. \(w\) is a weighted sum of these support vectors.
\item
  \textbf{Primal Problem:} Minimize \(||w||^2\) subject to constraints.
  Intuitive but depends on dimension \(d\).
\item
  \textbf{ERM connection:} SVM minimizes Regularization (\(||w||^2\)) +
  Loss (Hinge Loss).
\item
  \textbf{Dual Problem:} Maximizing \(\alpha\) depends only on dot
  products \(x_i^T x_j\). This is the key to the \textbf{Kernel
  Trick}.\\

  \newpage
\end{itemize}

\section{Non-Linear SVMs and The Kernel
Trick}\label{non-linear-svms-and-the-kernel-trick}

Standard SVMs (Linear SVMs) are powerful, but they have a fatal flaw:
they assume the data is \textbf{linearly separable}. In the real world,
data is rarely linear (e.g., the XOR problem, concentric circles).

To handle non-linear data, we use the \textbf{Kernel Trick}, a method
that allows linear algorithms to learn non-linear decision boundaries
efficiently.

\subsection{\texorpdfstring{1. The Intuition: Feature Mapping
(\(\phi\))}{1. The Intuition: Feature Mapping (\textbackslash phi)}}\label{the-intuition-feature-mapping-phi}

\subsubsection{1.1 The Problem}\label{the-problem}

Imagine a dataset in 1D space (\(\mathbb{R}^1\)) where:

\begin{itemize}
\item
  Class \(-1\) points are near the origin (e.g., \(x \in [-2, 2]\)).
\item
  Class \(+1\) points are far away (e.g., \(x < -5\) or \(x > 5\)).
\end{itemize}

A single linear threshold (a point in 1D) cannot separate these classes.
We need two cuts, which a linear classifier cannot provide.

\subsubsection{1.2 The Solution: Higher
Dimensions}\label{the-solution-higher-dimensions}

If we map the input data into a \textbf{higher-dimensional feature
space} \(\mathcal{H}\), the data might become linearly separable there.

Let \(\phi: \mathcal{X} \to \mathcal{H}\) be a \textbf{feature map}.

For our 1D example, let \(\phi(x) = [x, x^2]^T\).

Now the data exists in 2D. The parabolic relationship (\(y = x^2\))
makes the classes separable by a straight line in this new space.

\subsubsection{1.3 The Computational
Bottleneck}\label{the-computational-bottleneck}

If we map data to a very high dimension (e.g., \(10^6\) dimensions),
calculating the feature vector \(\phi(x)\) becomes computationally
prohibitive. This is known as the \textbf{Curse of Dimensionality}.

\subsection{2. The Kernel Trick}\label{the-kernel-trick}

The ``Trick'' is a mathematical shortcut that allows us to operate in
the high-dimensional space \(\mathcal{H}\) \textbf{without ever
explicitly calculating the coordinates \(\phi(x)\)}.

\subsubsection{2.1 Revisiting the Dual
Form}\label{revisiting-the-dual-form}

Recall the Dual Optimization Problem for SVM derived in the previous
section:

\[\underset{\alpha}{\text{maximize}} \sum_{i=1}^N \alpha_i - \frac{1}{2} \sum_{i=1}^N \sum_{j=1}^N \alpha_i \alpha_j y_i y_j (x_i^T x_j)\]

\textbf{Crucial Observation:} The optimization depends \emph{only} on
the \textbf{dot product} (scalar product) between data points
\(\langle x_i, x_j \rangle\), not on the individual vectors themselves.

\subsubsection{2.2 Substituting the Map}\label{substituting-the-map}

If we transform data using \(\phi\), the dot product becomes
\(\langle \phi(x_i), \phi(x_j) \rangle\).

We define a \textbf{Kernel Function} \(K\) such that:

\[K(x_i, x_j) = \langle \phi(x_i), \phi(x_j) \rangle\]

The Kernel Trick is the ability to compute this inner product directly
using the original inputs \(x_i, x_j\), without calculating \(\phi\).

\subsection{2.3 Why the Dot Product is
Everything}\label{why-the-dot-product-is-everything}

The key point of the Kernel Trick is not just that we avoid computing
the feature map \(\phi(x)\).\\
The deeper reason is that \textbf{SVMs only need dot products} between
points.

In the dual formulation, every training and prediction step depends on
quantities of the form

\[
\langle x_i, x_j \rangle
\]

and, after the feature mapping, on

\[
\langle \phi(x_i), \phi(x_j) \rangle.
\]

So the algorithm never needs the coordinates of \(\phi(x)\) by
themselves.\\
It only needs to know how pairs of transformed points interact through
their \textbf{inner product}.

This is crucial because the dot product already contains geometric
information:

\begin{itemize}
\tightlist
\item
  it measures how much two vectors are aligned;
\item
  it determines their angle;
\item
  it gives their norm, since
\end{itemize}

\[
\|\phi(x)\|^2 = \langle \phi(x), \phi(x) \rangle;
\]

\begin{itemize}
\tightlist
\item
  it also gives distances, because
\end{itemize}

\[
\|\phi(x)-\phi(z)\|^2
=
\|\phi(x)\|^2+\|\phi(z)\|^2-2\langle \phi(x),\phi(z)\rangle.
\]

Therefore, knowing all the inner products between points means knowing
the relevant geometry of the data in the feature space \(\mathcal H\),
even without explicitly knowing the coordinates of the mapped vectors.

This is exactly what the kernel function provides:

\[
K(x,z)=\langle \phi(x),\phi(z)\rangle.
\]

So the Kernel Trick works because it gives us direct access to the
\textbf{geometry of the high-dimensional space through dot products},
without ever constructing that space explicitly.

In other words:

\begin{itemize}
\tightlist
\item
  we do \textbf{not} explicitly see the coordinates of the data in
  \(\mathcal H\);
\item
  but we do know how points relate to each other geometrically in
  \(\mathcal H\);
\item
  and for SVM, that is enough.
\end{itemize}

That is why the trick is possible:\\
\textbf{the learning algorithm depends on scalar products, and the
kernel computes those scalar products directly in the original space.}

\newpage

\subsection{2.4 Geometric
Interpretation}\label{geometric-interpretation}

A linear separator in the feature space \(\mathcal H\) has the form

\[
f(x)=\langle w,\phi(x)\rangle + b.
\]

Using the representer form of the solution,

\[
w=\sum_{i=1}^N \alpha_i y_i \phi(x_i),
\]

we get

\[
f(x)=\sum_{i=1}^N \alpha_i y_i \langle \phi(x_i),\phi(x)\rangle + b
=
\sum_{i=1}^N \alpha_i y_i K(x_i,x)+b.
\]

Again, only dot products appear.

So the classifier decides the label of a point \(x\) by comparing it,
through inner products, with the support vectors in feature space.\\
The kernel therefore lets us compute a non-linear boundary in the
original input space \(\mathcal X\), while still using a linear
separator in \(\mathcal H\).

\subsection{3. Concrete Example: Polynomial
Kernel}\label{concrete-example-polynomial-kernel}

Let \(x, z \in \mathbb{R}^2\) (2D input vectors).

Consider the kernel function:

\[K(x, z) = (x^T z)^2\]

Let's expand this to see what feature map \(\phi\) it corresponds to.

Let \(x = [x_1, x_2]^T\) and \(z = [z_1, z_2]^T\).

\[\begin{aligned} K(x, z) &= (x_1 z_1 + x_2 z_2)^2 \\ &= (x_1 z_1)^2 + (x_2 z_2)^2 + 2 x_1 z_1 x_2 z_2 \\ &= x_1^2 z_1^2 + x_2^2 z_2^2 + (\sqrt{2} x_1 x_2)(\sqrt{2} z_1 z_2) \end{aligned}\]

This is exactly the dot product of the mapping
\(\phi(x) = [x_1^2, x_2^2, \sqrt{2} x_1 x_2]^T\).

\begin{itemize}
\item
  \textbf{Explicit Mapping Cost:} calculating \(\phi(x)\) takes time
  proportional to the number of dimensions in \(\mathcal{H}\).
\item
  \textbf{Kernel Calculation:} calculating \((x^T z)^2\) takes time
  proportional to the original dimension \(d\).
\end{itemize}

\textbf{Result:} We get the benefit of working in a quadratic space
without the computational cost.

\subsection{4. Common Kernels}\label{common-kernels}

\subsubsection{4.1 Linear Kernel}\label{linear-kernel}

\[K(x, z) = x^T z\]

\begin{itemize}
\item
  Standard linear SVM.
\item
  Used when data is linearly separable or \(d\) is very large (e.g.,
  text classification).
\end{itemize}

\subsubsection{4.2 Polynomial Kernel}\label{polynomial-kernel}

\[K(x, z) = (x^T z + c)^d\]

\begin{itemize}
\item
  Maps data to a space of all monomials of degree up to \(d\).
\item
  Parameters: \(d\) (degree), \(c\) (bias).
\end{itemize}

\subsubsection{4.3 Radial Basis Function (RBF) / Gaussian
Kernel}\label{radial-basis-function-rbf-gaussian-kernel}

This is the most popular kernel.

\[K(x, z) = \exp\left(-\frac{||x - z||^2}{2\sigma^2}\right) = \exp(-\gamma ||x - z||^2)\]

\begin{itemize}
\item
  \textbf{Infinite Dimensions:} The Taylor expansion of the exponential
  function shows that the RBF kernel corresponds to a feature map
  \(\phi\) into an \textbf{infinite-dimensional} space.
\item
  \textbf{Intuition:} It measures \textbf{similarity}. If
  \(x \approx z\), \(K \approx 1\). If they are far apart,
  \(K \approx 0\).
\item
  \textbf{Hyperparameter:} \(\gamma\) (gamma).

  \begin{itemize}
  \item
    High \(\gamma\): Complex boundary, risks overfitting (only points
    very close influence prediction).
  \item
    Low \(\gamma\): Smoother boundary.
  \end{itemize}
\end{itemize}

\subsubsection{4.4 Sigmoid Kernel (Neural
Network)}\label{sigmoid-kernel-neural-network}

\[K(x, z) = \tanh(\alpha x^T z + c)\]

\begin{itemize}
\item
  Makes the SVM behave like a Neural Network (Multi-Layer Perceptron).
\item
  \emph{Note: Not always a valid kernel (not always Positive
  Semi-Definite).}
\end{itemize}

\subsection{5. Mercer's Theorem (Theory)}\label{mercers-theorem-theory}

Not any function \(K(x, z)\) can be a kernel. It must correspond to a
valid dot product in some space \(\mathcal{H}\).

\textbf{Mercer's Condition:}

A symmetric function \(K(x, z)\) is a valid kernel if and only if the
\textbf{Gram Matrix} (the \(N \times N\) matrix where
\(G_{ij} = K(x_i, x_j)\)) is \textbf{Positive Semi-Definite (PSD)} for
any dataset.

\[z^T G z \geq 0, \quad \forall z \in \mathbb{R}^N\]

This ensures the optimization problem remains \textbf{convex},
guaranteeing a unique global minimum.

\subsubsection{Key Takeaways}\label{key-takeaways-16}

\begin{itemize}
\item
  \textbf{Motivation:} Linear boundaries in high-dimensional space
  (\(\mathcal{H}\)) correspond to non-linear boundaries in input space
  (\(\mathcal{X}\)).
\item
  \textbf{The Trick:} We compute the high-dimensional dot product
  \(\langle \phi(x), \phi(z) \rangle\) using a function \(K(x, z)\) in
  the low-dimensional space.
\item
  \textbf{Efficiency:} We avoid the computational cost of the explicit
  map \(\phi(x)\).
\item
  \textbf{Flexibility:} By changing \(K\), we change the geometry of the
  decision boundary without changing the optimization algorithm (Dual
  Solver).
\item
  \textbf{RBF Kernel:} Maps to infinite dimensions; acts as a similarity
  measure.

  \newpage
\end{itemize}

\section{Decision Trees}\label{decision-trees}

Decision Trees are a non-parametric supervised learning method used for
classification and regression. They model decisions as a tree-like
structure of binary (or multi-way) choices, making them highly
\textbf{interpretable}.

\subsection{1. Example of Usage}\label{example-of-usage}

Consider a credit card approval system. We want to predict if a customer
is a \textbf{``High Risk''} or \textbf{``Low Risk''} borrower based on
their attributes.

\begin{itemize}
\item
  \textbf{Root Node:} \emph{Income \textgreater{} \$50k?}

  \begin{itemize}
  \item
    \textbf{Yes (Right Branch):} \emph{Has Criminal Record?}

    \begin{itemize}
    \item
      \textbf{Yes:} \(\rightarrow\) \textbf{High Risk} (Leaf Node)
    \item
      \textbf{No:} \(\rightarrow\) \textbf{Low Risk} (Leaf Node)
    \end{itemize}
  \item
    \textbf{No (Left Branch):} \emph{Employment Status = Employed?}

    \begin{itemize}
    \item
      \textbf{Yes:} \(\rightarrow\) \textbf{Low Risk}
    \item
      \textbf{No:} \(\rightarrow\) \textbf{High Risk}
    \end{itemize}
  \end{itemize}
\end{itemize}

The model partitions the input space into rectangular regions, assigning
a constant prediction (class or mean value) to each region.

\subsection{2. The Training Algorithm (Recursive
Partitioning)}\label{the-training-algorithm-recursive-partitioning}

The standard algorithm for building a decision tree (e.g., ID3, C4.5,
CART) is a \textbf{greedy top-down} approach.

\subsubsection{2.1 The Steps}\label{the-steps}

\begin{enumerate}
\def\labelenumi{\arabic{enumi}.}
\item
  Start with the entire dataset \(\mathcal{D}\) at the root node.
\item
  \textbf{Feature Selection:} Evaluate all possible features (and split
  points) to find the one that best separates the classes (maximizes
  purity).
\item
  \textbf{Split:} Partition the data into subsets based on the selected
  feature.
\item
  \textbf{Recurse:} Repeat the process recursively for each child node.
\item
  \textbf{Stop:} Terminate when:

  \begin{itemize}
  \item
    All samples in a node belong to the same class (Pure node).
  \item
    Maximum depth is reached.
  \item
    The number of samples is below a threshold.
  \end{itemize}
\end{enumerate}

\subsection{3. Splitting Criteria: Measuring
``Gain''}\label{splitting-criteria-measuring-gain}

How do we decide which feature is the ``best'' to split on? We need a
mathematical measure of \textbf{impurity}.

\subsubsection{3.1 Entropy (Shannon
Entropy)}\label{entropy-shannon-entropy}

Entropy measures the uncertainty or disorder in a dataset \(S\). For a
binary classification (\(p_+\) and \(p_-\)):

\[H(S) = -p_+ \log_2(p_+) - p_- \log_2(p_-)\]

\begin{itemize}
\item
  \textbf{Low Entropy (\(H \approx 0\)):} The node is pure (all samples
  are one class).
\item
  \textbf{High Entropy (\(H \approx 1\)):} The node is a 50/50 mix.
\end{itemize}

\subsubsection{3.2 Information Gain (IG)}\label{information-gain-ig}

Information Gain measures the expected reduction in entropy after
splitting set \(S\) on attribute \(A\).

\[IG(S, A) = H(S) - \sum_{v \in \text{Values}(A)} \frac{|S_v|}{|S|} H(S_v)\]

\begin{itemize}
\item
  \(H(S)\): Entropy of the parent node.
\item
  The second term is the \textbf{Weighted Average Entropy} of the
  children nodes.
\item
  \textbf{Goal:} Choose attribute \(A\) that maximizes \(IG(S, A)\).
\end{itemize}

\subsubsection{3.3 Gini Impurity (Used in
CART)}\label{gini-impurity-used-in-cart}

An alternative to Entropy (computationally faster as it avoids
logarithms).

\[Gini(S) = 1 - \sum_{k=1}^K p_k^2\]

It measures the probability of incorrectly classifying a randomly chosen
element if it were labeled according to the distribution of labels in
the subset.

\subsection{4. VC Dimension of Decision
Trees}\label{vc-dimension-of-decision-trees}

The Vapnik-Chervonenkis (VC) dimension measures the capacity of the
hypothesis class.

\subsubsection{4.1 Unbounded Trees}\label{unbounded-trees}

\begin{itemize}
\item
  If the tree depth is unlimited, a decision tree can memorize
  \emph{any} dataset with distinct examples (by creating a specific leaf
  for every single data point).
\item
  \textbf{VC Dimension = \(\infty\)}.
\item
  This implies that without regularization (limits), Decision Trees are
  prone to massive \textbf{overfitting}.
\end{itemize}

\subsubsection{4.2 Bounded Trees}\label{bounded-trees}

To guarantee generalization, we must restrict the complexity.

For a binary decision tree with \(N\) nodes (or depth \(d\)) defined on
inputs in \(\mathbb{R}^n\):

\begin{itemize}
\item
  The VC dimension is roughly proportional to the number of nodes (or
  \(2^d\)).
\item
  \(VC \approx O(\text{number of leaves} \cdot n)\).
\end{itemize}

This theoretical result justifies why we must limit the size of the
tree.

\subsection{5. Pruning (Regularization)}\label{pruning-regularization}

Because standard trees grow until training error is zero (High
Variance), we use \textbf{Pruning} to reduce complexity and improve
generalization.

\subsubsection{5.1 Pre-Pruning (Early
Stopping)}\label{pre-pruning-early-stopping}

Stop the tree growth \emph{during} the training process before it
perfectly fits the data.

\begin{itemize}
\item
  \textbf{Max Depth:} Limit the depth to \(d\) (e.g., 5).
\item
  \textbf{Min Samples Split:} Do not split a node if it has fewer than
  \(k\) samples.
\item
  \textbf{Min Impurity Decrease:} Do not split if the Information Gain
  is below a threshold \(\epsilon\).
\item
  \textbf{Pros:} Fast.
\item
  \textbf{Cons:} Can be ``short-sighted'' (greedy). It might stop too
  early, missing a very good split that would have occurred deeper in
  the tree.
\end{itemize}

\subsubsection{5.2 Post-Pruning (Cost-Complexity
Pruning)}\label{post-pruning-cost-complexity-pruning}

Grow the full tree \(T_0\) (overfit), then cut back branches.

We minimize a loss function that penalizes size:

\[R_\alpha(T) = R(T) + \alpha \cdot |T|\]

\begin{itemize}
\item
  \(R(T)\): The training error (misclassification rate).
\item
  \(|T|\): Number of leaf nodes (complexity).
\item
  \(\alpha \geq 0\): Complexity parameter (like \(\lambda\) in Lasso).
\end{itemize}

\textbf{Algorithm:}

\begin{enumerate}
\def\labelenumi{\arabic{enumi}.}
\item
  Compute the effective \(\alpha\) for every subtree.
\item
  Remove the subtree that gives the smallest increase in error per leaf
  removed.
\item
  Use Cross-Validation to select the best \(\alpha\).
\end{enumerate}

\subsubsection{Key Takeaways}\label{key-takeaways-17}

\begin{itemize}
\item
  \textbf{Intuition:} Recursive binary splitting of the input space.
\item
  \textbf{Information Gain:} The metric used to greedily select the best
  split (max reduction in Entropy).
\item
  \textbf{VC Dimension:} Infinite for unrestricted trees; requires
  constraints to generalize.
\item
  \textbf{Overfitting:} Trees naturally overfit. A fully grown tree has
  low Bias but very High Variance.
\item
  \textbf{Pruning:} The critical technique to balance the tree.
  Post-pruning is generally more robust than pre-pruning.

  \newpage
\end{itemize}

\section{Random Forests and Ensemble
Learning}\label{random-forests-and-ensemble-learning}

Random Forest (Breiman, 2001) is an \textbf{Ensemble Learning} method
that constructs a multitude of Decision Trees at training time and
outputs the class that is the mode of the classes (classification) or
mean prediction (regression) of the individual trees.

It is designed specifically to address the main weakness of Decision
Trees: \textbf{High Variance} (overfitting).

\subsection{1. The Core Mechanism: Bagging (Bootstrap
Aggregation)}\label{the-core-mechanism-bagging-bootstrap-aggregation}

Random Forest is a modification of \textbf{Bagging}. To understand
Random Forest, we must first understand Bagging.

\subsubsection{1.1 Bootstrapping}\label{bootstrapping}

Given a training set \(D\) of size \(N\), Bagging generates \(B\) new
training sets \(D_1, D_2, \dots, D_B\), each of size \(N\), by sampling
from \(D\) \textbf{uniformly and with replacement}.

\begin{itemize}
\item
  Some observations may be repeated in a set \(D_b\).
\item
  Some observations may be omitted (on average, \(D_b\) contains only
  \textasciitilde63.2\% of unique original examples).
\end{itemize}

\subsubsection{1.2 Aggregation}\label{aggregation}

We train a separate model (tree) \(T_b\) on each bootstrap sample
\(D_b\).

The final prediction \(\hat{f}_{rf}(x)\) is the average:

\[\hat{f}_{rf}(x) = \frac{1}{B} \sum_{b=1}^{B} T_b(x)\]

\subsection{2. The ``Random'' in Random Forest
(Decorrelation)}\label{the-random-in-random-forest-decorrelation}

If we simply apply Bagging to Decision Trees, we get a ``Bagged Tree''
model. Random Forest adds one crucial modification to
\textbf{decorrelate} the trees.

\subsubsection{2.1 The Correlation
Problem}\label{the-correlation-problem}

In standard trees, if one feature is very strong (e.g., ``Income'' in
credit scoring), \textbf{every} tree in the bag will likely choose
``Income'' as the top split.

Consequently, all trees will look structurally similar. Their
predictions will be highly \textbf{correlated}.

Averaging highly correlated variables does not significantly reduce
variance.

\subsubsection{2.2 Feature Randomness}\label{feature-randomness}

To fix this, Random Forest forces trees to be different.

\begin{itemize}
\item
  \textbf{Standard Tree:} At each split, search through \textbf{all}
  \(p\) features to find the best split.
\item
  \textbf{Random Forest:} At each split, select a \textbf{random subset}
  of \(m\) features (\(m < p\)). The best split is chosen \emph{only}
  from this subset.
\end{itemize}

\subsubsection{\texorpdfstring{2.3 Hyperparameter
\(m\)}{2.3 Hyperparameter m}}\label{hyperparameter-m}

\begin{itemize}
\item
  \textbf{Classification:} \(m \approx \sqrt{p}\)
\item
  \textbf{Regression:} \(m \approx p/3\)
\end{itemize}

By restricting the features, we force the trees to consider other
(weaker) features, making the individual trees less correlated
(\(\rho \downarrow\)).

\subsection{3. The Algorithm}\label{the-algorithm-3}

\begin{enumerate}
\def\labelenumi{\arabic{enumi}.}
\item
  \textbf{Input:} Dataset \(D\), number of trees \(B\), feature subset
  size \(m\).
\item
  \textbf{For \(b = 1\) to \(B\):}

  \begin{enumerate}
  \def\labelenumii{\alph{enumii}.}
  \item
    Draw a bootstrap sample \(Z^*_b\) of size \(N\) from \(D\).
  \item
    Grow a random-forest tree \(T_b\) to the bootstrapped data.
  \end{enumerate}

  \begin{itemize}
  \item
    Recursively repeat until minimum node size is reached.
  \item
    \textbf{Do not prune} the tree (grow it deep/overfit).
  \item
    At each node, select \(m\) variables at random from the \(p\)
    features.
  \item
    Pick the best variable/split-point among the \(m\).
  \end{itemize}
\item
  \textbf{Output:} The ensemble of trees \(\{T_b\}_1^B\).

  \begin{itemize}
  \item
    \textbf{Regression:}
    \(\hat{f}_{rf}(x) = \frac{1}{B} \sum_{b=1}^{B} T_b(x)\)
  \item
    \textbf{Classification:} Majority Vote.
  \end{itemize}
\end{enumerate}

\subsection{4. Mathematical Justification: Variance
Reduction}\label{mathematical-justification-variance-reduction}

Why does this work? Let the prediction of a single tree be a random
variable \(X_i\) with variance \(\sigma^2\).

We average \(B\) trees: \(\bar{X} = \frac{1}{B} \sum X_i\).

If the trees are \textbf{independent} (identically distributed, i.i.d),
the variance of the average decreases linearly:

\[Var(\bar{X}) = \frac{\sigma^2}{B}\]

However, trees in a forest are \textbf{not independent}; they are
positively correlated with coefficient \(\rho\).

The variance of the average of correlated variables is:

\[Var(\bar{X}) = \rho \sigma^2 + \frac{1-\rho}{B} \sigma^2\]

\subsubsection{Interpretation}\label{interpretation}

\begin{itemize}
\item
  \textbf{Term 2 (\(\frac{1-\rho}{B}\sigma^2\)):} Vanishes as
  \(B \to \infty\). Adding more trees eliminates the variance caused by
  the sampling noise.
\item
  \textbf{Term 1 (\(\rho \sigma^2\)):} This is the limit. The variance
  cannot drop below \(\rho \sigma^2\).
\end{itemize}

\textbf{Key Insight:}

\begin{itemize}
\item
  \textbf{Bagging} reduces the second term by increasing \(B\).
\item
  \textbf{Random Forest} reduces the first term by \textbf{reducing
  \(\rho\)} (via feature randomization), allowing the ensemble to
  achieve lower total error than standard bagging.
\end{itemize}

\subsection{5. Out-of-Bag (OOB) Error}\label{out-of-bag-oob-error}

Random Forests have a built-in validation mechanism, eliminating the
need for a separate Hold-out set (Cost-free Cross-Validation).

\subsubsection{5.1 Mechanism}\label{mechanism-3}

For each tree \(T_b\), about 1/3 of the training data was \textbf{not}
used in the bootstrap sample. These are the \textbf{Out-of-Bag} examples
for that tree.

\subsubsection{5.2 OOB Estimation}\label{oob-estimation}

\begin{enumerate}
\def\labelenumi{\arabic{enumi}.}
\item
  For every observation \(x_i\), identify all trees \(T_b\) where
  \(x_i\) was OOB (not in the training set).
\item
  Aggregate the predictions of these specific trees to predict
  \(\hat{y}_i\).
\item
  Calculate the error (OOB Error) based on these aggregated predictions.
\end{enumerate}

\textbf{Theorem:} The OOB error is an unbiased estimate of the
generalization error.

\subsection{6. Feature Importance}\label{feature-importance}

Random Forests are less interpretable than single trees (black box), but
they provide \textbf{Feature Importance} metrics.

\subsubsection{6.1 Mean Decrease Impurity
(MDI)}\label{mean-decrease-impurity-mdi}

Calculated during training.

\begin{itemize}
\item
  Every time a node splits on feature \(j\), we measure the decrease in
  impurity (Gini or Entropy) weighted by the number of samples reaching
  that node.
\item
  We sum these decreases for feature \(j\) across all trees in the
  forest.
\item
  The features with the highest total decrease are the most important.
\end{itemize}

\subsubsection{6.2 Mean Decrease Accuracy (Permutation
Importance)}\label{mean-decrease-accuracy-permutation-importance}

Calculated after training (using OOB data).

\begin{enumerate}
\def\labelenumi{\arabic{enumi}.}
\item
  Record the OOB accuracy of the forest.
\item
  Randomly \textbf{shuffle} the values of feature \(j\) in the OOB data
  (breaking the relationship between \(x_j\) and \(y\)).
\item
  Measure the new OOB accuracy.
\item
  \textbf{Importance:} The drop in accuracy. If shuffling a feature
  doesn't hurt accuracy, that feature was not important.
\end{enumerate}

\subsubsection{Key Takeaways}\label{key-takeaways-18}

\begin{itemize}
\item
  \textbf{Ensemble:} Combines many weak(er) learners to form a strong
  learner.
\item
  \textbf{Variance Reduction:} The primary goal. It turns high-variance
  trees into a low-variance forest.
\item
  \textbf{Bagging + Random Subspace:} The two pillars. Bagging reduces
  variance; Random features reduce correlation (\(\rho\)).
\item
  \textbf{No Overfitting (mostly):} As \(B\) (number of trees)
  increases, the model does \emph{not} overfit. It converges to a stable
  error rate.
\item
  \textbf{Parallelization:} Unlike Boosting (which is sequential),
  Random Forest trees can be trained in parallel on different CPU cores.
\item
  \textbf{OOB Error:} A crucial tool for validation without splitting
  the dataset.

  \newpage
\end{itemize}

\section{K-Nearest Neighbors (KNN)}\label{k-nearest-neighbors-knn}

K-Nearest Neighbors (KNN) is a \textbf{non-parametric},
\textbf{instance-based} (lazy) learning algorithm.

\begin{itemize}
\item
  \textbf{Non-parametric:} It makes no assumptions about the underlying
  data distribution (unlike Linear Regression which assumes linearity).
\item
  \textbf{Instance-based:} It does not learn a fixed model parameters
  \(\theta\) during a training phase. Instead, it simply memorizes the
  entire training dataset \(\mathcal{D}\). ``Training'' is \(O(1)\) (or
  \(O(N)\) for indexing), while prediction is computationally expensive
  \(O(N \cdot d)\).
\end{itemize}

\subsection{1. The Algorithm}\label{the-algorithm-4}

Given a query point \(x_q\) and a training set
\(\mathcal{D} = \{(x_i, y_i)\}_{i=1}^N\):

\begin{enumerate}
\def\labelenumi{\arabic{enumi}.}
\item
  \textbf{Compute Distances:} Calculate the distance \(d(x_q, x_i)\)
  between \(x_q\) and every training example \(x_i\).
\item
  \textbf{Sort:} Find the set \(\mathcal{N}_k\) containing the \(k\)
  training examples with the smallest distances to \(x_q\).
\item
  \textbf{Aggregate:} Combine the labels \(y_i\) of the neighbors in
  \(\mathcal{N}_k\) to predict \(\hat{y}_q\).
\end{enumerate}

\subsubsection{1.1 Classification}\label{classification}

For discrete labels (e.g., \(y \in \{0, 1, \dots, C\}\)), the
aggregation strategy is \textbf{Majority Voting}.

\[\hat{y}_q = \underset{c}{\text{argmax}} \sum_{(x_i, y_i) \in \mathcal{N}_k} \mathbb{I}(y_i = c)\]

\begin{itemize}
\item
  The class with the highest frequency among the \(k\) neighbors wins.
\item
  To break ties (e.g., \(k=2\), one neighbor is Class A, one is Class
  B), we often decrease \(k\) by 1 or weight votes by
  \(1/\text{distance}\).
\end{itemize}

\subsubsection{1.2 Regression}\label{regression}

For continuous targets (e.g., \(y \in \mathbb{R}\)), the aggregation
strategy is \textbf{Averaging}.

\[\hat{y}_q = \frac{1}{k} \sum_{(x_i, y_i) \in \mathcal{N}_k} y_i\]

\begin{itemize}
\item
  \emph{Variation:} \textbf{Weighted KNN Regression}, where neighbors
  closer to the query point contribute more to the average:

  \[\hat{y}_q = \frac{\sum w_i y_i}{\sum w_i} \quad \text{where} \quad w_i = \frac{1}{d(x_q, x_i)^2}\]
\end{itemize}

\subsection{\texorpdfstring{2. The Hyperparameter \(k\): Bias-Variance
Trade-off}{2. The Hyperparameter k: Bias-Variance Trade-off}}\label{the-hyperparameter-k-bias-variance-trade-off}

The choice of \(k\) drastically alters the model's behavior, controlling
the complexity.

\subsubsection{\texorpdfstring{2.1 Case: \(k=1\)
(Overfitting)}{2.1 Case: k=1 (Overfitting)}}\label{case-k1-overfitting}

\begin{itemize}
\item
  \textbf{Mechanism:} The model predicts the label of the \emph{single}
  closest point.
\item
  \textbf{Decision Boundary:} extremely jagged and complex. It forms a
  \textbf{Voronoi Tessellation} around the training points.
\item
  \textbf{Noise Sensitivity:} Highly sensitive. If an outlier (noise)
  exists (e.g., a red point inside a blue cluster), the model creates a
  specific ``island'' of red prediction around that noise point.
\item
  \textbf{Performance:}

  \begin{itemize}
  \item
    \textbf{Training Error:} Always 0 (the nearest neighbor to a
    training point is itself).
  \item
    \textbf{Test Error:} High.
  \item
    \textbf{Regime:} \textbf{Low Bias, High Variance}.
  \end{itemize}
\end{itemize}

\subsubsection{\texorpdfstring{2.2 Case: High \(k\)
(Underfitting)}{2.2 Case: High k (Underfitting)}}\label{case-high-k-underfitting}

\begin{itemize}
\item
  \textbf{Mechanism:} The model averages over a large region of space.
\item
  \textbf{Decision Boundary:} Smooth and simple (approaches a linear
  boundary).
\item
  \textbf{Limit (\(k=N\)):} The model predicts the global majority class
  (or global mean) for \emph{any} input.
\item
  \textbf{Performance:}

  \begin{itemize}
  \item
    \textbf{Training Error:} High.
  \item
    \textbf{Test Error:} High (fails to capture local structure).
  \item
    \textbf{Regime:} \textbf{High Bias, Low Variance}.
  \end{itemize}
\end{itemize}

\subsection{3. Distance Metrics}\label{distance-metrics}

The definition of ``nearest'' depends on the metric used. Let
\(x, z \in \mathbb{R}^d\).

\subsubsection{\texorpdfstring{3.1 Minkowski Distance (\(L_p\)
Norm)}{3.1 Minkowski Distance (L\_p Norm)}}\label{minkowski-distance-l_p-norm}

The generalized form for distance:

\[D(x, z) = \left( \sum_{i=1}^d |x_i - z_i|^p \right)^{1/p}\]

\subsubsection{3.2 Common Special Cases}\label{common-special-cases}

\begin{itemize}
\item
  \textbf{Euclidean Distance (\(p=2\)):}

  \[d(x, z) = \sqrt{\sum_{i=1}^d (x_i - z_i)^2}\]

  \begin{itemize}
  \item
    Standard geometric distance (``as the crow flies'').
  \item
    Sensitive to outliers (due to squaring differences).
  \end{itemize}
\item
  \textbf{Manhattan Distance (\(p=1\)):}

  \[d(x, z) = \sum_{i=1}^d |x_i - z_i|\]

  \begin{itemize}
  \item
    Also called ``Taxicab geometry.''
  \item
    Often preferred in very high-dimensional spaces (less sensitive to
    outliers than \(L_2\)).
  \end{itemize}
\item
  \textbf{Chebyshev Distance (\(p=\infty\)):}

  \[d(x, z) = \max_i |x_i - z_i|\]
\end{itemize}

\subsubsection{3.3 Critical Preprocessing: Feature
Scaling}\label{critical-preprocessing-feature-scaling}

KNN is extremely sensitive to the scale of features.

\begin{itemize}
\item
  \emph{Example:} Feature \(x_1\) is ``Income'' (range 0-100,000) and
  \(x_2\) is ``Age'' (range 0-100).
\item
  The Euclidean distance will be dominated entirely by Income:
  \((100,000 - 90,000)^2 \gg (30 - 40)^2\).
\item
  \textbf{Requirement:} Always normalize (Min-Max) or standardize
  (Z-score) data before applying KNN.
\end{itemize}

\subsection{4. The Curse of
Dimensionality}\label{the-curse-of-dimensionality}

KNN fails catastrophically in high-dimensional spaces (e.g.,
\(d > 20\)). This phenomenon is known as the Curse of Dimensionality
(Bellman, 1961).

\subsubsection{4.1 Theoretical Problem: Distance
Concentration}\label{theoretical-problem-distance-concentration}

As the dimension \(d \to \infty\), the distance between \emph{any} two
random points approaches a constant value. The contrast between the
``nearest'' and ``farthest'' neighbor vanishes.

\textbf{Theorem (Distance Concentration):}

Let \(D_{max}\) be the distance to the farthest point and \(D_{min}\) be
the distance to the nearest point. As \(d \to \infty\):

\[\lim_{d \to \infty} \frac{D_{max} - D_{min}}{D_{min}} \to 0\]

\begin{itemize}
\tightlist
\item
  \textbf{Intuition:} In high dimensions, all points are ``equally far''
  from each other. The concept of ``neighbor'' becomes meaningless.
\end{itemize}

\subsubsection{4.2 Geometric Intuition:
Sparsity}\label{geometric-intuition-sparsity}

The volume of the space increases exponentially with dimension.

\begin{itemize}
\item
  To cover a fraction \(r\) of the volume in 1D, we need \(r\) length.
\item
  To cover the same fraction \(r\) in \(d\) dimensions, we need a
  hypercube of edge length \(e_d = r^{1/d}\).
\item
  \emph{Example:} To capture 10\% of the data (\(r=0.1\)) in \(d=10\)
  dimensions:

  \[e_{10} = (0.1)^{1/10} \approx 0.8\]

  We need to cover 80\% of the range of \emph{every} feature just to
  find 10\% of the neighbors.
\item
  \textbf{Result:} Data becomes incredibly sparse. To maintain the same
  density as in 1D, the number of training samples \(N\) must grow
  exponentially (\(N \propto 10^d\)).
\end{itemize}

\subsubsection{Key Takeaways}\label{key-takeaways-19}

\begin{itemize}
\item
  \textbf{Lazy Learning:} KNN stores data and computes only at query
  time. Fast training, slow prediction.
\item
  \textbf{The Role of \(k\):}

  \begin{itemize}
  \item
    \(k=1\): Memorizes noise, High Variance (Jagged boundary).
  \item
    \(k \to N\): Ignores local structure, High Bias (Flat boundary).
  \end{itemize}
\item
  \textbf{Scaling:} Must normalize data (Z-score) so features with large
  magnitudes don't dominate the distance metric.
\item
  \textbf{Curse of Dimensionality:}

  \begin{itemize}
  \item
    In high \(d\), points become equidistant.
  \item
    Euclidean distance loses meaning.
  \item
    For high-dimensional data (e.g., images, text), standard KNN is
    useless without dimensionality reduction (PCA) or specialized
    metrics (Cosine Similarity).

    \newpage
  \end{itemize}
\end{itemize}

\section{\texorpdfstring{Optimizing \(k\): The ``Elbow'' of the Error
Curve}{Optimizing k: The ``Elbow'' of the Error Curve}}\label{optimizing-k-the-elbow-of-the-error-curve}

Choosing the optimal \(k\) is crucial for the performance of the KNN
algorithm. Since \(k\) controls the model complexity, we cannot learn it
from the training data (setting \(k=1\) always minimizes training error
to 0).

Instead, we use \textbf{Cross-Validation} and visual inspection of the
\textbf{Validation Error Curve}.

\subsection{1. The Error Rate Plot}\label{the-error-rate-plot}

To find the optimal \(k\), we employ an empirical approach:

\begin{enumerate}
\def\labelenumi{\arabic{enumi}.}
\item
  \textbf{Range:} Define a range of possible \(k\) values (e.g.,
  \(k \in \{1, 3, 5, \dots, \sqrt{N}\}\)).
\item
  \textbf{Cross-Validation:} For each \(k\), perform \(M\)-fold
  Cross-Validation and compute the average \textbf{Validation Error}.
\item
  \textbf{Plot:} Plot the Error Rate (Y-axis) against \(k\) (X-axis).
\end{enumerate}

\subsubsection{1.1 Interpreting the Curve}\label{interpreting-the-curve}

The resulting curve typically follows a U-shape (or a reverse-J shape):

\begin{itemize}
\item
  \textbf{Region A (Low \(k\), e.g., 1-3):} High Variance. The error is
  high because the model is capturing noise (overfitting).
\item
  \textbf{Region B (The ``Sweet Spot''):} The error drops as \(k\)
  increases, smoothing out the decision boundary. The curve reaches a
  minimum.
\item
  \textbf{Region C (High \(k\)):} High Bias. The error starts to
  increase again (or plateaus at a high level) because the model becomes
  too simple and ignores local structure (underfitting).
\end{itemize}

\subsection{2. The ``Elbow'' Rule for KNN}\label{the-elbow-rule-for-knn}

While theoretically we simply pick the \(k\) with the \textbf{Global
Minimum} error, in practice, we often use a heuristic similar to the
``Elbow Method'':

\textbf{The Rule:} If the error curve drops sharply at first and then
plateaus (flattens out), we choose the \(k\) value at the ``elbow'' of
the curve---the point where adding more neighbors yields diminishing
returns in accuracy.

\subsubsection{Why not always the Global
Minimum?}\label{why-not-always-the-global-minimum}

\begin{itemize}
\item
  \textbf{Computational Cost:} A smaller \(k\) is computationally faster
  for prediction (fewer neighbors to process/sort).
\item
  \textbf{Parsimony:} If \(k=5\) yields 95.0\% accuracy and \(k=50\)
  yields 95.1\% accuracy, we prefer \(k=5\) (Occham's Razor).
\end{itemize}

\subsection{3. Crucial Distinction: KNN
vs.~K-Means}\label{crucial-distinction-knn-vs.-k-means}

{\def\LTcaptype{none} % do not increment counter
\begin{longtable}[]{@{}
  >{\raggedright\arraybackslash}p{(\linewidth - 4\tabcolsep) * \real{0.3333}}
  >{\raggedright\arraybackslash}p{(\linewidth - 4\tabcolsep) * \real{0.3333}}
  >{\raggedright\arraybackslash}p{(\linewidth - 4\tabcolsep) * \real{0.3333}}@{}}
\toprule\noalign{}
\begin{minipage}[b]{\linewidth}\raggedright
\textbf{Feature}
\end{minipage} & \begin{minipage}[b]{\linewidth}\raggedright
\textbf{KNN (Supervised)}
\end{minipage} & \begin{minipage}[b]{\linewidth}\raggedright
\textbf{K-Means (Unsupervised)}
\end{minipage} \\
\midrule\noalign{}
\endhead
\bottomrule\noalign{}
\endlastfoot
\textbf{Objective} & Classification / Regression & Clustering
(Grouping) \\
\textbf{Y-Axis Metric} & \textbf{Validation Error} (Misclassification
Rate) & \textbf{WCSS} (Within-Cluster Sum of Squares) \\
\textbf{Goal} & Find \(k\) that minimizes Error (U-shape curve). & Find
\(k\) where WCSS stops dropping fast (L-shape/Elbow). \\
\textbf{Typical Curve} & U-shaped (High \(\to\) Low \(\to\) High due to
Bias-Variance tradeoff) & Monotonically Decreasing (Always drops as
\(k \uparrow\)) \\
\end{longtable}
}

\subsubsection{Summary for Exam}\label{summary-for-exam}

\begin{itemize}
\item
  In \textbf{KNN}, we look for the \textbf{minimum} of the validation
  error curve.
\item
  In \textbf{Clustering}, we look for the \textbf{elbow} (kink) of the
  inertia curve.
\end{itemize}

\subsubsection{Key Takeaways}\label{key-takeaways-20}

\begin{itemize}
\item
  \textbf{Start Small:} Test odd values of \(k\) (to avoid ties in
  binary classification).
\item
  \textbf{The Heuristic:} Choose \(k \approx \sqrt{N}\) as a starting
  point, but always validate.
\item
  \textbf{The Plot:} Visualizing Error vs.~\(k\) is the standard way to
  diagnose Bias/Variance in KNN.

  \begin{itemize}
  \item
    Left side of plot (\(k=1\)): High Variance.
  \item
    Right side of plot (\(k=N\)): High Bias.

    \newpage
  \end{itemize}
\end{itemize}

\section{Clustering Algorithms: K-Means \& Hierarchical
Methods}\label{clustering-algorithms-k-means-hierarchical-methods}

Clustering is an \textbf{Unsupervised Learning} task where the goal is
to group similar data points together based on a distance metric,
without using any external labels (\(y\)).

\subsection{1. K-Means Clustering}\label{k-means-clustering}

K-Means is a centroid-based algorithm that partitions a dataset
\(\mathcal{D}\) into \(K\) distinct, non-overlapping clusters.

\subsubsection{1.1 The Algorithm (Lloyd's
Algorithm)}\label{the-algorithm-lloyds-algorithm}

\begin{enumerate}
\def\labelenumi{\arabic{enumi}.}
\item
  \textbf{Initialization:} Randomly select \(K\) points as initial
  centroids \(\mu_1, \dots, \mu_K\).
\item
  \textbf{Assignment Step:} Assign every point \(x_i\) to the closest
  centroid.

  \[c^{(i)} := \underset{j}{\text{argmin}} \ ||x_i - \mu_j||^2\]
\item
  \textbf{Update Step:} Recompute the centroids by taking the mean of
  all points assigned to that cluster.

  \[\mu_j := \frac{1}{|S_j|} \sum_{i \in S_j} x_i\]
\item
  \textbf{Convergence:} Repeat steps 2 and 3 until centroids do not
  change significantly or a maximum number of iterations is reached.
\end{enumerate}

\subsubsection{1.2 Objective Function
(Inertia)}\label{objective-function-inertia}

K-Means minimizes the \textbf{Within-Cluster Sum of Squares (WCSS)}:

\[J(c, \mu) = \sum_{i=1}^N ||x_i - \mu_{c^{(i)}}||^2\]

This creates clusters that are as compact (spherical) as possible.

\subsubsection{1.3 Limitations}\label{limitations}

\begin{itemize}
\item
  \textbf{Convex Shapes:} It assumes clusters are convex (spherical). It
  fails on elongated or irregular shapes (e.g., concentric circles).
\item
  \textbf{Initialization Sensitivity:} Random initialization can lead to
  bad local minima. \textbf{K-Means++} is used to initialize centroids
  far apart to mitigate this.
\item
  \textbf{Fixed K:} Requires specifying \(K\) in advance (use Elbow
  Method or Silhouette Score to find it).
\end{itemize}

\subsection{2. Hierarchical Clustering}\label{hierarchical-clustering}

Unlike K-Means, Hierarchical Clustering does not require specifying
\(K\) upfront. It builds a hierarchy of clusters, usually represented as
a tree structure called a \textbf{Dendrogram}.

There are two approaches:

\begin{enumerate}
\def\labelenumi{\arabic{enumi}.}
\item
  \textbf{Agglomerative (Bottom-Up):} Start with \(N\) clusters (each
  point is a cluster). Merge the two closest clusters iteratively.
\item
  \textbf{Divisive (Top-Down):} Start with 1 cluster. Split iteratively.
\end{enumerate}

\subsubsection{2.1 The ``Linkage'' Problem}\label{the-linkage-problem}

In Agglomerative clustering, the key decision is \textbf{how to measure
the distance between two clusters} \(A\) and \(B\). This metric is
called \textbf{Linkage}.

\subsubsection{2.2 Single Linkage (Min
Linkage)}\label{single-linkage-min-linkage}

The distance between two clusters is the distance between their
\textbf{closest} points.

\[D(A, B) = \min \{ d(a, b) : a \in A, b \in B \}\]

\begin{itemize}
\item
  \textbf{Behavior:} It can follow ``chains'' of points.
\item
  \textbf{Pros:} Can detect non-elliptical, irregular shapes (e.g., two
  concentric half-moons).
\item
  \textbf{Cons:} Suffer from the \textbf{Chaining Effect}. A single
  noise point between two distinct clusters can merge them into one.
  Very sensitive to outliers.
\end{itemize}

\subsubsection{2.3 Complete Linkage (Max
Linkage)}\label{complete-linkage-max-linkage}

The distance between two clusters is the distance between their
\textbf{farthest} points.

\[D(A, B) = \max \{ d(a, b) : a \in A, b \in B \}\]

\begin{itemize}
\item
  \textbf{Behavior:} It prefers compact, spherical clusters with similar
  diameters.
\item
  \textbf{Pros:} Less sensitive to outliers/noise than Single Linkage.
  Avoids chaining.
\item
  \textbf{Cons:} Tends to break large clusters; biased towards finding
  globular shapes.
\end{itemize}

\subsubsection{2.4 Average Linkage}\label{average-linkage}

The distance is the \textbf{average} distance between all pairs of
points in the two clusters.

\[D(A, B) = \frac{1}{|A| \cdot |B|} \sum_{a \in A} \sum_{b \in B} d(a, b)\]

\begin{itemize}
\item
  \textbf{Behavior:} A compromise between Single and Complete.
\item
  \textbf{Pros:} Robust to noise and outliers.
\item
  \textbf{Cons:} Computationally expensive (\(O(N^2)\) distance
  calculations).
\end{itemize}

\subsubsection{2.5 Ward's Linkage}\label{wards-linkage}

Minimizes the variance (increase in WCSS) when merging two clusters.
This makes Hierarchical Clustering behave similarly to K-Means (favoring
equal-sized, spherical clusters).

\subsection{3. Performance and
Comparison}\label{performance-and-comparison}

{\def\LTcaptype{none} % do not increment counter
\begin{longtable}[]{@{}
  >{\raggedright\arraybackslash}p{(\linewidth - 4\tabcolsep) * \real{0.3333}}
  >{\raggedright\arraybackslash}p{(\linewidth - 4\tabcolsep) * \real{0.3333}}
  >{\raggedright\arraybackslash}p{(\linewidth - 4\tabcolsep) * \real{0.3333}}@{}}
\toprule\noalign{}
\begin{minipage}[b]{\linewidth}\raggedright
\textbf{Feature}
\end{minipage} & \begin{minipage}[b]{\linewidth}\raggedright
\textbf{K-Means}
\end{minipage} & \begin{minipage}[b]{\linewidth}\raggedright
\textbf{Hierarchical (Agglomerative)}
\end{minipage} \\
\midrule\noalign{}
\endhead
\bottomrule\noalign{}
\endlastfoot
\textbf{Complexity} & \(O(N \cdot K \cdot I)\) (Linear). Very Fast. &
\(O(N^2 \log N)\) or \(O(N^3)\). Very Slow for large \(N\). \\
\textbf{Shape Assumption} & Spherical (Convex). & Arbitrary (Single
Linkage) or Spherical (Ward/Complete). \\
\textbf{Parameters} & Must define \(K\) upfront. & No \(K\) needed; cut
the dendrogram at any height. \\
\textbf{Output} & Flat partition. & Hierarchy (Dendrogram). \\
\textbf{Determinism} & Stochastic (Random Initialization). &
Deterministic. \\
\end{longtable}
}

\subsubsection{3.1 Evaluation Metrics
(Internal)}\label{evaluation-metrics-internal}

Since we lack labels, we use internal metrics to evaluate cluster
quality:

\begin{itemize}
\item
  \textbf{Silhouette Coefficient:} Measures how similar a point is to
  its own cluster (cohesion) compared to other clusters (separation).
  Range \([-1, 1]\).

  \begin{itemize}
  \item
    \(+1\): Correctly clustered.
  \item
    \(0\): Overlapping clusters.
  \item
    \(-1\): Incorrectly clustered.
  \end{itemize}
\item
  \textbf{Davies-Bouldin Index:} The average similarity ratio of each
  cluster with its most similar cluster. Lower is better.
\end{itemize}

\subsubsection{Key Takeaways}\label{key-takeaways-21}

\begin{itemize}
\item
  \textbf{K-Means:} Best for general-purpose clustering on large
  datasets where clusters are roughly spherical.
\item
  \textbf{Hierarchical:} Best for small datasets (\(N < 10,000\)) or
  when the hierarchy/taxonomy is important (e.g., biological species).
\item
  \textbf{Linkage matters:}

  \begin{itemize}
  \item
    \textbf{Single:} Good for weird shapes, bad for noise (Chaining).
  \item
    \textbf{Complete:} Good for compact blobs, breaks large clusters.
  \item
    \textbf{Average/Ward:} Balanced and robust.
  \end{itemize}
\item
  \textbf{Dendrogram:} The visual tool used to decide the number of
  clusters in hierarchical methods by making a horizontal ``cut.''

  \newpage
\end{itemize}

\section{Feedforward Neural Networks
(FNNs)}\label{feedforward-neural-networks-fnns}

A Feedforward Neural Network is the foundational architecture of Deep
Learning. It is a biologically inspired computational model capable of
learning complex non-linear approximations of functions.

``Feedforward'' means information flows in one direction: from Input
\(\to\) Hidden Layers \(\to\) Output. There are no loops (unlike
Recurrent Neural Networks).

\subsection{1. Architecture and Forward
Propagation}\label{architecture-and-forward-propagation}

An FNN consists of layers of composed non-linear functions.

\subsubsection{1.1 The Neuron (Perceptron)}\label{the-neuron-perceptron}

The basic unit is a neuron. It performs two steps:

\begin{enumerate}
\def\labelenumi{\arabic{enumi}.}
\item
  \textbf{Linear Transformation:} Computes a weighted sum of inputs plus
  a bias.

  \[z = w^T x + b\]
\item
  \textbf{Activation:} Applies a non-linear function \(\sigma(\cdot)\).

  \[a = \sigma(z)\]
\end{enumerate}

\subsubsection{1.2 The Network Structure}\label{the-network-structure}

Let \(L\) be the number of layers.

\begin{itemize}
\item
  \textbf{Layer 0:} Input layer \(a^{(0)} = x\).
\item
  \textbf{Hidden Layers (\(l = 1 \dots L-1\)):} Perform feature
  extraction.
\item
  \textbf{Output Layer (\(L\)):} Produces the final prediction
  \(\hat{y}\).
\end{itemize}

\subsubsection{1.3 Mathematical Formulation
(Vectorized)}\label{mathematical-formulation-vectorized}

For any layer \(l\), let \(W^{(l)}\) be the weight matrix and
\(b^{(l)}\) be the bias vector.

The Forward Propagation steps are:

\begin{enumerate}
\def\labelenumi{\arabic{enumi}.}
\item
  \textbf{Pre-activation:}

  \[z^{(l)} = W^{(l)} a^{(l-1)} + b^{(l)}\]
\item
  \textbf{Activation:}

  \[a^{(l)} = \sigma(z^{(l)})\]
\end{enumerate}

\subsubsection{\texorpdfstring{1.4 The Importance of Non-Linearity
(\(\sigma\))}{1.4 The Importance of Non-Linearity (\textbackslash sigma)}}\label{the-importance-of-non-linearity-sigma}

If we did not use non-linear activation functions (like ReLU, Sigmoid,
Tanh), the entire network would collapse into a single linear
transformation (matrix multiplication is associative).

\[W_2(W_1 x) = (W_2 W_1)x = W_{new} x\]

\textbf{Non-linearity allows the network to learn complex boundaries.}

\subsection{2. Geometric Intuition}\label{geometric-intuition}

How does an FNN solve non-linear problems (like XOR or concentric
circles) using linear operations (\(Wx+b\))?

\subsubsection{2.1 Feature Space
Transformation}\label{feature-space-transformation}

We can view each hidden layer as performing a \textbf{Change of Basis}
(coordinate transformation).

\begin{itemize}
\item
  \textbf{Input Space:} Data classes are entangled and not linearly
  separable.
\item
  \textbf{Hidden Space:} The network warps, stretches, and folds the
  space.
\item
  \textbf{Output Space:} The manifold is ``unraveled'' such that the
  classes become linearly separable by the final layer.
\end{itemize}

\subsubsection{2.2 The Universal Approximation
Theorem}\label{the-universal-approximation-theorem}

\textbf{Theorem (Cybenko, 1989):} A feedforward network with a single
hidden layer containing a finite number of neurons can approximate
continuous functions on compact subsets of \(\mathbb{R}^n\), under mild
assumptions on the activation function.

\begin{itemize}
\item
  \textbf{Intuition:} FNNs are ``Universal Function Approximators.''
\item
  \textbf{Caveat:} The theorem says a solution \emph{exists}, but it
  doesn't guarantee that we can \emph{find} it (learn it) or that the
  layer size is feasible.
\end{itemize}

\subsection{3. VC Dimension of Neural
Networks}\label{vc-dimension-of-neural-networks}

The Vapnik-Chervonenkis (VC) dimension measures the capacity of the
network (its ability to shatter points).

\subsubsection{3.1 Estimates}\label{estimates}

For a network with \(W\) total weights (parameters) and \(L\) layers
using ReLU or Sigmoid activations:

\begin{itemize}
\item
  The VC dimension is roughly proportional to the number of weights.

  \[VC_{dim} \approx O(W \cdot \log W)\]
\item
  Or sometimes cited as \(O(W \cdot L)\).
\end{itemize}

\subsubsection{3.2 Implication: The Capacity
Problem}\label{implication-the-capacity-problem}

\begin{itemize}
\item
  Neural Networks often have \(W \gg N\) (more parameters than data
  points).
\item
  High VC Dimension \(\implies\) \textbf{Massive Overfitting Risk}.
\item
  \textbf{Deep Learning Paradox:} Despite \(VC \to \infty\), deep nets
  generalize well. This is usually achieved through strong
  \textbf{Regularization} (Dropout, Weight Decay, Early Stopping) and
  the inductive bias of Stochastic Gradient Descent.
\end{itemize}

\subsection{4. Backpropagation (The Learning
Algorithm)}\label{backpropagation-the-learning-algorithm}

We train the network by minimizing a Cost Function \(J(\theta)\) using
Gradient Descent.

\textbf{Backpropagation} is the algorithm used to efficiently calculate
the gradient \(\nabla_\theta J(\theta)\). It is essentially the
\textbf{Chain Rule of Calculus} applied recursively.

\subsubsection{4.1 The Objective}\label{the-objective-2}

We want to find how much the Total Error (\(J\)) changes if we change a
specific weight \(W_{ij}^{(l)}\).

\[\frac{\partial J}{\partial W_{ij}^{(l)}}\]

\subsubsection{4.2 The Four Fundamental
Equations}\label{the-four-fundamental-equations}

Let \(\delta^{(l)}\) be the ``error term'' (gradient of cost w.r.t.
\(z\)) for layer \(l\).

\begin{enumerate}
\def\labelenumi{\arabic{enumi}.}
\item
  \textbf{Error at Output Layer (\(L\)):}

  Depends on the Loss function (e.g., MSE or Cross-Entropy).

  \[\delta^{(L)} = \nabla_a J \odot \sigma'(z^{(L)})\]

  \emph{(Where \(\odot\) is the Hadamard/element-wise product)}.
\item
  \textbf{Error Propagated to Hidden Layers (\(l\)):}

  We pull the error backward from layer \(l+1\) through the weights
  \(W^{(l+1)}\).

  \[\delta^{(l)} = ((W^{(l+1)})^T \delta^{(l+1)}) \odot \sigma'(z^{(l)})\]

  \emph{Intuition: The error at layer \(l\) is the weighted sum of
  errors at \(l+1\), gated by the derivative of the activation
  function.}
\item
  \textbf{Gradient of Weights:}

  \[\frac{\partial J}{\partial W^{(l)}} = \delta^{(l)} (a^{(l-1)})^T\]
\item
  \textbf{Gradient of Biases:}

  \[\frac{\partial J}{\partial b^{(l)}} = \delta^{(l)}\]
\end{enumerate}

\subsubsection{4.3 The Algorithm
Step-by-Step}\label{the-algorithm-step-by-step}

\begin{enumerate}
\def\labelenumi{\arabic{enumi}.}
\item
  \textbf{Forward Pass:} Compute all activations \(a^{(l)}\) and inputs
  \(z^{(l)}\) up to output.
\item
  \textbf{Compute Loss:} \(J(\hat{y}, y)\).
\item
  \textbf{Backward Pass:}

  \begin{itemize}
  \item
    Compute \(\delta^{(L)}\) at the output.
  \item
    Iteratively compute \(\delta^{(L-1)}, \dots, \delta^{(1)}\) moving
    backwards.
  \end{itemize}
\item
  \textbf{Update Parameters (Gradient Descent):}

  \[W^{(l)} := W^{(l)} - \alpha \frac{\partial J}{\partial W^{(l)}}\]
\end{enumerate}

\subsubsection{Key Takeaways}\label{key-takeaways-22}

\begin{itemize}
\item
  \textbf{Architecture:} Input \(\to\) Hidden (Non-linear Features)
  \(\to\) Output.
\item
  \textbf{Non-linearity:} Essential. Without \(\sigma(\cdot)\), a Deep
  Network is just Linear Regression.
\item
  \textbf{Geometric View:} Layers perform a diffeomorphic transformation
  to make data linearly separable.
\item
  \textbf{Capacity:} High VC dimension (\(\approx\) num\_params).
  Requires regularization.
\item
  \textbf{Backprop:} The efficient application of the Chain Rule to
  compute gradients \(\frac{\partial J}{\partial W}\). It propagates the
  error signal from output to input.****

  \newpage
\end{itemize}

\section{Convolutional Neural Networks
(CNNs)}\label{convolutional-neural-networks-cnns}

Convolutional Neural Networks (CNNs) are a specialized class of deep
neural networks designed for processing data that has a known grid-like
topology, such as time-series data (1D grid) and image data (2D grid).

Unlike FNNs, which treat inputs as independent vectors, CNNs exploit the
\textbf{spatial structure} of the input.

\subsection{1. FNN vs.~CNN: Why do we need
ConvNets?}\label{fnn-vs.-cnn-why-do-we-need-convnets}

If we were to use a standard Feedforward Neural Network (FNN) to
classify a high-resolution image, we would face insurmountable problems.

\subsubsection{1.1 The Problems with FNNs on
Images}\label{the-problems-with-fnns-on-images}

\begin{enumerate}
\def\labelenumi{\arabic{enumi}.}
\item
  \textbf{Parameter Explosion:}

  \begin{itemize}
  \item
    Consider an image of size \(1000 \times 1000 \times 3\) (RGB). The
    input vector \(x\) has \(3,000,000\) elements.
  \item
    If the first hidden layer has 1,000 neurons, the weight matrix \(W\)
    would have \(3 \times 10^9\) parameters.
  \item
    \textbf{Result:} Impossible to train; massive overfitting.
  \end{itemize}
\item
  \textbf{Loss of Spatial Topology:}

  \begin{itemize}
  \item
    FNNs require ``flattening'' the image into a 1D vector.
  \item
    This destroys the spatial relationships (e.g., a pixel at \((0,0)\)
    is related to \((0,1)\), but in a flattened vector, this proximity
    is lost).
  \end{itemize}
\item
  \textbf{Translation Variance:}

  \begin{itemize}
  \tightlist
  \item
    If a cat appears in the top-left corner during training, an FNN
    might not recognize the same cat in the bottom-right corner.
  \end{itemize}
\end{enumerate}

\subsubsection{1.2 The CNN Solution}\label{the-cnn-solution}

CNNs introduce three key ideas to solve this:

\begin{enumerate}
\def\labelenumi{\arabic{enumi}.}
\item
  \textbf{Sparse Connectivity:} A neuron is connected only to a small
  local region of the input (Local Receptive Field).
\item
  \textbf{Parameter Sharing:} The \emph{same} weights (filter) are used
  across the entire image. If a feature (e.g., a vertical edge) is
  useful in one part of the image, it is useful everywhere.
\item
  \textbf{Equivariance to Translation:} If the input shifts, the output
  feature map shifts by the same amount.
\end{enumerate}

{\def\LTcaptype{none} % do not increment counter
\begin{longtable}[]{@{}
  >{\raggedright\arraybackslash}p{(\linewidth - 4\tabcolsep) * \real{0.3333}}
  >{\raggedright\arraybackslash}p{(\linewidth - 4\tabcolsep) * \real{0.3333}}
  >{\raggedright\arraybackslash}p{(\linewidth - 4\tabcolsep) * \real{0.3333}}@{}}
\toprule\noalign{}
\begin{minipage}[b]{\linewidth}\raggedright
\textbf{Feature}
\end{minipage} & \begin{minipage}[b]{\linewidth}\raggedright
\textbf{FNN (Dense)}
\end{minipage} & \begin{minipage}[b]{\linewidth}\raggedright
\textbf{CNN (Convolutional)}
\end{minipage} \\
\midrule\noalign{}
\endhead
\bottomrule\noalign{}
\endlastfoot
\textbf{Connectivity} & All-to-All (Dense) & Local (Sparse
Connectivity) \\
\textbf{Parameters} & Unique weight for every connection & Shared
weights (Kernels/Filters) \\
\textbf{Input Structure} & Ignored (Input is flattened) & Preserved
(2D/3D Grid Structure) \\
\textbf{Invariance} & None (Spatial info is lost) & Translation
Equivariance (via Weight Sharing) \\
\end{longtable}
}

\subsection{2. The Convolutional Layer
(CONV)}\label{the-convolutional-layer-conv}

The core building block of a CNN. It performs the feature extraction.

\subsubsection{2.1 Filters (Kernels)}\label{filters-kernels}

A \textbf{Filter} (or Kernel) \(K\) is a small learnable tensor of
weights (e.g., \(3 \times 3 \times Depth\)).

\begin{itemize}
\item
  \textbf{Depth:} The filter depth must match the input depth. For an
  RGB image, a \(3 \times 3\) filter is actually
  \(3 \times 3 \times 3\).
\item
  \textbf{Role:} Each filter learns to detect a specific pattern (edges,
  textures, blobs).
\end{itemize}

\subsubsection{2.2 The Convolution
Operation}\label{the-convolution-operation}

We ``slide'' the filter over the input image \(I\). At every position,
we compute the \textbf{element-wise product} between the filter and the
overlapping input patch, then sum the results to get a single number.

\textbf{Mathematical Formula (2D):}

\[S(i, j) = (I * K)(i, j) = \sum_{m} \sum_{n} I(i+m, j+n) K(m, n)\]

The output of this operation is a \textbf{Feature Map} (or Activation
Map).

\subsubsection{2.3 Multiple Filters}\label{multiple-filters}

A generic CNN layer does not use just one filter. It uses \(F\)
different filters.

\begin{itemize}
\item
  Input: \(H \times W \times C_{in}\)
\item
  Filters: \(F\) kernels of size \(k \times k \times C_{in}\)
\item
  Output: \(H_{out} \times W_{out} \times F\)
\end{itemize}

\subsection{3. Hyperparameters: Stride, Padding,
Dilation}\label{hyperparameters-stride-padding-dilation}

The dimensions of the output feature map are controlled by three
hyperparameters.

\subsubsection{\texorpdfstring{3.1 Stride
(\(S\))}{3.1 Stride (S)}}\label{stride-s}

Stride controls the \textbf{step size} of the filter as it slides.

\begin{itemize}
\item
  \textbf{Stride = 1:} The filter moves 1 pixel at a time. Output size
  is roughly equal to input size.
\item
  \textbf{Stride = 2:} The filter jumps 2 pixels. This
  \textbf{downsamples} the spatial dimensions by a factor of 2 (halves
  the height and width).
\item
  \textbf{Function:} An alternative to Pooling for dimensionality
  reduction.
\end{itemize}

\subsubsection{\texorpdfstring{3.2 Padding
(\(P\))}{3.2 Padding (P)}}\label{padding-p}

As we convolve, the borders of the image get ``eaten away,'' and the
spatial size decreases (\(N - k + 1\)).

\begin{itemize}
\item
  \textbf{Valid Padding:} No padding. Drop the borders. Output shrinks.
\item
  \textbf{Same Padding:} Pad the input with zeros around the border so
  that Output Size = Input Size (when Stride=1). Preserves spatial
  information.
\end{itemize}

\textbf{Output Dimension Formula:}

\[W_{out} = \lfloor \frac{W_{in} - k + 2P}{S} \rfloor + 1\]

\subsubsection{\texorpdfstring{3.3 Dilation
(\(D\))}{3.3 Dilation (D)}}\label{dilation-d}

Also known as ``Atrous Convolution.'' It introduces gaps (holes) between
the kernel elements.

\begin{itemize}
\item
  \textbf{Standard Convolution:} Dilation = 1.
\item
  \textbf{Dilated Convolution:} The kernel covers a larger area of the
  input without increasing the number of parameters.
\item
  \textbf{Purpose:} Expands the \textbf{Receptive Field} exponentially.
  Used heavily in Semantic Segmentation to capture global context
  without losing resolution.
\end{itemize}

\subsection{4. The Pooling Layer (POOL)}\label{the-pooling-layer-pool}

Pooling layers are usually inserted between successive Conv layers.
Their function is to progressively reduce the spatial size
(downsampling) to reduce the amount of parameters and computation.

\subsubsection{4.1 Max Pooling}\label{max-pooling}

Selects the \textbf{maximum} value from the window (e.g., \(2 \times 2\)
region).

\begin{itemize}
\item
  \textbf{Intuition:} ``If a feature was detected anywhere in this
  region, keep it.'' It captures the most prominent features.
\item
  \textbf{Invariant:} Provides small invariance to translation and
  rotation.
\end{itemize}

\subsubsection{4.2 Average Pooling}\label{average-pooling}

Calculates the average value in the window.

\begin{itemize}
\tightlist
\item
  \textbf{Intuition:} ``Smooths'' the features. Used less frequently in
  modern intermediate layers, but often used as \textbf{Global Average
  Pooling} before the final classification layer.
\end{itemize}

\subsection{5. Typical Architecture: Putting it
Together}\label{typical-architecture-putting-it-together}

A classic CNN (like LeNet-5, AlexNet, VGG) follows a predictable
pattern:

\begin{enumerate}
\def\labelenumi{\arabic{enumi}.}
\item
  \textbf{Input:} Raw pixel values (\(H \times W \times 3\)).
\item
  \textbf{Feature Learning (The ``Body''):}

  \begin{itemize}
  \item
    \textbf{CONV:} Learn features.
  \item
    \textbf{ReLU:} Apply non-linearity (\(\max(0, x)\)).
  \item
    \textbf{POOL:} Downsample.
  \item
    \emph{(Repeat Stack multiple times)}
  \end{itemize}
\item
  \textbf{Classification (The ``Head''):}

  \begin{itemize}
  \item
    \textbf{Flatten:} Convert the 3D tensor (\(h \times w \times c\))
    into a 1D vector.
  \item
    \textbf{Fully Connected (FC):} Standard FNN layers to combine
    features.
  \item
    \textbf{Softmax:} Output probabilities for \(N\) classes.
  \end{itemize}
\end{enumerate}

\subsubsection{Key Takeaways}\label{key-takeaways-23}

\begin{itemize}
\item
  \textbf{Spatial Hierarchy:} CNNs learn a hierarchy of features.

  \begin{itemize}
  \item
    Early layers: Edges, colors, gradients.
  \item
    Middle layers: Textures, corners, parts of objects (eyes, wheels).
  \item
    Deep layers: Whole objects (faces, cars).
  \end{itemize}
\item
  \textbf{Parameter Sharing:} The main reason CNNs are efficient. The
  same edge detector is useful everywhere.
\item
  \textbf{Stride:} Downsampling mechanism (aggressive step).
\item
  \textbf{Dilation:} Expanding the field of view without extra cost.
\item
  \textbf{Pooling:} Makes the representation robust to small shifts and
  distortions.

  \newpage
\end{itemize}

\section{Activation Functions in Neural
Networks}\label{activation-functions-in-neural-networks}

Activation functions are the mathematical components that introduce
\textbf{non-linearity} into neural networks. Without them, a neural
network---regardless of how many layers it has---would behave exactly
like a single Linear Regression model.

\subsection{1. Why Are They Necessary?}\label{why-are-they-necessary}

Consider a two-layer neural network with \textbf{linear} activations
(identity function):

\begin{enumerate}
\def\labelenumi{\arabic{enumi}.}
\item
  \(z^{(1)} = W^{(1)} x + b^{(1)}\)
\item
  \(\hat{y} = W^{(2)} z^{(1)} + b^{(2)}\)
\end{enumerate}

If we substitute the first equation into the second:

\[\hat{y} = W^{(2)} (W^{(1)} x + b^{(1)}) + b^{(2)} = (W^{(2)} W^{(1)}) x + (W^{(2)} b^{(1)} + b^{(2)})\]

Since the product of two matrices \(W^{(2)} W^{(1)}\) is simply another
matrix \(W'\), the entire network collapses into a single linear
transformation: \(\hat{y} = W' x + b'\).

To learn complex patterns (e.g., separating data distributed in
concentric circles or spirals), we must apply a non-linear function
\(\sigma(\cdot)\) after every affine transformation.

\subsection{2. Classical Functions (and their
problems)}\label{classical-functions-and-their-problems}

Before the advent of modern Deep Learning, the most common functions
were the Sigmoid and Hyperbolic Tangent.

\subsubsection{2.1 Sigmoid (Logistic)
Function}\label{sigmoid-logistic-function}

\[\sigma(z) = \frac{1}{1 + e^{-z}}\]

\begin{itemize}
\item
  \textbf{Range:} \((0, 1)\).
\item
  \textbf{Usage:} Probabilistic interpretation (output layer for binary
  classification).
\item
  \textbf{Problem 1 (Vanishing Gradient):}

  \begin{itemize}
  \item
    When \(|z|\) is large (very positive or very negative), the curve
    flattens out (``saturates'').
  \item
    The derivative \(\sigma'(z) \approx 0\).
  \item
    During Backpropagation, we multiply local gradients via the Chain
    Rule. If we multiply many small numbers close to zero, the gradient
    ``vanishes'' before reaching the early layers. The network stops
    learning.
  \end{itemize}
\item
  \textbf{Problem 2 (Not Zero-Centered):} The output is always positive.
  This can cause ``zig-zagging'' dynamics during gradient descent
  optimization.
\end{itemize}

\subsubsection{2.2 Hyperbolic Tangent
(Tanh)}\label{hyperbolic-tangent-tanh}

\[\tanh(z) = \frac{e^z - e^{-z}}{e^z + e^{-z}}\]

\begin{itemize}
\item
  \textbf{Range:} \((-1, 1)\).
\item
  \textbf{Improvement:} It is \textbf{zero-centered} (mean output is
  close to 0), which makes optimization easier than Sigmoid.
\item
  \textbf{Flaw:} It still suffers heavily from the \textbf{Vanishing
  Gradient} problem in the saturation regions.
\end{itemize}

\subsection{3. ReLU (Rectified Linear
Unit)}\label{relu-rectified-linear-unit}

ReLU has become the default activation function for \textbf{hidden
layers} in Deep Learning.

\subsubsection{3.1 Mathematical
Definition}\label{mathematical-definition-1}

\[f(z) = \max(0, z)\]

It is a piecewise linear function:

\begin{itemize}
\item
  If \(z < 0\), \(f(z) = 0\).
\item
  If \(z \geq 0\), \(f(z) = z\) (identity).
\end{itemize}

\subsubsection{3.2 Why is ReLU Chosen?
(Advantages)}\label{why-is-relu-chosen-advantages}

\begin{enumerate}
\def\labelenumi{\arabic{enumi}.}
\item
  \textbf{Solution to Vanishing Gradient (in positive domain):}

  \begin{itemize}
  \item
    The derivative of ReLU is:

    \[f'(z) = \begin{cases} 1 & \text{if } z > 0 \\ 0 & \text{if } z < 0 \end{cases}\]
  \item
    For all positive inputs, the gradient is exactly \textbf{1}. It does
    not decay, no matter how deep the network is. This enables the
    training of very Deep Networks.
  \end{itemize}
\item
  \textbf{Computational Efficiency:}

  \begin{itemize}
  \item
    It does not require expensive exponential calculations (\(e^x\))
    like Sigmoid or Tanh.
  \item
    It is a simple threshold operation:
    \texttt{if\ z\ \textgreater{}\ 0\ return\ z\ else\ 0}.
  \end{itemize}
\item
  \textbf{Sparsity:}

  \begin{itemize}
  \item
    Since ReLU forces all negative inputs to zero, a significant portion
    of neurons in a layer will have 0 activation.
  \item
    This leads to \textbf{sparse representations} (few active neurons
    for a given input), which are often biologically plausible and
    computationally efficient.
  \end{itemize}
\end{enumerate}

\subsubsection{3.3 The Problem: ``Dying
ReLU''}\label{the-problem-dying-relu}

A potential disadvantage is the death of neurons.

\begin{itemize}
\item
  If a neuron falls into a configuration of weights such that for
  \emph{all} inputs in the dataset \(W^T x + b < 0\), the output will
  always be 0.
\item
  Consequently, the gradient will always be 0 (\(f'(z)=0\)).
\item
  The weights of that neuron will never be updated again. The neuron is
  ``dead.''
\item
  \textbf{Cause:} High Learning Rate or unlucky initialization.
\end{itemize}

\subsection{4. Variants of ReLU}\label{variants-of-relu}

To mitigate the ``Dying ReLU'' problem, several variants have been
introduced.

\subsubsection{4.1 Leaky ReLU}\label{leaky-relu}

Instead of absolute zero for \(z < 0\), we allow a small slope.

\[f(z) = \max(\alpha z, z)\]

\begin{itemize}
\item
  Where \(\alpha\) is a small constant (e.g., \(\alpha = 0.01\)).
\item
  \textbf{Advantage:} The gradient for \(z < 0\) is \(\alpha\)
  (non-zero). The neuron can still ``learn'' and recover even if
  inactive.
\end{itemize}

\subsubsection{4.2 Parametric ReLU (PReLU)}\label{parametric-relu-prelu}

Similar to Leaky ReLU, but \(\alpha\) is not a fixed hyperparameter; it
is a \textbf{learnable parameter} that the network adjusts via
Backpropagation.

\subsubsection{4.3 ELU (Exponential Linear
Unit)}\label{elu-exponential-linear-unit}

\[f(z) = \begin{cases} z & \text{if } z > 0 \\ \alpha (e^z - 1) & \text{if } z \leq 0 \end{cases}\]

\begin{itemize}
\item
  Combines the benefits of ReLU (no saturation on the right) with Tanh
  (zero-centered, negative outputs).
\item
  \textbf{Cons:} Computationally more expensive (requires \texttt{exp}).
\end{itemize}

\subsubsection{4.4 GELU / Swish (Modern)}\label{gelu-swish-modern}

Used in Transformer architectures (e.g., BERT, GPT). They are ``smooth''
versions of ReLU.

\begin{itemize}
\item
  \textbf{Swish:} \(f(z) = z \cdot \sigma(z)\)
\item
  Non-monotonic, smooth derivative everywhere. Often outperforms ReLU in
  very deep networks.
\end{itemize}

\subsection{5. Output Layer Functions}\label{output-layer-functions}

While ReLU is standard for hidden layers, the final layer requires
specific functions based on the task.

\begin{itemize}
\item
  \textbf{Regression:} No activation (Identity) or ReLU (if output must
  be positive).
\item
  \textbf{Binary Classification:} Sigmoid.
\item
  \textbf{Multi-class Classification:} \textbf{Softmax}.
\end{itemize}

\subsubsection{5.1 Softmax}\label{softmax}

A generalization of Sigmoid for \(K\) classes. It transforms a vector of
real numbers (logits) into a probability distribution.

\[\text{Softmax}(z)_i = \frac{e^{z_i}}{\sum_{j=1}^K e^{z_j}}\]

\begin{itemize}
\tightlist
\item
  Ensures that \(\sum \hat{y}_i = 1\) and \(\hat{y}_i \geq 0\).
\end{itemize}

\subsubsection{Key Takeaways}\label{key-takeaways-24}

\begin{itemize}
\item
  \textbf{Necessity:} Activations break linearity, enabling the learning
  of complex functions.
\item
  \textbf{Sigmoid/Tanh:} Historically significant but avoided in hidden
  layers due to \textbf{Vanishing Gradient}.
\item
  \textbf{ReLU:} The industry standard.

  \begin{itemize}
  \item
    \textbf{Pros:} Fast, constant gradient (1) for \(x>0\), promotes
    sparsity.
  \item
    \textbf{Cons:} Dying ReLU problem (dead neurons).
  \end{itemize}
\item
  \textbf{Leaky ReLU:} Simple variant to prevent dead neurons.
\item
  \textbf{Rule of Thumb:} Start with \textbf{ReLU}. If the network
  doesn't converge or many neurons die, try Leaky ReLU or reduce the
  Learning Rate. Use \textbf{Softmax} only for the final output.
\end{itemize}

Here are the complete study notes on \textbf{Principal Component
Analysis (PCA)}, including the mathematical derivation via Lagrangian
multipliers, and a comparative analysis with \textbf{Autoencoders}.

\begin{verbatim}
\newpage
\end{verbatim}

\section{Principal Component Analysis (PCA) \&
Autoencoders}\label{principal-component-analysis-pca-autoencoders}

Principal Component Analysis (PCA) is the most fundamental algorithm for
\textbf{Dimensionality Reduction} in Unsupervised Learning.

\textbf{The Goal:} Given a dataset \(X \in \mathbb{R}^{N \times d}\),
project it onto a lower-dimensional subspace of dimension \(k\)
(\(k < d\)) such that we retain as much information (variance) as
possible.

\subsection{1. Geometric Intuition}\label{geometric-intuition-1}

PCA can be interpreted in two mathematically equivalent ways:

\begin{enumerate}
\def\labelenumi{\arabic{enumi}.}
\item
  \textbf{Maximum Variance Formulation:} Find a direction (vector) in
  space along which the data varies the most. We want to preserve the
  ``spread'' of the data.
\item
  \textbf{Minimum Error Formulation:} Find a linear subspace
  (line/plane) such that the average squared distance (reconstruction
  error) between the original points and their projections is minimized.
\end{enumerate}

We will focus on the \textbf{Maximum Variance} approach as it leads
directly to the eigenvalue problem.

\subsection{2. Mathematical Derivation (The
Lagrangian)}\label{mathematical-derivation-the-lagrangian}

\subsubsection{2.1 Assumptions and Setup}\label{assumptions-and-setup}

Let our dataset be \(X = \{x^{(1)}, \dots, x^{(N)}\}\), where
\(x^{(i)} \in \mathbb{R}^d\).

\textbf{Pre-requisite:} The data must be \textbf{centered} (mean is
zero).

\[\mu = \frac{1}{N} \sum_{i=1}^N x^{(i)} = 0\]

We want to find a unit vector \(u\) (\(||u||^2 = u^T u = 1\)) such that
the variance of the data projected onto \(u\) is maximized.

The projection of a point \(x^{(i)}\) onto \(u\) is scalar
\(z_i = u^T x^{(i)}\).

\subsubsection{2.2 The Variance Objective}\label{the-variance-objective}

The variance of the projected points \(z\) is:

\[\text{Var}(z) = \frac{1}{N} \sum_{i=1}^N (z_i - \bar{z})^2\]

Since the data is centered, \(\bar{z} = 0\).

\[\text{Var}(z) = \frac{1}{N} \sum_{i=1}^N (u^T x^{(i)})^2 = \frac{1}{N} \sum_{i=1}^N u^T x^{(i)} (x^{(i)})^T u\]

We can factor out \(u\):

\[\text{Var}(z) = u^T \left( \frac{1}{N} \sum_{i=1}^N x^{(i)} (x^{(i)})^T \right) u\]

The term in the parenthesis is exactly the \textbf{Covariance Matrix},
denoted as \(\Sigma\) (or \(S\)).

\textbf{Objective Function:}

\[\max_u \ u^T \Sigma u \quad \text{subject to} \quad u^T u = 1\]

\subsubsection{2.3 Solving with Lagrangian
Multipliers}\label{solving-with-lagrangian-multipliers}

To solve this constrained optimization, we construct the Lagrangian
\(\mathcal{L}(u, \lambda)\):

\[\mathcal{L}(u, \lambda) = u^T \Sigma u - \lambda (u^T u - 1)\]

We take the derivative with respect to the vector \(u\) and set it to 0:

\[\nabla_u \mathcal{L} = \frac{\partial}{\partial u} (u^T \Sigma u) - \frac{\partial}{\partial u} (\lambda u^T u)\]

Recall vector calculus rules:
\(\frac{\partial}{\partial x} (x^T A x) = 2Ax\) (if A is symmetric) and
\(\frac{\partial}{\partial x} (x^T x) = 2x\).

\[2 \Sigma u - 2 \lambda u = 0\]

Dividing by 2, we get the fundamental result:

\[\Sigma u = \lambda u\]

\subsubsection{2.4 Interpretation}\label{interpretation-1}

This is the definition of the \textbf{Eigenvalue Problem}.

\begin{itemize}
\item
  \textbf{Conclusion:} The direction \(u\) that maximizes variance must
  be an \textbf{eigenvector} of the covariance matrix \(\Sigma\).
\item
  \textbf{The Variance:} If we plug this back into the objective
  function:
  \(u^T \Sigma u = u^T (\lambda u) = \lambda (u^T u) = \lambda\).
\item
  \textbf{Therefore:} The variance along direction \(u\) is exactly the
  eigenvalue \(\lambda\). To maximize variance, we pick the eigenvector
  corresponding to the \textbf{largest eigenvalue}.
\end{itemize}

\subsection{3. The PCA Algorithm:
Step-by-Step}\label{the-pca-algorithm-step-by-step}

\begin{enumerate}
\def\labelenumi{\arabic{enumi}.}
\item
  \textbf{Standardization:}

  \begin{itemize}
  \item
    Compute the mean \(\mu\) and subtract it:
    \(x^{(i)} := x^{(i)} - \mu\).
  \item
    (Optional but recommended) Divide by standard deviation to scale
    features to unit variance.
  \end{itemize}
\item
  \textbf{Covariance Matrix:}

  \begin{itemize}
  \tightlist
  \item
    Compute \(\Sigma = \frac{1}{N-1} X^T X\) (where \(X\) is the
    centered data matrix).
  \end{itemize}
\item
  \textbf{Eigendecomposition:}

  \begin{itemize}
  \item
    Find eigenvectors \(u_1, \dots, u_d\) and eigenvalues
    \(\lambda_1, \dots, \lambda_d\) of \(\Sigma\).
  \item
    \emph{Note:} In practice, we use \textbf{SVD (Singular Value
    Decomposition)} on \(X\) directly for numerical stability, rather
    than calculating \(\Sigma\) explicitly.
  \end{itemize}
\item
  \textbf{Sort and Select:}

  \begin{itemize}
  \item
    Sort eigenvalues in descending order:
    \(\lambda_1 \geq \lambda_2 \geq \dots \geq \lambda_d\).
  \item
    Choose the top \(k\) eigenvectors. These form the \textbf{projection
    matrix} \(W_k = [u_1, \dots, u_k]\).
  \end{itemize}
\item
  \textbf{Projection:}

  \begin{itemize}
  \item
    Transform the original data \(X\) into the new subspace \(Z\).
  \item
    \(Z = X W_k\).
  \end{itemize}
\end{enumerate}

\subsubsection{\texorpdfstring{3.1 Choosing \(k\) (Explained
Variance)}{3.1 Choosing k (Explained Variance)}}\label{choosing-k-explained-variance}

We select \(k\) by looking at the ratio of variance preserved.

\[\text{Explained Variance Ratio} = \frac{\sum_{i=1}^k \lambda_i}{\sum_{j=1}^d \lambda_j}\]

We typically pick \(k\) such that we preserve 95\% or 99\% of the
variance.

\subsection{4. PCA vs.~Autoencoders (AE)}\label{pca-vs.-autoencoders-ae}

An Autoencoder is a neural network trained to copy its input to its
output (\(y=x\)) through a ``bottleneck'' layer.

\subsubsection{4.1 The Architecture}\label{the-architecture}

\begin{itemize}
\item
  \textbf{Encoder:} Maps input \(x \in \mathbb{R}^d\) to code
  \(z \in \mathbb{R}^k\).
\item
  \textbf{Decoder:} Maps code \(z\) back to reconstruction
  \(\hat{x} \in \mathbb{R}^d\).
\item
  \textbf{Loss:} Minimize \(||x - \hat{x}||^2\) (MSE).
\end{itemize}

\subsubsection{4.2 The Connection (Linear
AE)}\label{the-connection-linear-ae}

If we build an Autoencoder with:

\begin{enumerate}
\def\labelenumi{\arabic{enumi}.}
\item
  \textbf{Linear} activation functions only.
\item
  \textbf{MSE} Loss.
\item
  A bottleneck of size \(k\).
\end{enumerate}

\textbf{Theorem:} The Linear Autoencoder learns the \textbf{same
subspace} as PCA. The hidden layer activations will span the same space
as the first \(k\) principal components.

\subsubsection{4.3 The Difference
(Non-Linearity)}\label{the-difference-non-linearity}

The limitation of PCA is that it is a \textbf{Linear Projection}. It can
only rotate and stretch the axes.

\begin{itemize}
\item
  \textbf{PCA Limit:} If the data lies on a curved surface (a non-linear
  manifold), PCA fails to unfold it. It will squash the data onto a flat
  plane.
\item
  \textbf{Deep Autoencoder:} By using non-linear activations (ReLU,
  Sigmoid), Autoencoders can learn \textbf{non-linear manifolds}. They
  can ``bend'' the projection space to fit the curve of the data.
\end{itemize}

\subsubsection{4.4 Example: The Swiss
Roll}\label{example-the-swiss-roll}

Imagine data shaped like a rolled-up carpet (Swiss Roll) in 3D.

\begin{itemize}
\item
  \textbf{PCA:} Would squash the roll flat like a pancake, merging
  points that are far apart on the roll but close in Euclidean 3D space.
\item
  \textbf{Autoencoder:} Can ``unroll'' the carpet, preserving the
  intrinsic geometry.
\end{itemize}

{\def\LTcaptype{none} % do not increment counter
\begin{longtable}[]{@{}
  >{\raggedright\arraybackslash}p{(\linewidth - 4\tabcolsep) * \real{0.3333}}
  >{\raggedright\arraybackslash}p{(\linewidth - 4\tabcolsep) * \real{0.3333}}
  >{\raggedright\arraybackslash}p{(\linewidth - 4\tabcolsep) * \real{0.3333}}@{}}
\toprule\noalign{}
\begin{minipage}[b]{\linewidth}\raggedright
\textbf{Feature}
\end{minipage} & \begin{minipage}[b]{\linewidth}\raggedright
\textbf{PCA}
\end{minipage} & \begin{minipage}[b]{\linewidth}\raggedright
\textbf{Autoencoder (Deep)}
\end{minipage} \\
\midrule\noalign{}
\endhead
\bottomrule\noalign{}
\endlastfoot
\textbf{Transformation} & Linear (\(W^T x\)) & Non-Linear (e.g.,
\(f(W x + b)\)) \\
\textbf{Optimization} & Convex (Closed-form SVD) & Non-Convex (Gradient
Descent) \\
\textbf{Global/Local} & Global optimum guaranteed & Prone to Local
Minima \\
\textbf{Computational} & Fast for small \(d\) & Slow (requires
training) \\
\textbf{Flexibility} & Rigid rotation & Can learn any complex
topology \\
\end{longtable}
}

\subsubsection{Key Takeaways}\label{key-takeaways-25}

\begin{itemize}
\item
  \textbf{PCA Goal:} Maximize Variance = Minimize Reconstruction Error.
\item
  \textbf{Covariance Matrix:} Encodes the spread and correlations of the
  data.
\item
  \textbf{Lagrangian:} Proves that the optimal directions are the
  \textbf{eigenvectors} of the Covariance Matrix.
\item
  \textbf{Decorrelation:} PCA rotates data so that features become
  uncorrelated (diagonal covariance).
\item
  \textbf{Autoencoders:} A generalization of PCA.

  \begin{itemize}
  \item
    \textbf{Linear AE \(\approx\) PCA.}
  \item
    \textbf{Non-Linear AE \textgreater{} PCA} for complex manifolds
    (like images or curved geometric structures).
  \item
    PCA is preferred for simple noise reduction; AEs are preferred for
    complex generative tasks or compression.
  \end{itemize}
\end{itemize}

\newpage

\section{Autoencoders (AE) and the
Bottleneck}\label{autoencoders-ae-and-the-bottleneck}

An Autoencoder is a specific type of \textbf{Neural Network} trained to
reconstruct its own input.\\
Given an input \(x\), the model produces an output \(\hat{x}\) that
should be as close as possible to \(x\).

Although Autoencoders are commonly described as \textbf{Unsupervised}
models, it is often more precise to call them \textbf{Self-Supervised},
because the target is automatically derived from the input itself.

The fundamental idea is to train a network to ``copy'' its input to its
output, but under a constraint that forces it to learn meaningful
structure rather than simply memorizing the data.

\subsection{1. The Architecture}\label{the-architecture-1}

An Autoencoder consists of two distinct components connected through a
latent representation.

\subsubsection{1.1 The Components}\label{the-components}

\begin{enumerate}
\def\labelenumi{\arabic{enumi}.}
\item
  \textbf{The Encoder (\(f_\phi\))}

  \begin{itemize}
  \item
    A neural network function that maps the input \(x\) into a
    latent-space representation \(z\).
  \item
    \(z = f_\phi(x)\)
  \end{itemize}
\item
  \textbf{The Bottleneck (Latent Space)}

  \begin{itemize}
  \item
    The internal representation where the information is compressed.
  \item
    In the classical setting, the latent dimension \(k\) is smaller than
    the input dimension \(d\).
  \end{itemize}
\item
  \textbf{The Decoder (\(g_\theta\))}

  \begin{itemize}
  \item
    A neural network function that reconstructs the input from the
    latent representation \(z\).
  \item
    \(\hat{x} = g_\theta(z)\)
  \end{itemize}
\end{enumerate}

\subsubsection{1.2 Mathematical
Formulation}\label{mathematical-formulation-3}

Let \(x \in \mathbb{R}^d\) be the input.

Let \(z \in \mathbb{R}^k\) be the latent code.

Let \(\hat{x} \in \mathbb{R}^d\) be the reconstruction.

The full forward pass is:

\[
\hat{x} = g_\theta(f_\phi(x))
\]

We train the parameters \(\phi\) (encoder weights) and \(\theta\)
(decoder weights) jointly to minimize the reconstruction error.

\subsection{2. The Bottleneck: The Core
Constraint}\label{the-bottleneck-the-core-constraint}

The \textbf{Bottleneck} is the hidden representation with dimension
\(k\), typically such that:

\[
k < d
\]

This architecture is called an \textbf{Undercomplete Autoencoder}.

\subsubsection{2.1 Why is the Bottleneck
necessary?}\label{why-is-the-bottleneck-necessary}

If we allowed the latent space to have dimension \(k \ge d\), and the
network had enough capacity, the model could simply learn the
\textbf{Identity Function}:

\[
f(x) = x
\]

In that case, the network would just copy the input to the output
without extracting useful structure or features. It would be useless for
compression or representation learning.

\subsubsection{2.2 Forced Compression}\label{forced-compression}

By imposing \(k < d\), we force the network to compress the input.

To minimize reconstruction error, the model must learn to:

\begin{itemize}
\item
  \textbf{Retain:} the most important features, correlations, and
  patterns
\item
  \textbf{Discard:} redundant information and noise
\end{itemize}

This forces the latent variable \(z\) to capture the most relevant
structure of the data.

\subsubsection{2.3 Lossy Compression}\label{lossy-compression}

Because \(k < d\), the reconstruction is generally not perfect:

\[
\hat{x} \approx x
\]

Autoencoders therefore perform \textbf{Lossy Compression}, meaning that
some information is lost while the essential structure is preserved.

\subsection{3. The Bottleneck Alone Is Not Always
Enough}\label{the-bottleneck-alone-is-not-always-enough}

A bottleneck is a strong constraint, but it is not always sufficient to
guarantee useful representations.

If the encoder and decoder are very expressive, the model may still
learn solutions that minimize reconstruction loss without learning
semantically meaningful features.

For this reason, Autoencoders are often combined with additional
regularization techniques such as:

\begin{itemize}
\item
  noise injection
\item
  sparsity constraints
\item
  contractive penalties
\item
  limited architectural capacity
\item
  weight decay or dropout
\end{itemize}

So, the bottleneck is one important mechanism, but not the only one.

\subsection{4. Types of Autoencoders}\label{types-of-autoencoders}

\subsubsection{4.1 Undercomplete
Autoencoder}\label{undercomplete-autoencoder}

This is the classical case:

\[
k < d
\]

The model is forced to compress the input into a smaller latent code.

\subsubsection{4.2 Overcomplete
Autoencoder}\label{overcomplete-autoencoder}

In this case:

\[
k \ge d
\]

Here, dimensional compression alone does not prevent the Autoencoder
from learning the identity function.

To make an overcomplete Autoencoder useful, we need extra constraints
such as:

\begin{itemize}
\item
  sparsity
\item
  denoising
\item
  contractive regularization
\item
  limited model capacity
\end{itemize}

So, an Autoencoder can still learn useful features even without a narrow
bottleneck, but only if it is properly regularized.

\subsection{5. The Objective Function
(Loss)}\label{the-objective-function-loss}

We train the Autoencoder using Backpropagation and Gradient Descent to
minimize the discrepancy between the input \(x\) and its reconstruction
\(\hat{x}\).

Given a dataset \(\{x^{(i)}\}_{i=1}^N\), we optimize:

\[
\min_{\phi,\theta} \frac{1}{N} \sum_{i=1}^N L\Big(x^{(i)}, g_\theta(f_\phi(x^{(i)}))\Big)
\]

where \(L\) is the reconstruction loss.

\subsection{6. Common Loss Functions}\label{common-loss-functions}

The choice of loss depends on the data type and on how we interpret the
output.

\subsubsection{6.1 Mean Squared Error
(MSE)}\label{mean-squared-error-mse}

For continuous data, we commonly use the squared Euclidean distance:

\[
\mathcal{L}_{\text{MSE}} = \frac{1}{N} \sum_{i=1}^N \|x^{(i)} - \hat{x}^{(i)}\|^2
\]

This is typical when:

\begin{itemize}
\item
  the input is real-valued
\item
  the output is real-valued
\item
  we want to strongly penalize large reconstruction errors
\end{itemize}

\subsubsection{6.2 Binary Cross-Entropy
(BCE)}\label{binary-cross-entropy-bce}

If the input data is binary or normalized to \([0,1]\), we may interpret
the output as probabilities:

\[
\mathcal{L}_{\text{BCE}} =
- \sum_{i=1}^N \sum_{j=1}^d
\left[
x_{ij}\log(\hat{x}_{ij}) + (1 - x_{ij})\log(1 - \hat{x}_{ij})
\right]
\]

This is common when:

\begin{itemize}
\item
  the data is binary
\item
  image pixels are normalized to \([0,1]\)
\item
  the output layer uses a \textbf{sigmoid}
\end{itemize}

\subsubsection{6.3 Output Activation
Matters}\label{output-activation-matters}

The loss must be consistent with the final output activation:

\begin{itemize}
\item
  \textbf{Linear output} + MSE for unconstrained real-valued outputs
\item
  \textbf{Sigmoid output} + BCE for values in \([0,1]\)
\item
  \textbf{Tanh output} in some cases for data scaled to \([-1,1]\)
\end{itemize}

This consistency is important for stable and meaningful training.

\subsection{7. Relationship with PCA}\label{relationship-with-pca}

This is an important theoretical point.

If an Autoencoder has:

\begin{itemize}
\item
  a linear encoder
\item
  a linear decoder
\item
  no nonlinear activation functions
\item
  MSE loss
\end{itemize}

then it learns, in essence, the same principal subspace found by
\textbf{PCA}.

\subsubsection{Key Difference}\label{key-difference}

\begin{itemize}
\item
  \textbf{PCA} is a linear dimensionality reduction method
\item
  \textbf{Autoencoders} can be nonlinear
\end{itemize}

Because Autoencoders can use nonlinear activations, they can model more
complex structures and curved manifolds that PCA cannot capture.

For this reason, Autoencoders are often seen as a nonlinear
generalization of PCA.

\subsection{8. Nonlinearity and Representation
Learning}\label{nonlinearity-and-representation-learning}

Unlike PCA, Autoencoders can use nonlinear activations such as:

\begin{itemize}
\item
  ReLU
\item
  Sigmoid
\item
  Tanh
\end{itemize}

These nonlinearities allow the model to learn more expressive mappings
between the input space and the latent space.

This is why Autoencoders can represent data lying on complicated
low-dimensional manifolds embedded in high-dimensional spaces.

\subsection{9. Denoising Autoencoders
(DAE)}\label{denoising-autoencoders-dae}

A standard Autoencoder may still learn a near-identity mapping,
especially if it has too much capacity.\\
To make the representation more robust, we use \textbf{Denoising
Autoencoders}.

\subsubsection{9.1 The Mechanism}\label{the-mechanism}

\begin{enumerate}
\def\labelenumi{\arabic{enumi}.}
\item
  \textbf{Corrupt:} Start from a clean input \(x\) and add noise to
  obtain \(\tilde{x}\)
\item
  \textbf{Feed:} Pass the corrupted input \(\tilde{x}\) into the encoder
\item
  \textbf{Target:} Force the decoder to reconstruct the original clean
  input \(x\)
\end{enumerate}

The loss becomes:

\[
\mathcal{L} = \|x - g_\theta(f_\phi(\tilde{x}))\|^2
\]

\subsubsection{9.2 Types of Corruption}\label{types-of-corruption}

Common corruption methods include:

\begin{itemize}
\item
  Gaussian noise
\item
  masking noise
\item
  salt-and-pepper noise
\item
  dropout on the input
\end{itemize}

\subsubsection{9.3 Intuition}\label{intuition-1}

Since the input is corrupted, the network cannot simply memorize and
copy it.\\
It must learn the stable structure of the data and ignore accidental
perturbations.

This usually leads to a more robust latent representation.

\subsection{10. Intuition from Manifold
Learning}\label{intuition-from-manifold-learning}

A useful geometric interpretation is that real-world data does not
occupy all of \(\mathbb{R}^d\), but lies near a lower-dimensional
\textbf{manifold}.

The Autoencoder learns:

\begin{itemize}
\item
  to map valid inputs into a latent representation
\item
  to reconstruct them back into the original space
\end{itemize}

In the denoising case, if noise pushes a sample away from the manifold,
the model learns to project it back toward a plausible clean point.

This interpretation helps explain why Autoencoders are effective for
denoising and feature extraction.

\subsection{11. Other Important
Variants}\label{other-important-variants}

\subsubsection{11.1 Sparse Autoencoder}\label{sparse-autoencoder}

A \textbf{Sparse Autoencoder} does not necessarily require a small
latent dimension.\\
Instead, it imposes that only a few latent units are active for each
input.

This is usually done by adding a sparsity regularization term, such as:

\begin{itemize}
\item
  \(L_1\) penalty on activations
\item
  KL-divergence-based sparsity constraints
\end{itemize}

Effect:

\begin{itemize}
\item
  the representation becomes more selective
\item
  the model is discouraged from using all latent units at once
\item
  trivial copying becomes harder
\end{itemize}

\subsubsection{11.2 Contractive
Autoencoder}\label{contractive-autoencoder}

A \textbf{Contractive Autoencoder} adds a penalty that reduces the
sensitivity of the latent representation to small input perturbations:

\[
\lambda \left\| \frac{\partial f_\phi(x)}{\partial x} \right\|^2
\]

This encourages local robustness and smoother learned representations.

\subsubsection{11.3 Variational Autoencoder
(VAE)}\label{variational-autoencoder-vae}

This is a very important extension.

In a \textbf{Variational Autoencoder}, the encoder does not map \(x\) to
a single deterministic latent vector, but to a probability distribution:

\[
q_\phi(z \mid x)
\]

Typically, the encoder outputs the parameters of a Gaussian
distribution:

\[
\mu(x), \sigma(x)
\]

A latent sample \(z\) is then drawn from this distribution and passed to
the decoder.

The loss has two components:

\begin{enumerate}
\def\labelenumi{\arabic{enumi}.}
\item
  reconstruction loss
\item
  a KL-divergence term that regularizes the latent distribution
\end{enumerate}

\subsubsection{Why is the VAE
different?}\label{why-is-the-vae-different}

\begin{itemize}
\item
  \textbf{Standard AE:} deterministic latent code, mainly for
  reconstruction and feature learning
\item
  \textbf{VAE:} probabilistic latent space, explicitly designed also for
  generation
\end{itemize}

A classical Autoencoder does not generally define a well-structured
latent space for sampling new data, while a VAE does.

\subsection{12. Capacity and
Overfitting}\label{capacity-and-overfitting}

Autoencoders can also \textbf{overfit}.

If the model has too many parameters relative to the amount of training
data, it may reconstruct the training set extremely well but generalize
poorly to unseen inputs.

This is especially important in applications such as:

\begin{itemize}
\item
  anomaly detection
\item
  denoising
\item
  learned compression
\end{itemize}

To reduce overfitting, common strategies include:

\begin{itemize}
\item
  early stopping
\item
  weight decay
\item
  dropout
\item
  noise injection
\item
  validation sets
\item
  appropriate model size
\end{itemize}

So, low training reconstruction error alone is not enough to conclude
that the model learned something useful.

\subsection{\texorpdfstring{13. Choosing the Latent Dimension
\(k\)}{13. Choosing the Latent Dimension k}}\label{choosing-the-latent-dimension-k}

The latent dimension \(k\) is a crucial hyperparameter.

If \(k\) is too small:

\begin{itemize}
\item
  too much information is lost
\item
  reconstruction quality becomes poor
\end{itemize}

If \(k\) is too large:

\begin{itemize}
\item
  the model may behave almost like the identity function
\item
  the latent representation may be less informative
\end{itemize}

In practice, \(k\) is chosen based on:

\begin{itemize}
\item
  validation reconstruction error
\item
  usefulness of the latent features for downstream tasks
\item
  latent space visualization
\item
  desired compression level
\end{itemize}

There is no universal rule for selecting \(k\).

\subsection{14. Typical Architectures in
Practice}\label{typical-architectures-in-practice}

The architecture depends on the type of data.

\subsubsection{14.1 Tabular or Vector
Data}\label{tabular-or-vector-data}

For vector inputs, we often use fully connected networks:

\begin{itemize}
\item
  encoder with dense layers
\item
  bottleneck
\item
  decoder with dense layers
\end{itemize}

\subsubsection{14.2 Image Data}\label{image-data}

For images, we usually use \textbf{Convolutional Autoencoders}:

\begin{itemize}
\item
  convolutional layers in the encoder
\item
  compressed feature maps in the latent representation
\item
  transposed convolutions or upsampling in the decoder
\end{itemize}

Convolutional Autoencoders are much more suitable for images because
they exploit spatial structure.

\subsection{15. Applications of
Autoencoders}\label{applications-of-autoencoders}

\subsubsection{15.1 Dimensionality
Reduction}\label{dimensionality-reduction}

The encoder produces a compact representation:

\[
z = f_\phi(x)
\]

This can be used instead of the original input for:

\begin{itemize}
\item
  visualization
\item
  clustering
\item
  preprocessing
\item
  downstream supervised tasks
\end{itemize}

\subsubsection{15.2 Feature Learning}\label{feature-learning}

Autoencoders can learn useful hidden representations that capture
higher-level structure better than the raw input.

\subsubsection{15.3 Anomaly Detection}\label{anomaly-detection}

A common strategy is:

\begin{enumerate}
\def\labelenumi{\arabic{enumi}.}
\item
  train the Autoencoder only on normal data
\item
  reconstruct new samples at test time
\item
  compute the reconstruction error
\end{enumerate}

If the error is high:

\[
\|x - \hat{x}\|^2
\]

then the sample may be anomalous, because the network has not learned to
represent that type of pattern well.

\paragraph{Important Limitation}\label{important-limitation}

This works well only if:

\begin{itemize}
\item
  the training set contains mostly normal data
\item
  anomalies are sufficiently different from normal samples
\item
  the model is not too expressive
\end{itemize}

A very powerful Autoencoder may reconstruct anomalies well too, reducing
detection performance.

\subsubsection{15.4 Image Denoising}\label{image-denoising}

Denoising Autoencoders can learn to recover clean images from noisy
inputs.

\subsubsection{15.5 Learned Compression}\label{learned-compression}

The latent representation can serve as a compressed version of the data.

Unlike classical codecs such as JPEG, this compression is
\textbf{data-dependent}: it works best on data similar to the training
distribution.

\subsubsection{15.6 Super-Resolution and
Restoration}\label{super-resolution-and-restoration}

Autoencoders can also be trained to reconstruct higher-quality outputs
from degraded or low-resolution inputs.

\subsection{16. Limitations of
Autoencoders}\label{limitations-of-autoencoders}

Although Autoencoders are powerful, they also have important
limitations.

\subsubsection{16.1 Good Reconstruction Does Not Guarantee Semantic
Understanding}\label{good-reconstruction-does-not-guarantee-semantic-understanding}

A model may reconstruct the input well without learning features that
are truly meaningful or interpretable.

\subsubsection{16.2 The Latent Space May Be Poorly
Structured}\label{the-latent-space-may-be-poorly-structured}

In a standard Autoencoder, nearby latent points do not necessarily
correspond to semantically similar samples in a smooth or organized way.

\subsubsection{16.3 Poor Generative
Ability}\label{poor-generative-ability}

Sampling random points from the latent space of a classical Autoencoder
does not usually generate realistic new data.

For proper generative modeling, VAEs or other generative models are
better suited.

\subsubsection{16.4 MSE Can Produce Blurry
Reconstructions}\label{mse-can-produce-blurry-reconstructions}

In image-related tasks, MSE often encourages average-looking outputs,
which may appear blurry.

\subsection{17. Training Procedure}\label{training-procedure}

Training an Autoencoder follows the usual neural network optimization
pipeline, except that the target is the input itself:

\begin{enumerate}
\def\labelenumi{\arabic{enumi}.}
\item
  take an input batch \(x\)
\item
  compute the latent code \(z = f_\phi(x)\)
\item
  reconstruct the input \(\hat{x} = g_\theta(z)\)
\item
  compute the reconstruction loss
\item
  update parameters via backpropagation
\end{enumerate}

Common optimizers include:

\begin{itemize}
\item
  SGD
\item
  Adam
\end{itemize}

During training, we usually monitor:

\begin{itemize}
\item
  training loss
\item
  validation loss
\item
  reconstruction quality
\item
  usefulness of latent features for downstream tasks
\end{itemize}

\subsection{18. Summary}\label{summary}

An Autoencoder learns a transformation:

\[
x \mapsto z \mapsto \hat{x}
\]

where:

\begin{itemize}
\item
  \(z\) is a compressed or regularized latent representation
\item
  \(\hat{x}\) is a reconstruction of the original input
\end{itemize}

The \textbf{Bottleneck} is one of the main mechanisms that prevents
trivial copying and forces the model to learn structure.

However, the bottleneck is not the only useful constraint. In practice,
robust Autoencoders often rely on additional regularization such as
noise, sparsity, or contractive penalties.

Autoencoders are useful for:

\begin{itemize}
\item
  dimensionality reduction
\item
  feature learning
\item
  denoising
\item
  anomaly detection
\item
  learned compression
\end{itemize}

At the same time, they have clear limitations, especially in semantic
interpretability and generative modeling.

\subsection{Key Takeaways}\label{key-takeaways-26}

\begin{itemize}
\item
  \textbf{Structure:} Encoder (Compress) \(\rightarrow\) Latent Space
  (Bottleneck) \(\rightarrow\) Decoder (Reconstruct)
\item
  \textbf{Bottleneck:} the central constraint that limits information
  flow and forces compression
\item
  \textbf{Identity Trap:} without constraints, the model may simply
  learn \(\hat{x} = x\)
\item
  \textbf{Loss:} measures how close the reconstruction is to the
  original input
\item
  \textbf{PCA Connection:} a linear Autoencoder with MSE is closely
  related to PCA
\item
  \textbf{Nonlinearity:} makes Autoencoders more expressive than PCA
\item
  \textbf{Variants:} DAE, Sparse AE, Contractive AE, and VAE add useful
  properties
\item
  \textbf{Applications:} compression, denoising, anomaly detection,
  feature extraction
\item
  \textbf{Limitation:} good reconstruction does not automatically mean a
  good semantic representation
\end{itemize}

\end{document}
